% Interaction gravitationnelle et champ de gravitation — Physique Tle E — Cours signé OnBuch+
% Programme officiel MINESEC. Compiler avec : tectonic P02-champ-de-gravitation.tex
\newcommand{\DOCMATIERE}{Physique}
\newcommand{\DOCNIVEAU}{Tle E}
\newcommand{\DOCMODULE}{Module 2 : Mouvements et interactions — évolutions temporelles des systèmes mécaniques}
\newcommand{\DOCLECON}{P02}
\newcommand{\DOCTITRE}{Interaction gravitationnelle et champ de gravitation}
% OnBuch+ — préambule « cours signés » ENRICHI (compilé avec tectonic / XeLaTeX)
% Système visuel « Pop » d'OnBuch+ : fond crème (jamais blanc), bordures encre
% épaisses, ombres « dures » décalées, titres Archivo Black en capitales.
% Version MATHÉMATIQUES : graphes (pgfplots/tikz), boîtes pédagogiques
% (définition, propriété, méthode, attention, le savais-tu, exercice résolu,
% exercice, TP, bilan), page de garde et signature OnBuch+ ancrée en bas.
%
% Macros attendues dans main.tex AVANT \input{preamble} :
%   \DOCMATIERE  \DOCNIVEAU  \DOCMODULE  \DOCLECON (numéro)  \DOCTITRE
\documentclass[11pt]{article}

\usepackage{fontspec}
\usepackage[french]{babel}
\usepackage[a4paper,margin=2cm,top=3.4cm,bottom=2.8cm,headheight=1.4cm,footskip=1.3cm]{geometry}
\usepackage{amsmath,amssymb,amsfonts}
\usepackage{mathtools} % \xmapsto, \coloneqq... souvent utilisés par les modèles
\usepackage{cancel} % simplifications barrées (bilans de forces)
% \abs{z} (module, valeur absolue) et \norm{u} (norme), utilisés par les modèles
\providecommand{\abs}{}\let\abs\relax\DeclarePairedDelimiter{\abs}{\lvert}{\rvert}
\providecommand{\norm}{}\let\norm\relax\DeclarePairedDelimiter{\norm}{\lVert}{\rVert}
\usepackage[table]{xcolor} % [table] : fournit aussi \rowcolors (alternance de lignes)
\usepackage{pagecolor}
\usepackage{tikz}
\usetikzlibrary{arrows.meta,positioning,calc,shapes.geometric,shapes.misc,shapes.arrows,decorations.pathmorphing,decorations.markings,patterns,shadows,mindmap,trees,fit,backgrounds,matrix,intersections,angles,quotes}
\usepackage{pgfplots}
\usepackage[version=4]{mhchem} % équations nucléaires \ce{^{235}_{92}U}
\usepackage[european,siunitx]{circuitikz} % schémas électriques
\pgfplotsset{compat=1.18}
\usepgfplotslibrary{fillbetween,groupplots} % aires entre courbes (calcul intégral)
\usepackage{enumitem}
\usepackage{fancyhdr}
% [most] charge toutes les bibliothèques tcolorbox (listings, minted, raster,
% documentation...) et gonfle inutilement la mémoire TeX sur un document long
% (~300 boîtes) : on ne charge que ce qui est réellement utilisé.
\usepackage[skins,breakable]{tcolorbox}
\usepackage{graphicx}
\usepackage{colortbl}
\usepackage{booktabs}
\usepackage{tabularx}
\usepackage{diagbox} % case d'en-tête diagonale des tableaux à double entrée
% Colonnes extensibles centrée / gauche / droite (C, L, R), souvent utilisées par les modèles
\newcolumntype{C}{>{\centering\arraybackslash}X}
\newcolumntype{L}{>{\raggedright\arraybackslash}X}
\newcolumntype{R}{>{\raggedleft\arraybackslash}X}
\usepackage{array}
\usepackage{multicol}
\usepackage{siunitx}
% parse-numbers=false : accepte aussi \SI{2,5 \times 10^{-3}}{...}
\sisetup{locale=FR,output-decimal-marker={,},group-minimum-digits=4,parse-numbers=false}
\DeclareSIUnit\ohm{\ensuremath{\symup{\Omega}}} % U+2126 absent de la police
\DeclareSIUnit\division{div} % réglages d'oscilloscope (ms/div, V/div)
\usepackage{qrcode}
% \degree : commande fréquemment utilisée par les modèles pour un angle ou une
% température en dehors de \SI{}{} (non fournie par les paquets chargés).
\providecommand{\degree}{^{\circ}}
% \qed (fin de démonstration), utilisé par les modèles sans amsthm
\providecommand{\qed}{\hfill$\square$}
% \barre : double barre des valeurs interdites dans un tableau de variations (tkz-tab)
\providecommand{\barre}{\ensuremath{\|}}
% environnement proof (amsthm non chargé) utilisé par les modèles
\ifdefined\proof\else\newenvironment{proof}[1][Démonstration]{\par\noindent\textit{#1.}\ }{\qed\par}\fi
\usepackage{lastpage}
\usepackage{hyperref}
\hypersetup{hidelinks}

% ============================================================
% Palette « Pop » OnBuch+ — mêmes hex que la classe P (plus_ui.dart)
% ============================================================
\definecolor{poporange}{HTML}{F59321}
\definecolor{popdark}{HTML}{E07A0C}
\definecolor{popcream}{HTML}{FFF6E8}
\definecolor{popink}{HTML}{1C1714}
\definecolor{poppurple}{HTML}{7A5AE0}
\definecolor{popgreen}{HTML}{2E9E6A}
\definecolor{poppink}{HTML}{E0457B}
\definecolor{popblue}{HTML}{2F7DE1}
\definecolor{popgold}{HTML}{C98A12}
\definecolor{popmuted}{HTML}{8A7F76}
\definecolor{popline}{HTML}{EADFD2}
\definecolor{popgray}{HTML}{8A7F76}
\colorlet{popgrey}{popgray}
% Alias tolérants (le modèle nomme parfois les couleurs différemment)
\colorlet{popwhite}{white}
\colorlet{popblack}{popink}
\colorlet{popred}{poppink}
\colorlet{popyellow}{popgold}
% Teintes claires pour fonds de tableaux / zones
\colorlet{popblueL}{popblue!10!white}
\colorlet{poporangeL}{poporange!14!white}
\colorlet{popgreenL}{popgreen!12!white}
\colorlet{poppurpleL}{poppurple!12!white}
\colorlet{poppinkL}{poppink!12!white}
\colorlet{popgoldL}{popgold!14!white}

\pagecolor{popcream}

% ============================================================
% Typographie
% ============================================================
\defaultfontfeatures{Ligatures=TeX}
\setmainfont{PlusJakartaSans-Medium.ttf}[Path=../../../fonts/,BoldFont=PlusJakartaSans-Bold.ttf]
% Maths en sans-serif (Fira Math) avec chiffres et lettres droites de Plus
% Jakarta, pour que les formules se fondent dans le texte.
\usepackage[math-style=ISO,bold-style=ISO]{unicode-math}
\setmathfont{FiraMath-Regular.otf}[Path=../../../fonts/]
\setmathfont{PlusJakartaSans-Medium.ttf}[Path=../../../fonts/,range={up/num,up/{Latin,latin},\mathup}]
\usepackage{icomma} % virgule décimale sans espace en mode math
% Caractères mathématiques Unicode absents des polices de texte (Plus Jakarta,
% Archivo Black) : rendus par leur équivalent mathématique, en texte comme en
% formule — sinon ils s'affichent en carré vide (ex. « Arithmétique dans ℤ »).
\usepackage{newunicodechar}
\newunicodechar{ℝ}{\ensuremath{\mathbb{R}}}
\newunicodechar{ℤ}{\ensuremath{\mathbb{Z}}}
\newunicodechar{ℂ}{\ensuremath{\mathbb{C}}}
\newunicodechar{ℕ}{\ensuremath{\mathbb{N}}}
\newunicodechar{ℚ}{\ensuremath{\mathbb{Q}}}
\newunicodechar{𝔻}{\ensuremath{\mathbb{D}}}
\newunicodechar{α}{\ensuremath{\alpha}}
\newunicodechar{β}{\ensuremath{\beta}}
\newunicodechar{γ}{\ensuremath{\gamma}}
\newunicodechar{δ}{\ensuremath{\delta}}
\newunicodechar{ε}{\ensuremath{\varepsilon}}
\newunicodechar{θ}{\ensuremath{\theta}}
\newunicodechar{λ}{\ensuremath{\lambda}}
\newunicodechar{μ}{\ensuremath{\mu}}
\newunicodechar{σ}{\ensuremath{\sigma}}
\newunicodechar{τ}{\ensuremath{\tau}}
\newunicodechar{φ}{\ensuremath{\varphi}}
\newunicodechar{ψ}{\ensuremath{\psi}}
\newunicodechar{ω}{\ensuremath{\omega}}
\newunicodechar{Σ}{\ensuremath{\Sigma}}
% majuscules produites par \MakeUppercase dans les titres (α-aminés -> Α)
\newunicodechar{Α}{\ensuremath{\alpha}}
\newunicodechar{Β}{\ensuremath{\beta}}
\newunicodechar{Γ}{\ensuremath{\Gamma}}
\newunicodechar{Δ}{\ensuremath{\Delta}}
\newunicodechar{Ε}{\ensuremath{\varepsilon}}
\newunicodechar{Θ}{\ensuremath{\Theta}}
\newunicodechar{Λ}{\ensuremath{\Lambda}}
\newunicodechar{Μ}{\ensuremath{\mu}}
\newunicodechar{Τ}{\ensuremath{\tau}}
\newunicodechar{Φ}{\ensuremath{\Phi}}
\newunicodechar{Ψ}{\ensuremath{\Psi}}
\newunicodechar{Ω}{\ensuremath{\Omega}}
\newunicodechar{↦}{\ensuremath{\mapsto}}
\newunicodechar{∈}{\ensuremath{\in}}
\newunicodechar{∉}{\ensuremath{\notin}}
\newunicodechar{∧}{\ensuremath{\wedge}}
\newunicodechar{⇔}{\ensuremath{\Leftrightarrow}}
\newunicodechar{⟺}{\ensuremath{\Longleftrightarrow}}
\newunicodechar{⇒}{\ensuremath{\Rightarrow}}
\newunicodechar{⟹}{\ensuremath{\Longrightarrow}}
\newunicodechar{∘}{\ensuremath{\circ}}
\newunicodechar{∪}{\ensuremath{\cup}}
\newunicodechar{∩}{\ensuremath{\cap}}
\newunicodechar{≡}{\ensuremath{\equiv}}
\newunicodechar{⊂}{\ensuremath{\subset}}
\newunicodechar{⊥}{\ensuremath{\perp}}
\newunicodechar{∃}{\ensuremath{\exists}}
\newunicodechar{∀}{\ensuremath{\forall}}
\newunicodechar{∛}{\ensuremath{\sqrt[3]{\,}}}
\newunicodechar{∜}{\ensuremath{\sqrt[4]{\,}}}
\newunicodechar{⃗}{\ensuremath{{}^{\to}}}
\newunicodechar{ⁿ}{\ifmmode^{n}\else\textsuperscript{\ensuremath{n}}\fi}
\newunicodechar{ˣ}{\ifmmode^{x}\else\textsuperscript{\ensuremath{x}}\fi}
\newunicodechar{ᵇ}{\ifmmode^{b}\else\textsuperscript{\ensuremath{b}}\fi}
\newunicodechar{ʳ}{\ifmmode^{r}\else\textsuperscript{\ensuremath{r}}\fi}
\newunicodechar{ᵘ}{\ifmmode^{u}\else\textsuperscript{\ensuremath{u}}\fi}
\newunicodechar{ʸ}{\ifmmode^{y}\else\textsuperscript{\ensuremath{y}}\fi}
\newunicodechar{ᵃ}{\ifmmode^{a}\else\textsuperscript{\ensuremath{a}}\fi}
\newunicodechar{ᵉ}{\ifmmode^{e}\else\textsuperscript{\ensuremath{e}}\fi}
\newunicodechar{ᵏ}{\ifmmode^{k}\else\textsuperscript{\ensuremath{k}}\fi}
\newunicodechar{ᵖ}{\ifmmode^{p}\else\textsuperscript{\ensuremath{p}}\fi}
\newunicodechar{ᴺ}{\ifmmode^{N}\else\textsuperscript{\ensuremath{N}}\fi}
\newunicodechar{ᶜ}{\ifmmode^{c}\else\textsuperscript{\ensuremath{c}}\fi}
\newunicodechar{ʰ}{\ifmmode^{h}\else\textsuperscript{\ensuremath{h}}\fi}
\newunicodechar{ᵠ}{\ifmmode^{q}\else\textsuperscript{\ensuremath{q}}\fi}
\newunicodechar{ₙ}{\ifmmode_{n}\else\textsubscript{\ensuremath{n}}\fi}
\newunicodechar{ₐ}{\ifmmode_{a}\else\textsubscript{\ensuremath{a}}\fi}
\newunicodechar{ᵢ}{\ifmmode_{i}\else\textsubscript{\ensuremath{i}}\fi}
\newunicodechar{ᵣ}{\ifmmode_{r}\else\textsubscript{\ensuremath{r}}\fi}
\newunicodechar{ₖ}{\ifmmode_{k}\else\textsubscript{\ensuremath{k}}\fi}
\newunicodechar{ᵦ}{\ifmmode_{\beta}\else\textsubscript{\ensuremath{\beta}}\fi}
\newunicodechar{ₓ}{\ifmmode_{x}\else\textsubscript{\ensuremath{x}}\fi}
\newfontfamily\popdisplay{ArchivoBlack-Regular.ttf}[Path=../../../fonts/]
\newfontfamily\poplogo{SpaceGrotesk-Bold.ttf}[Path=../../../fonts/]
\newfontfamily\popsemi{PlusJakartaSans-SemiBold.ttf}[Path=../../../fonts/]
\newfontfamily\popxbold{PlusJakartaSans-ExtraBold.ttf}[Path=../../../fonts/]
\newfontfamily\popxboldls{PlusJakartaSans-ExtraBold.ttf}[Path=../../../fonts/,LetterSpace=3.0]

\setlist[enumerate]{leftmargin=1.7em,itemsep=5pt,topsep=4pt}
\setlist[itemize]{leftmargin=1.7em,itemsep=4pt,topsep=4pt,label={\color{poporange}\textbullet}}
\setlength{\parindent}{0pt}
\setlength{\parskip}{5pt}
\renewcommand{\arraystretch}{1.25}

% pgfplots : style Pop par défaut
\pgfplotsset{
  every axis/.append style={axis line style={popink,line width=1pt},tick style={popink},
    grid style={popline,line width=0.5pt},label style={font=\small},tick label style={font=\footnotesize}},
  popcurve/.style={poporange,line width=1.8pt,no markers,smooth},
}
% Même style disponible dans un \draw TikZ simple (hors pgfplots)
\tikzset{popcurve/.style={poporange,line width=1.8pt}}
\tikzset{popgrid/.style={popline,very thin}}
\colorlet{popcurve}{poporange} % le nom du style est parfois utilisé comme couleur

% ============================================================
% Pill / squiggle
% ============================================================
\newcommand{\poppill}[3]{%
  \tikz[baseline=-0.55ex]\node[draw=popink,line width=1.2pt,rounded corners=2.6pt,%
    fill=#2,inner xsep=6.5pt,inner ysep=3pt]{%
    \popxboldls\fontsize{7}{7}\selectfont\color{#3}\MakeUppercase{#1}};%
}
\newcommand{\popsquiggle}[1]{%
  \tikz[baseline=-0.4ex]{
    \draw[#1,line width=2.6pt,line cap=round]
      plot [smooth] coordinates {(0,0.09) (0.34,0.01) (0.68,0.21) (1.02,0.01) (1.36,0.10)};
  }%
}
% Mot-clé surligné (marqueur orange sous le texte)
\newcommand{\cle}[1]{{\popxbold\color{popdark}#1}}

% ============================================================
% En-tête / pied
% ============================================================
\pagestyle{fancy}
\fancyhf{}
\renewcommand{\headrulewidth}{0pt}
\renewcommand{\footrulewidth}{0pt}
\lhead{\raisebox{-0.15ex}{\poplogo\fontsize{13}{13}\selectfont\color{popink}OnBuch\color{poporange}+}}
\rhead{\poppill{\DOCMATIERE\ \textbullet\ \DOCNIVEAU\ \textbullet\ Leçon \DOCLECON}{white}{popink}}
\lfoot{\popxbold\fontsize{7.6}{7.6}\selectfont\color{popmuted}\MakeUppercase{Cours signé OnBuch+}\ \textbullet\ {\popsemi\DOCTITRE}}
\rfoot{\tikz[baseline=-0.6ex]\node[rounded corners=8.5pt,fill=popink,minimum height=17pt,minimum width=17pt,inner xsep=5pt,inner sep=0]{\popxbold\fontsize{8}{8}\selectfont\color{popcream}\thepage\,/\,\pageref*{LastPage}};}

% ============================================================
% Titres
% ============================================================
\newcommand{\chaptitre}[2]{%
  \begin{tcolorbox}[enhanced,colback=poporange,colframe=popink,boxrule=2.6pt,arc=11pt,
      drop shadow=popink,left=15pt,right=15pt,top=12pt,bottom=11pt,before skip=3pt,after skip=18pt]
    \poppill{#2}{popink}{popcream}\par\vspace{9pt}
    {\raggedright\hyphenpenalty=10000\exhyphenpenalty=10000\popdisplay\fontsize{22}{24}\selectfont\color{popcream}\MakeUppercase{#1}\par}
  \end{tcolorbox}
}
% Section numérotée : pastille orange avec le numéro + titre Archivo
\newcounter{popsec}
\newcommand{\coursec}[1]{%
  \stepcounter{popsec}\setcounter{popsub}{0}%
  \par\vspace{16pt}\noindent
  \begin{minipage}{\linewidth}
  \tikz[baseline=-0.7ex]\node[draw=popink,line width=1.6pt,fill=poporange,rounded corners=4pt,minimum size=22pt,inner sep=0]{\popdisplay\fontsize{12}{12}\selectfont\color{popcream}\thepopsec};%
  \hspace{8pt}{\popdisplay\fontsize{15}{17}\selectfont\color{popink}\MakeUppercase{#1}}\par
  \vspace{4pt}\noindent{\color{poporange}\rule{36pt}{3.6pt}}
  \end{minipage}\par\nopagebreak\vspace{8pt}
}
\newcounter{popsub}
\newcommand{\courssub}[1]{%
  \stepcounter{popsub}%
  \par\vspace{9pt}\noindent{\popxbold\fontsize{11.5}{13}\selectfont\color{popink}{\color{poporange}\thepopsec.\thepopsub}\hspace{6pt}#1}\par\nopagebreak\vspace{4pt}
}

% ============================================================
% Boîtes pédagogiques — même recette : fond blanc, bordure épaisse colorée,
% ombre dure encre, badge plein. Titre optionnel [#1].
% ============================================================
\tcbset{popbase/.style={breakable,enhanced,colback=white,boxrule=2.3pt,arc=9pt,
  shadow={3pt}{-3pt}{0pt}{popink},left=14pt,right=14pt,top=11pt,bottom=11pt,before skip=11pt,after skip=15pt,
  fontupper=\popsemi\fontsize{10.6}{14.5}\selectfont\color{popink},
  fontlower=\popsemi\fontsize{10.6}{14.5}\selectfont\color{popink}}}
% Coche dessinée (le glyphe ✓ n'existe pas dans les polices du document)
\newcommand{\popcheck}[1]{\tikz[baseline=-0.5ex]\draw[#1,line width=1.6pt,line cap=round,line join=round](-0.13,0.0)--(-0.04,-0.09)--(0.13,0.1);}
\providecommand{\coche}{\popcheck{popgreen}}
\providecommand{\resultat}[1]{\par\smallskip\noindent\fcolorbox{popgreen}{popgreenL}{\parbox{\dimexpr\linewidth-2\fboxsep-2\fboxrule\relax}{\textbf{Résultat.} #1}}\par\smallskip}
\newcommand{\popboxhead}[3]{\poppill{#1}{#2}{white}\ifx\relax#3\relax\else\hspace{6pt}{\popxbold\fontsize{10.6}{12}\selectfont\color{popink}#3}\fi\par\vspace{7pt}}

\newtcolorbox{aretenir}[1][]{popbase,colframe=popgreen,before upper={\popboxhead{À retenir}{popgreen}{#1}}}
\newtcolorbox{exemplebox}[1][]{popbase,colframe=poporange,before upper={\popboxhead{Exemple}{poporange}{#1}}}
\newtcolorbox{definition}[1][]{popbase,colframe=popblue,before upper={\popboxhead{Définition}{popblue}{#1}}}
\newtcolorbox{propriete}[1][]{popbase,colframe=poppurple,before upper={\popboxhead{Propriété}{poppurple}{#1}}}
\newtcolorbox{methode}[1][]{popbase,colframe=popgold,before upper={\popboxhead{Méthode}{popgold}{#1}}}
\newtcolorbox{attention}[1][]{popbase,colframe=poppink,before upper={\popboxhead{Attention}{poppink}{#1}}}
\newtcolorbox{experience}[1][]{popbase,colframe=popblue,colback=popblueL,before upper={\popboxhead{Expérience}{popblue}{#1}}}
% « Le savais-tu ? » : carte violette pleine (contexte, culture, Cameroun)
\newtcolorbox{savaistu}[1][]{popbase,colback=poppurple,colframe=popink,
  fontupper=\popsemi\fontsize{10.4}{14.2}\selectfont\color{white},
  before upper={\poppill{Le savais-tu\,?}{white}{poppurple}\ifx\relax#1\relax\else\hspace{6pt}{\popxbold\color{white}#1}\fi\par\vspace{7pt}}}
% Exercice résolu : énoncé en haut, \tcblower puis la solution sur fond crème
\newcounter{popexo}\newcounter{popexr}
\newtcolorbox{exoresolu}[1][]{popbase,colframe=poporange,colbacklower=popcream,
  segmentation style={popink,line width=1.2pt,dashed},
  before upper={\stepcounter{popexr}\popboxhead{Exercice résolu \thepopexr}{poporange}{#1}},
  before lower={\poppill{Solution}{popink}{popcream}\par\vspace{6pt}}}
% Exercice (non résolu en place) — la correction est dans la partie Corrigés
\newtcolorbox{exercice}[1][]{popbase,colframe=popink,
  before upper={\stepcounter{popexo}\popboxhead{Exercice \thepopexo}{popink}{#1}}}
% Corrigé
\newtcolorbox{corrige}[1][]{popbase,colframe=popgreen,colback=popgreenL,
  before upper={\popboxhead{Corrigé}{popgreen}{#1}}}
% Objectifs / prérequis / bilan
\newtcolorbox{objectifs}{popbase,colframe=popink,colback=poporangeL,
  before upper={\popboxhead{Objectifs de la leçon}{popink}{}}}
\newtcolorbox{prerequis}{popbase,colframe=popink,colback=popblueL,
  before upper={\popboxhead{Prérequis}{popblue}{}}}
\newtcolorbox{bilan}{popbase,colframe=popink,colback=white,boxrule=2.8pt,
  before upper={\popboxhead{Fiche bilan}{poporange}{L'essentiel à connaître pour l'examen}}}
% Figure encadrée avec légende
\newtcolorbox{popfigure}[1][]{enhanced,breakable=false,colback=white,colframe=popline,boxrule=1.4pt,arc=8pt,
  left=10pt,right=10pt,top=10pt,bottom=8pt,before skip=10pt,after skip=12pt,halign=center,
  lower separated=false,halign lower=center,
  fontlower=\popsemi\fontsize{9}{11.5}\selectfont\color{popmuted},
  #1}
\newcommand{\legende}[1]{\par\vspace{6pt}{\popsemi\fontsize{9}{11.5}\selectfont\color{popmuted}#1}}

% ============================================================
% Signature OnBuch+
% ============================================================
\newcommand{\poplogoinline}[1]{{\poplogo\fontsize{#1}{#1}\selectfont\color{popink}OnBuch\color{poporange}+}}
\newcommand{\signatureonbuch}{%
  \begin{tcolorbox}[enhanced,colback=popink,colframe=popink,boxrule=2.2pt,arc=10pt,
      drop shadow=poporange,left=14pt,right=14pt,top=10pt,bottom=10pt,before skip=0pt,after skip=0pt]
    \tikz[baseline=-0.7ex]\node[circle,fill=popgreen,draw=popcream,line width=1.3pt,minimum size=22pt,inner sep=0]{\popcheck{white}};%
    \hspace{9pt}\begin{minipage}[c]{0.62\linewidth}
      {\popxbold\fontsize{12}{13}\selectfont\color{popcream}Signé\ \textbullet\ {\poplogo OnBuch\color{poporange}+}}\par\vspace{2pt}
      {\raggedright\popsemi\fontsize{8}{10.5}\selectfont\color{popline}\MakeUppercase{Équipe pédagogique OnBuch+}\\\MakeUppercase{Conforme au programme officiel MINESEC}\par}
    \end{minipage}\hfill
    {\poplogo\fontsize{22}{22}\selectfont\color{popcream}OnBuch\color{poporange}+}
  \end{tcolorbox}%
}

% ============================================================
% Page de garde
% #1 titre · #2 matière · #3 niveau · #4 module · #5 n° leçon · #6 durée ·
% #7 description / objectifs (texte ou itemize)
% ============================================================
\newcommand{\pagegarde}[7]{%
  \thispagestyle{empty}
  \newgeometry{a4paper,margin=1.7cm,top=1.6cm,bottom=1.4cm}%
  \begin{tikzpicture}[remember picture,overlay]
    \fill[poporange!18] ([xshift=-3.2cm,yshift=-3.6cm]current page.north east) circle (4.6cm);
    \fill[poppurple!14] ([xshift=2.4cm,yshift=7.6cm]current page.south west) circle (3.8cm);
  \end{tikzpicture}%
  {\poplogo\fontsize{30}{30}\selectfont\color{popink}OnBuch\color{poporange}+}\hfill
  \raisebox{0.6ex}{\poppill{Cours signé \textbullet\ Premium}{popink}{popcream}}\par
  \vspace{1pt}\popsquiggle{poppurple}\par
  \vspace{8pt}
  \begin{tcolorbox}[enhanced,colback=poporange,colframe=popink,boxrule=2.8pt,arc=13pt,
      drop shadow=popink,left=18pt,right=18pt,top=12pt,bottom=13pt,after skip=10pt]
    \poppill{#2 \textbullet\ #3}{popink}{popcream}\hspace{5pt}\poppill{#4}{white}{popink}\par\vspace{8pt}
    {\popdisplay\fontsize{14}{14}\selectfont\color{popink}LEÇON #5}\par\vspace{4pt}
    {\raggedright\hyphenpenalty=10000\exhyphenpenalty=10000\popdisplay\fontsize{25}{27}\selectfont\color{popcream}\MakeUppercase{#1}\par}
    \vspace{8pt}
    {\popxbold\fontsize{9.5}{9.5}\selectfont\color{popink}\MakeUppercase{Durée conseillée : #6}}
  \end{tcolorbox}
  \begin{tcolorbox}[enhanced,colback=white,colframe=popink,boxrule=2.2pt,arc=10pt,
      drop shadow=popink,left=16pt,right=16pt,top=10pt,bottom=10pt,after skip=8pt]
    \poppill{Dans cette leçon}{poppurple}{white}\par\vspace{5pt}
    \popsemi\fontsize{9.3}{12.6}\selectfont\color{popink}#7
  \end{tcolorbox}
  \noindent\begin{minipage}[t]{0.487\linewidth}
    \begin{tcolorbox}[enhanced,colback=popgreen,colframe=popink,boxrule=2.2pt,arc=10pt,
        drop shadow=popink,left=11pt,right=11pt,top=10pt,bottom=10pt,height=3.1cm,valign=center]
      \tcbox[colback=white,colframe=popink,boxrule=1.3pt,arc=5pt,boxsep=0pt,nobeforeafter,
        left=3pt,right=3pt,top=3pt,bottom=3pt,valign=center]{\qrcode[height=1.9cm]{https://play.google.com/store/apps/details?id=com.luvvix.onbuch&hl=fr}}%
      \hspace{8pt}\begin{minipage}[c]{0.5\linewidth}
        {\popdisplay\fontsize{11}{12.5}\selectfont\color{white}\MakeUppercase{Télécharge OnBuch}}\par\vspace{4pt}
        {\raggedright\popsemi\fontsize{8.4}{11}\selectfont\color{white}Gratuit \textbullet\ Google Play\\Tuteur IA, annales, concours, résultats\par}
      \end{minipage}
    \end{tcolorbox}
  \end{minipage}\hfill
  \begin{minipage}[t]{0.487\linewidth}
    \begin{tcolorbox}[enhanced,colback=poppurple,colframe=popink,boxrule=2.2pt,arc=10pt,
        drop shadow=popink,left=11pt,right=11pt,top=10pt,bottom=10pt,height=3.1cm,valign=center]
      \tikz[baseline=-0.7ex]\node[circle,fill=white,draw=popink,line width=1.3pt,minimum size=1.6cm,inner sep=0]{\popdisplay\fontsize{24}{24}\selectfont\color{poppurple}+};%
      \hspace{8pt}\begin{minipage}[c]{0.6\linewidth}
        {\popdisplay\fontsize{11}{12.5}\selectfont\color{white}\MakeUppercase{Passe à OnBuch+}}\par\vspace{4pt}
        {\raggedright\popsemi\fontsize{8.4}{11}\selectfont\color{white}Tous les cours signés, Tuteur IA illimité, fiches PDF — onglet OnBuch+ de l'app\par}
      \end{minipage}
    \end{tcolorbox}
  \end{minipage}
  \vspace{8pt}
  \popcontenu
  \vfill
  \signatureonbuch
  \restoregeometry
  \newpage
}

% Rangée « Ce cours contient » (page de garde)
\newcommand{\poptile}[3]{% #1 couleur #2 gros texte #3 légende
  \begin{tcolorbox}[enhanced,colback=#1,colframe=popink,boxrule=2pt,arc=9pt,shadow={2.5pt}{-2.5pt}{0pt}{popink},
      left=6pt,right=6pt,top=9pt,bottom=9pt,halign=center,valign=center,height=1.75cm,width=\linewidth,nobeforeafter]
    {\popdisplay\fontsize{12.5}{13}\selectfont\color{white}\MakeUppercase{#2}}\par\vspace{3pt}
    {\centering\popsemi\fontsize{7.8}{9.5}\selectfont\color{white}#3\par}
  \end{tcolorbox}}
\newcommand{\popcontenu}{%
  \par\noindent{\popxboldls\fontsize{7.5}{7.5}\selectfont\color{popmuted}CE COURS SIGNÉ CONTIENT}\par\vspace{7pt}
  \noindent\begin{minipage}[t]{0.235\linewidth}\poptile{popblue}{Cours}{complet, illustré, pas à pas}\end{minipage}\hfill
  \begin{minipage}[t]{0.235\linewidth}\poptile{poporange}{Méthodes}{et exercices résolus}\end{minipage}\hfill
  \begin{minipage}[t]{0.235\linewidth}\poptile{poppink}{Exercices}{type examen + corrigés}\end{minipage}\hfill
  \begin{minipage}[t]{0.235\linewidth}\poptile{popgreen}{Fiche bilan}{l'essentiel à retenir}\end{minipage}\par}

% Sommaire
\newtcolorbox{sommaire}{enhanced,colback=white,colframe=popink,boxrule=2.2pt,arc=10pt,
  drop shadow=popink,left=18pt,right=18pt,top=13pt,bottom=13pt,before skip=6pt,after skip=20pt,
  before upper={\poppill{Sommaire}{poppurple}{white}\par\vspace{9pt}\popsemi\fontsize{10.6}{14.5}\selectfont\color{popink}}}

% Signature de fin de cours — poussée tout en bas de la dernière page
\newcommand{\findecours}{%
  \par\vfill\nopagebreak
  \begin{center}\popsquiggle{poporange}\end{center}\vspace{4pt}
  \signatureonbuch
}
\AtBeginDocument{\def\llbracket{\mathopen{[\mkern-3.5mu[}}\def\rrbracket{\mathclose{]\mkern-3.5mu]}}} % intervalles d'entiers (glyphe ⟦ absent de la police)
% Glyphes absents de Fira Math : redéfinis à partir de glyphes disponibles
\AtBeginDocument{%
  \def\setminus{\mathbin{\backslash}}\let\smallsetminus\setminus
  \def\vdots{\mathord{\vbox{\baselineskip3.2pt\lineskiplimit0pt\kern4pt\hbox{.}\hbox{.}\hbox{.}}}}%
  \def\ddots{\mathinner{\mkern1mu\raise6.5pt\hbox{.}\mkern2mu\raise3.6pt\hbox{.}\mkern2mu\raise0.7pt\hbox{.}\mkern1mu}}%
  \def\triangle{\mathord{\scalebox{0.9}{$\symup{\Delta}$}}}%
  \def\nabla{\mathord{\raisebox{\depth}{\scalebox{1}[-1]{$\symup{\Delta}$}}}}%
  \def\checkmark{\popcheck{popdark}}%
  \def\star{\mathbin{\ast}}\def\bigstar{\mathord{\ast}}%
  \def\diamond{\mathbin{\scalebox{0.6}{$\Diamond$}}}%
}

%% FIGFIX-BEGIN v4 — correctifs globaux des figures (docs/onbuch-generation/tools/figfix)
\usepackage{adjustbox}\usepackage{etoolbox}
% toute figure TikZ/pgfplots plus large que la ligne est réduite à la largeur disponible
\BeforeBeginEnvironment{tikzpicture}{\begin{adjustbox}{max width=\linewidth}}
\AfterEndEnvironment{tikzpicture}{\end{adjustbox}}
% légendes pgfplots lisibles par-dessus les courbes
\pgfplotsset{every axis/.append style={legend style={fill=white,fill opacity=0.92,text opacity=1,draw=black!15,inner sep=3pt}}}
% étiquettes d'arêtes (syntaxe edge["..."]) avec fond blanc
\tikzset{every edge quotes/.append style={fill=white,inner sep=1.2pt}}
%% FIGFIX-END
\begin{document}

\pagegarde{Interaction gravitationnelle et champ de gravitation}{Physique}{Tle E}{Module 2 : Mouvements et interactions — évolutions temporelles des systèmes mécaniques}{P02}{4 h}{Cette leçon établit la loi de l'attraction universelle de Newton et introduit la notion de champ de gravitation créé par un point matériel ou une sphère homogène. Elle aboutit à l'expression du champ de gravitation terrestre en fonction de l'altitude et à la distinction entre poids et force de gravitation.
\vspace{4pt}
\begin{multicols}{2}\small
\begin{itemize}
\item À la fin de cette leçon, je sais énoncer la loi de l'attraction universelle de Newton et identifier chaque grandeur
\item À la fin de cette leçon, je sais représenter les forces gravitationnelles s'exerçant entre deux corps
\item À la fin de cette leçon, je sais calculer l'intensité d'une force gravitationnelle et la comparer à d'autres forces
\item À la fin de cette leçon, je sais définir le vecteur champ de gravitation créé par un point matériel et le représenter
\item À la fin de cette leçon, je sais utiliser le théorème de la sphère homogène pour exprimer le champ à l'extérieur d'un astre
\item À la fin de cette leçon, je sais tracer et interpréter des lignes de champ de gravitation
\end{itemize}\end{multicols}}

\begin{sommaire}
\begin{multicols}{2}
\begin{enumerate}
\item La loi de l'attraction universelle de Newton
\item Champ de gravitation créé par un point matériel
\item Champ créé par une sphère homogène
\item Champ de gravitation terrestre et variation avec l'altitude
\item Champ uniforme au voisinage du sol ; poids et force de gravitation
\item Activité d'intégration
\item Méthodes et exercices résolus
\item Exercices
\item Corrigés des exercices
\item Fiche bilan
\end{enumerate}
\end{multicols}
\end{sommaire}

\begin{objectifs}
\begin{itemize}
\item À la fin de cette leçon, je sais énoncer la loi de l'attraction universelle de Newton et identifier chaque grandeur
\item À la fin de cette leçon, je sais représenter les forces gravitationnelles s'exerçant entre deux corps
\item À la fin de cette leçon, je sais calculer l'intensité d'une force gravitationnelle et la comparer à d'autres forces
\item À la fin de cette leçon, je sais définir le vecteur champ de gravitation créé par un point matériel et le représenter
\item À la fin de cette leçon, je sais utiliser le théorème de la sphère homogène pour exprimer le champ à l'extérieur d'un astre
\item À la fin de cette leçon, je sais tracer et interpréter des lignes de champ de gravitation
\item À la fin de cette leçon, je sais établir et utiliser l'expression g(h) = g0 R²/(R+h)²
\item À la fin de cette leçon, je sais justifier que le champ de gravitation est uniforme au voisinage du sol
\item À la fin de cette leçon, je sais distinguer le poids d'un corps de la force de gravitation qu'il subit
\end{itemize}
\end{objectifs}

\begin{prerequis}
\begin{itemize}
\item Forces : représentation vectorielle, principe des actions réciproques (3e loi de Newton) — classe de Première
\item Notion de vecteur : composantes, norme, vecteur unitaire
\item Poids P = m g et pesanteur terrestre g ≈ 9,8 N/kg (classes de collège et Seconde)
\item Puissances de 10, notation scientifique et calcul littéral
\item Modèle du point matériel et repérage d'un point dans l'espace
\end{itemize}
\end{prerequis}

\newpage

\coursec{La loi de l'attraction universelle de Newton}

\courssub{Accroche et problématique}

Le 6 août 2012, après un voyage de huit mois et demi dans l'espace, le robot \emph{Curiosity} se pose en douceur sur le sol martien. Pour guider cette prouesse technologique sur des centaines de millions de kilomètres, les ingénieurs de la NASA n'ont pas eu besoin de nouvelles physiques exotiques : ils ont utilisé la même loi fondamentale qui explique pourquoi un fruit de manguier tombe inéluctablement vers le sol à Douala, pourquoi les marées montent et descendent sur les côtes de Kribi, et pourquoi la Lune, attirée par la Terre, ne s'écrase pas sur elle mais décrit une orbite quasi circulaire depuis des milliards d'années.

Comment une même interaction, la \cle{gravitation}, peut-elle régir à la fois la chute d'une pomme, le mouvement des satellites artificiels (comme ceux de MTN ou CAMTEL en orbite géostationnaire) et la danse des planètes autour du Soleil ? Pourquoi la pesanteur \emph{g} vaut-elle environ 9,81 N·kg⁻¹ au niveau de la mer à Douala, mais devient-elle légèrement plus faible au sommet du mont Cameroun (4 095 m d'altitude) ? Cette première section pose les bases universelles de cette interaction : la loi de Newton, sa formulation vectorielle, la constante \cle{G} et ses conséquences immédiates.

\courssub{Mise en évidence qualitative : une interaction universelle}

Depuis l'Antiquité, on distingue le \emph{monde sublunaire} (où les corps tombent) du \emph{monde supralunaire} (où les astres tournent « parfaitement »). La génie de Newton (1687, \emph{Philosophiæ Naturalis Principia Mathematica}) a été d'unifier ces deux domaines en une seule hypothèse audacieuse : \textbf{tout corps massique attire tout autre corps massique}.

Cette attraction se manifeste à trois échelles emblématiques :
\begin{itemize}
    \item \textbf{Échelle locale :} La chute des corps. Lâchez un caillou à Yaoundé, il tombe vers le centre de la Terre. La Terre l'attire, et — point crucial souvent oublié — le caillou attire aussi la Terre (principe des actions réciproques).
    \item \textbf{Échelle planétaire :} Les marées. L'attraction de la Lune (et du Soleil) déforme les océans terrestres. À Kribi, la mer monte et descend deux fois par jour sous l'effet de cette force à distance.
    \item \textbf{Échelle cosmique :} Le mouvement de la Lune. Elle est attirée par la Terre, mais sa vitesse tangentielle fait qu'elle « tombe sans cesse à côté de la Terre », décrivant une orbite stable. C'est la même loi qui retient les satellites de télécommunication (MTN, Orange, CAMTEL) en orbite géostationnaire à 36 000 km d'altitude.
\end{itemize}

\begin{savaistu}[Newton et la pomme : mythe ou réalité ?]
L'histoire de la pomme qui tombe sur la tête de Newton dans son jardin de Woolsthorpe (1666) est en grande partie une construction a posteriori (racontée par Voltaire). Newton lui-même a écrit avoir été inspiré en voyant une pomme tomber, se demandant : « Pourquoi cette pomme tombe-t-elle toujours perpendiculairement au sol ? Pourquoi ne va-t-elle pas latéralement ou vers le haut ? » L'intuition majeure fut de réaliser que \emph{la même force} qui tire la pomme vers le sol s'étend jusqu'à la Lune et la retient en orbite. Il a fallu attendre 1687 pour que la loi soit publiée sous sa forme mathématique aboutie.
\end{savaistu}

\courssub{Énoncé de la loi de l'attraction universelle}

Considérons deux \cle{points matériels} A et B, de masses respectives $m_A$ et $m_B$ (en kilogrammes), séparés par une distance $d = AB$ (en mètres). L'expérience et l'observation astronomique montrent que :
\begin{enumerate}
    \item A exerce sur B une force \textbf{attractive} $\vec{F}_{A/B}$.
    \item B exerce sur A une force \textbf{attractive} $\vec{F}_{B/A}$.
    \item Ces deux forces ont même direction (la droite $(AB)$), même sens (l'un vers l'autre), même intensité.
\end{enumerate}

\begin{definition}[Loi de l'attraction universelle (Newton, 1687)]
Deux points matériels de masses $m_A$ et $m_B$, séparés par une distance $d$, s'attirent mutuellement avec une force d'intensité :
\[
F = G\,\frac{m_A\,m_B}{d^2}
\]
où $G = \SI{6,67e-11}{N.m^2.kg^{-2}}$ est la \cle{constante de gravitation universelle}.
\end{definition}

\begin{popfigure}
\begin{tikzpicture}[scale=1.2, >=Stealth]
    % Points
    \node[circle, fill=popblue, inner sep=2pt, label=above:$A\ (m_A)$] (A) at (0,0) {};
    \node[circle, fill=poporange, inner sep=2pt, label=above:$B\ (m_B)$] (B) at (5,0) {};
    
    % Distance
    \draw[popmuted, <->] (A) -- (B) node[midway, above=2pt] {$d = AB$};
    
    % Vecteur unitaire u_AB (de A vers B)
    \draw[popblue, thick, ->] (A) -- ++(1,0) node[midway, above] {$\vec{u}_{AB}$};
    
    % Force sur B (exercée par A) : attractive, vers A
    \draw[poporange, very thick, ->] (B) -- ++(-1.5,0) node[midway, above] {$\vec{F}_{A/B}$};
    
    % Force sur A (exercée par B) : attractive, vers B
    \draw[popblue, very thick, ->] (A) -- ++(1.5,0) node[midway, above] {$\vec{F}_{B/A}$};
    
    % Legende vecteur unitaire
    \node[align=center] at (2.5, -1.2) {$\vec{u}_{AB} = \dfrac{\vec{AB}}{AB}$\\(vecteur unitaire de A vers B)};
\end{tikzpicture}
\legende{Représentation des forces gravitationnelles mutuelles entre deux points matériels A et B. Les forces sont attractives, colinéaires à la droite (AB), de même norme.}
\end{popfigure}

\courssub{Expression vectorielle et scalaire}

Pour travailler rigoureusement en mécanique, nous avons besoin de l'expression vectorielle. Introduisons le \cle{vecteur unitaire} $\vec{u}_{AB}$ dirigé de A vers B :
\[
\vec{u}_{AB} = \frac{\overrightarrow{AB}}{AB} = \frac{\overrightarrow{AB}}{d}
\]

La force que A exerce sur B est attractive : elle est dirigée de B vers A, c'est-à-dire \emph{opposée} à $\vec{u}_{AB}$. D'où l'expression vectorielle fondamentale :

\begin{propriete}[Expression vectorielle de la force gravitationnelle]
\[
\boxed{\vec{F}_{A/B} = -\,G\,\frac{m_A\,m_B}{d^2}\,\vec{u}_{AB}}
\]
Symétriquement : $\vec{F}_{B/A} = -\,G\,\frac{m_A\,m_B}{d^2}\,\vec{u}_{BA} = +\,G\,\frac{m_A\,m_B}{d^2}\,\vec{u}_{AB}$.
\end{propriete}

\begin{attention}[Pièges classiques sur l'expression vectorielle]
\begin{itemize}
    \item \textbf{Le signe moins :} Il est \textbf{obligatoire}. Il traduit l'attraction (force opposée au vecteur unitaire qui pointe de la source vers la cible). L'oublier donne une force répulsive (fausse).
    \item \textbf{Le vecteur unitaire :} $\vec{u}_{AB}$ va de A vers B. La force $\vec{F}_{A/B}$ s'applique \emph{en B} et pointe vers A.
    \item \textbf{Notations :} $\vec{F}_{A/B}$ se lit « force exercée \emph{par} A \emph{sur} B ». L'ordre des indices compte pour le signe vectoriel.
\end{itemize}
\end{attention}

L'expression scalaire de l'intensité (toujours positive) est :
\[
F = G\,\frac{m_A\,m_B}{d^2}
\]

\courssub{La constante de gravitation universelle G}

\begin{definition}[Constante de gravitation universelle]
\[
G = \SI{6,67e-11}{N.m^2.kg^{-2}} \quad (\text{valeur CODATA 2018 : } \SI{6,67430e-11}{N.m^2.kg^{-2}})
\]
\end{definition}

\noindent \textbf{Signification physique :} $G$ quantifie l'intensité intrinsèque de l'interaction gravitationnelle. C'est une \emph{constante universelle} : elle a la même valeur partout dans l'Univers, à tout instant, quelles que soient les masses considérées (électrons, pommes, planètes, trous noirs). C'est l'une des constantes les plus difficiles à mesurer avec précision car la gravitation est l'interaction la plus faible (environ $10^{36}$ fois plus faible que l'électromagnétisme).

\begin{savaistu}[La mesure de G : l'expérience de Cavendish (1798)]
Newton a énoncé sa loi en 1687, mais il ne connaissait pas la valeur de $G$ ! Il a fallu attendre 1798 pour que le physicien britannique \textbf{Henry Cavendish} la mesure avec une \emph{balance de torsion}. Il a suspendu une barre horizontale portant deux petites boules de plomb à un fil de quartz. Deux grosses boules de plomb placées à proximité attirent les petites, faisant pivoter la barre. En mesurant l'angle de torsion et la constante du fil, il a déduit $G$ et, du même coup, « pesé la Terre » (déduit sa masse $M_T$). Au Cameroun, ce type d'expérience de précision est enseigné en TP de physique à l'université (ENS Yaoundé, Faculté des Sciences) pour illustrer la mesure des constantes fondamentales.
\end{savaistu}

\courssub{Caractéristiques des forces gravitationnelles}

Rappelons les quatre caractéristiques d'une force (nécessaires pour tout bilan de forces au Bac) :

\begin{center}
\begin{tabularx}{\linewidth}{|l|X|}\hline
\textbf{Caractéristique} & \textbf{Force gravitationnelle $\vec{F}_{A/B}$} \\ \hline
\textbf{Direction} & Droite $(AB)$ joignant les deux centres de masse. \\ \hline
\textbf{Sens} & \textbf{Attractif} : de B vers A (pour $\vec{F}_{A/B}$). \\ \hline
\textbf{Point d'application} & Centre de masse du corps \emph{qui subit} la force (point B pour $\vec{F}_{A/B}$). \\ \hline
\textbf{Norme (intensité)} & $F = G\,\dfrac{m_A\,m_B}{d^2}$ (en newtons, N). \\ \hline
\end{tabularx}
\end{center}

\courssub{Vérification du principe des actions réciproques}

C'est une conséquence immédiate de l'expression vectorielle, souvent demandée en rédaction de Bac.

\begin{propriete}[Principe des actions réciproques (3\textsuperscript{e} loi de Newton)]
Les forces gravitationnelles mutuelles vérifient :
\[
\vec{F}_{A/B} = -\,\vec{F}_{B/A}
\]
\textbf{Preuve :} $\vec{F}_{A/B} = -G\frac{m_A m_B}{d^2}\vec{u}_{AB}$ et $\vec{F}_{B/A} = -G\frac{m_A m_B}{d^2}\vec{u}_{BA}$. Or $\vec{u}_{BA} = -\vec{u}_{AB}$. Donc $\vec{F}_{B/A} = +G\frac{m_A m_B}{d^2}\vec{u}_{AB} = -\vec{F}_{A/B}$. \qed
\end{propriete}

\begin{attention}[Rédaction attendue au Bac]
Ne vous contentez pas d'écrire « d'après le principe des actions réciproques ». \textbf{Justifiez} par l'expression vectorielle : « Les deux forces ont même direction (droite AB), même norme ($G m_A m_B/d^2$), sens opposés (l'une selon $\vec{u}_{AB}$, l'autre selon $\vec{u}_{BA} = -\vec{u}_{AB}$), donc $\vec{F}_{A/B} = -\vec{F}_{B/A}$. »
\end{attention}

\courssub{Ordres de grandeur : trois échelles comparées}

La force gravitationnelle est proportionnelle au produit des masses et inversement proportionnelle au \emph{carré} de la distance. Voyons ce que cela donne concrètement.

\begin{exemplebox}[Force entre deux élèves à 1 m]
Deux élèves de masse $m_1 = m_2 = \SI{60}{kg}$ sont assis à $d = \SI{1}{m}$ l'un de l'autre.
\[
F = G\,\frac{m_1 m_2}{d^2} = \num{6,67e-11} \times \frac{60 \times 60}{1^2} = \SI{2,4e-7}{N}
\]
C'est \emph{extrêmement faible} : équivalent au poids d'un grain de sable de $\approx \SI{24}{\micro g}$. C'est pourquoi nous ne ressentons pas l'attraction de nos camarades en classe !
\end{exemplebox}

\begin{exemplebox}[Force Terre–Lune]
\begin{itemize}
    \item Masse de la Terre : $M_T = \SI{6,0e24}{kg}$
    \item Masse de la Lune : $M_L = \SI{7,3e22}{kg}$
    \item Distance moyenne Terre–Lune : $d_{TL} = \SI{3,84e8}{m}$
\end{itemize}
\[
F_{T/L} = \num{6,67e-11} \times \frac{(\num{6,0e24})(\num{7,3e22})}{(\num{3,84e8})^2}
= \num{6,67e-11} \times \frac{\num{4,38e47}}{\num{1,47e17}}
\approx \SI{1,98e20}{N}
\]
Arrondi : $\boxed{F_{T/L} \approx \SI{2,0e20}{N}}$.
C'est cette force colossale qui maintient la Lune en orbite (force centripète).
\end{exemplebox}

\begin{exemplebox}[Force Soleil–Terre]
\begin{itemize}
    \item Masse du Soleil : $M_S = \SI{2,0e30}{kg}$
    \item Distance moyenne Soleil–Terre (1 UA) : $d_{ST} = \SI{1,50e11}{m}$
\end{itemize}
\[
F_{S/T} = \num{6,67e-11} \times \frac{(\num{2,0e30})(\num{6,0e24})}{(\num{1,50e11})^2}
= \num{6,67e-11} \times \frac{\num{1,2e55}}{\num{2,25e22}}
\approx \SI{3,56e22}{N}
\]
Arrondi : $\boxed{F_{S/T} \approx \SI{3,6e22}{N}}$.
Le Soleil attire la Terre environ 180 fois plus fort que la Lune ne le fait. Pourtant, les marées sont principalement dues à la Lune ! (Car les marées dépendent du \emph{gradient} du champ, $\propto 1/d^3$, pas de l'intensité du champ $\propto 1/d^2$).
\end{exemplebox}

\begin{popfigure}
\begin{tikzpicture}[scale=0.9, every node/.style={font=\small}]
    % Echelle 1 : Deux élèves
    \begin{scope}[yshift=0]
        \node[align=center, popblue] at (0,1.5) {\textbf{Échelle humaine}};
        \node[circle, fill=popblue, inner sep=3pt, label=below:60 kg] (E1) at (-1,0) {};
        \node[circle, fill=popblue, inner sep=3pt, label=below:60 kg] (E2) at (1,0) {};
        \draw[popmuted, <->] (E1) -- (E2) node[midway, above] {1 m};
        \draw[poporange, thick, <->] (E1) -- (E2);
        \node[align=center, poporange] at (0,-1) {$F \approx \SI{2,4e-7}{N}$\\(poids d'un grain de sable)};
    \end{scope}
    
    % Echelle 2 : Terre-Lune
    \begin{scope}[xshift=7cm, yshift=0]
        \node[align=center, popblue] at (0,1.5) {\textbf{Échelle Terre--Lune}};
        \node[circle, fill=popblue!80!black, minimum size=1.2cm, label={below:$M_T = 6,0\times10^{24}$} kg] (T) at (-1.5,0) {};
        \node[circle, fill=popgray, minimum size=0.4cm, label={below:$M_L = 7,3\times10^{22}$} kg] (L) at (2.5,0) {};
        \draw[popmuted, <->] (T.east) -- (L.west) node[midway, above] {$3,84\times10^8$ m};
        \draw[poporange, thick, <->] (T) -- (L);
        \node[align=center, poporange] at (0.5,-1.2) {$F \approx \SI{2,0e20}{N}$};
    \end{scope}
    
    % Echelle 3 : Soleil-Terre
    \begin{scope}[xshift=14cm, yshift=0]
        \node[align=center, popblue] at (0,1.5) {\textbf{Échelle Soleil--Terre}};
        \node[circle, fill=poporange!80!yellow, minimum size=2cm, label={below:$M_S = 2,0\times10^{30}$} kg] (S) at (-2,0) {};
        \node[circle, fill=popblue!80!black, minimum size=0.5cm, label={below:$M_T = 6,0\times10^{24}$} kg] (T2) at (2,0) {};
        \draw[popmuted, <->] (S.east) -- (T2.west) node[midway, above] {$1,50\times10^{11}$ m};
        \draw[poporange, thick, <->] (S) -- (T2);
        \node[align=center, poporange] at (0,-1.3) {$F \approx \SI{3,6e22}{N}$\\(180$\times$ Terre-Lune)};
    \end{scope}
\end{tikzpicture}
\legende{Comparaison des forces gravitationnelles à trois échelles. Notez l'écart énorme (43 ordres de grandeur entre l'échelle humaine et l'échelle Soleil-Terre) et le rôle prépondérant des masses.}
\end{popfigure}

\courssub{Extension aux corps à symétrie sphérique}

Jusqu'ici, nous avons raisonné sur des \emph{points matériels}. Or, la Terre, la Lune, le Soleil, les planètes sont des sphères (quasi) homogènes. Un théorème majeur (démontré par Newton lui-même avec le calcul infinitésimal, admis en Terminale) permet d'étendre la loi :

\begin{propriete}[Théorème de Newton (admis en Terminale C)]
Le champ de gravitation créé par une \textbf{sphère homogène} (creuse ou pleine) en un point \textbf{extérieur} à cette sphère est \textbf{identique} à celui qui serait créé si toute la masse de la sphère était concentrée en son centre.
\end{propriete}

\noindent \textbf{Conséquence immédiate :} Pour calculer la force entre deux sphères homogènes (ex: Terre et Lune), la distance $d$ à utiliser dans la formule $F = G M_1 M_2 / d^2$ est la distance entre \textbf{leurs centres} (et non entre leurs surfaces).

\begin{attention}[Erreur fréquente : distance entre surfaces vs centres]
Si un énoncé donne le \emph{rayon} de la Terre $R_T = \SI{6370}{km}$ et l'\emph{altitude} d'un satellite $h = \SI{400}{km}$, la distance $d$ à utiliser est $d = R_T + h = \SI{6770}{km} = \SI{6,77e6}{m}$.
\textbf{Jamais} $d = h$ seul, et \textbf{jamais} $d = R_T$ seul (sauf si l'objet est \emph{à la surface}).
\end{attention}

\courssub{Exercice résolu : Application numérique complète}

\begin{exoresolu}[Satellite en orbite basse — Calcul de la force gravitationnelle]
Un satellite de télécommunication de masse $m = \SI{1200}{kg}$ (type satellite d'observation pour l'agriculture ou la météo au Cameroun) est placé sur une orbite circulaire à une altitude $h = \SI{400}{km}$ au-dessus de la surface terrestre.
Données : $M_T = \SI{6,0e24}{kg}$ ; $R_T = \SI{6370}{km}$ ; $G = \SI{6,67e-11}{N.m^2.kg^{-2}}$.
\begin{enumerate}
    \item Calculer la distance $d$ entre le centre de la Terre et le satellite.
    \item Calculer l'intensité $F$ de la force gravitationnelle exercée par la Terre sur le satellite.
    \item Comparer cette force au poids du satellite au sol ($P_0 = m g_0$ avec $g_0 = \SI{9,81}{N.kg^{-1}}$).
\end{enumerate}

\tcblower
\textbf{Solution détaillée}

\noindent \textbf{1. Distance centre-à-centre}
L'altitude $h$ se mesure depuis la surface. La distance au centre est :
\[
d = R_T + h = \SI{6370}{km} + \SI{400}{km} = \SI{6770}{km} = \SI{6,77e6}{m}
\]
\emph{(Attention : conversion impérative en mètres pour utiliser $G$ en SI).}

\noindent \textbf{2. Force gravitationnelle}
\[
F = G\,\frac{M_T\,m}{d^2}
= \num{6,67e-11} \times \frac{(\num{6,0e24})(\num{1200})}{(\num{6,77e6})^2}
\]
Calcul du dénominateur : $d^2 = (\num{6,77e6})^2 = \SI{4,58e13}{m^2}$.
Calcul du numérateur : $M_T m = \SI{7,2e27}{kg^2}$.
\[
F = \frac{\num{6,67e-11} \times \num{7,2e27}}{\num{4,58e13}}
= \frac{\num{4,80e17}}{\num{4,58e13}}
\approx \SI{1,05e4}{N}
\]
\[
\boxed{F \approx \SI{10,5}{kN}}
\]

\noindent \textbf{3. Comparaison au poids au sol}
Poids au sol : $P_0 = m g_0 = \num{1200} \times \num{9,81} = \SI{11772}{N} \approx \SI{11,8}{kN}$.
Rapport : $\dfrac{F}{P_0} = \dfrac{10,5}{11,8} \approx 0,89$.
\textbf{Interprétation :} À 400 km d'altitude (orbite de la Station Spatiale Internationale ISS), la gravité vaut encore \textbf{89\%} de sa valeur au sol ! Les astronautes ne sont pas en « apesanteur » parce que la gravité a disparu, mais parce qu'ils sont en \textbf{chute libre} permanente autour de la Terre (vitesse orbitale $\approx \SI{7,7}{km.s^{-1}}$). C'est un point conceptuel crucial pour la suite du programme (mouvement circulaire, satellites).
\end{exoresolu}

\courssub{Synthèse des savoir-faire pour le Bac}

\begin{aretenir}[Check-list Section 1 — Loi de Newton]
\begin{itemize}
    \item \textbf{Énoncer la loi} (scalaire et vectorielle) en définissant tous les symboles ($m_A, m_B, d, G, \vec{u}_{AB}$).
    \item \textbf{Représenter} les deux forces $\vec{F}_{A/B}$ et $\vec{F}_{B/A}$ sur un schéma : même droites, sens opposés, même longueur.
    \item \textbf{Calculer} une force : convertir \textbf{systématiquement} les distances en \textbf{mètres}, les masses en \textbf{kg}, utiliser $G = \SI{6,67e-11}{N.m^2.kg^{-2}}$.
    \item \textbf{Justifier} le principe des actions réciproques à partir de l'expression vectorielle ($\vec{u}_{BA} = -\vec{u}_{AB}$).
    \item \textbf{Appliquer} le théorème de Newton pour les sphères : $d = \text{distance entre centres}$.
    \item \textbf{Donner des ordres de grandeur} : force humaine $\sim 10^{-7}$ N ; Terre-Lune $\sim 10^{20}$ N ; Soleil-Terre $\sim 10^{22}$ N.
\end{itemize}
\end{aretenir}

\begin{methode}[Calculer une force gravitationnelle entre deux sphères]
\begin{enumerate}
    \item Identifier les deux masses $M_1, M_2$ (en kg).
    \item Déterminer la distance $d$ entre les \textbf{centres} (en m). Si on donne un rayon $R$ et une altitude $h$ : $d = R + h$.
    \item Appliquer la formule scalaire : $F = G \dfrac{M_1 M_2}{d^2}$.
    \item Vérifier l'homogénéité : $[G] = \SI{}{N.m^2.kg^{-2}}$, $[M_1 M_2] = \SI{}{kg^2}$, $[d^2] = \SI{}{m^2}$ $\Rightarrow$ $[F] = \SI{}{N}$. \checkmark
    \item Exprimer le résultat en notation scientifique avec 2 ou 3 chiffres significatifs et l'unité (N ou kN).
\end{enumerate}
\end{methode}

\courssub{Exercices d'entraînement}

\begin{exercice}[Forces à l'échelle du système solaire]
La masse de Jupiter est $M_J = \SI{1,90e27}{kg}$. Sa distance moyenne au Soleil est $d_{SJ} = \SI{7,78e11}{m}$.
Calculer l'intensité de la force gravitationnelle que le Soleil exerce sur Jupiter. Comparer à la force Soleil-Terre calculée plus haut ($\SI{3,6e22}{N}$). Jupiter est-elle plus ou moins fortement attirée que la Terre ? Justifier par les masses et distances.
\end{exercice}

\begin{exercice}[Altitude d'un satellite géostationnaire]
Un satellite de télécommunication (type MTN/CAMTEL) est géostationnaire à $h = \SI{36000}{km}$ d'altitude.
Données : $M_T = \SI{6,0e24}{kg}$, $R_T = \SI{6370}{km}$.
Calculer la force gravitationnelle exercée par la Terre sur ce satellite de masse $m = \SI{2000}{kg}$.
Comparer cette force à son poids au sol. Que représente le rapport $F/P_0$ physiquement (rappel : $g(h)/g_0$) ?
\end{exercice}

\begin{corrige}[Exercice 1 — Force Soleil-Jupiter]
$F_{S/J} = G \frac{M_S M_J}{d_{SJ}^2} = \num{6,67e-11} \times \frac{(\num{2,0e30})(\num{1,90e27})}{(\num{7,78e11})^2}$.
Numérateur : $\num{2,53e47}$. Dénominateur : $\num{6,05e23}$.
$F_{S/J} \approx \SI{4,18e23}{N}$.
Rapport avec Terre : $F_{S/J} / F_{S/T} \approx 11,6$.
Jupiter est $\approx 318$ fois plus massive que la Terre ($M_J/M_T \approx 318$) mais $\approx 5,2$ fois plus loin ($d_{SJ}/d_{ST} \approx 5,2$). La force varie comme $M/d^2$ : $318 / 5,2^2 \approx 318 / 27 \approx 11,8$. Cohérent. Jupiter est plus fortement attirée en valeur absolue, mais son accélération $g = F/M_J$ est plus faible.
\end{corrige}

\begin{corrige}[Exercice 2 — Satellite géostationnaire]
$d = R_T + h = \SI{6370}{km} + \SI{36000}{km} = \SI{42370}{km} = \SI{4,237e7}{m}$.
$F = \num{6,67e-11} \times \frac{(\num{6,0e24})(\num{2000})}{(\num{4,237e7})^2} = \frac{\num{8,00e17}}{\num{1,795e15}} \approx \SI{446}{N}$.
Poids au sol $P_0 = \num{2000} \times \num{9,81} = \SI{19620}{N}$.
Rapport $F/P_0 \approx 0,0227 \approx 2,3\%$.
Cela signifie que $g(h) \approx 0,023 \times g_0 \approx \SI{0,22}{N.kg^{-1}}$. La gravité est encore présente mais très affaiblie (facteur $\approx 1/44$). C'est cette force qui fournit l'accélération centripète nécessaire au mouvement circulaire de période 24 h.
\end{corrige}

\coursec{Champ de gravitation créé par un point matériel}

Dans la section précédente, nous avons établi la loi de l'attraction universelle : deux corps ponctuels de masses $m_S$ et $m$, distants de $r$, exercent l'un sur l'autre des forces attractives de même norme $F = G\,\dfrac{m_S\, m}{r^2}$. Cette formulation met les deux corps sur un pied d'égalité. Or, dans de nombreuses situations, l'un des deux corps domine largement : la Terre attire la pomme, le Soleil attire les planètes. Il devient alors naturel de changer de point de vue : plutôt que de parler d'une interaction entre deux corps, on dit que la masse $m_S$ \emph{modifie l'espace qui l'entoure} en y créant un \cle{champ de gravitation}, et que toute masse $m$ placée dans cet espace « ressent » ce champ sous forme d'une force. C'est ce nouveau langage, d'une puissance considérable, que nous construisons dans cette section.

\courssub{De la force au champ : l'idée de champ de gravitation}

\textbf{Une observation pour commencer.}\par
La Lune tourne autour de la Terre sans jamais la toucher ; une pomme tombe de l'arbre sans que rien ne la pousse visiblement. Comment la Terre agit-elle « à distance » ? Les physiciens ont répondu à cette question en introduisant une notion intermédiaire : la Terre remplit l'espace autour d'elle d'une \emph{propriété invisible} — le champ de gravitation — qui existe en tout point, que la pomme soit là ou non. La pomme ne fait que révéler ce champ en subissant une force.

Reprenons la loi de Newton vectorielle. Soit $S$ un point matériel de masse $m_S$ (la \cle{source}) et $M$ un point de l'espace où se trouve une masse $m$ (la \cle{masse test}). Notons $\vec{u}$ le vecteur unitaire dirigé de $S$ vers $M$ et $r = SM$. La force subie par $m$ s'écrit :
\[ \vec{F}_{S/m} = -G\,\dfrac{m_S\, m}{r^2}\,\vec{u} \]
Regroupons astucieusement les facteurs : mettons la masse test $m$ en évidence.
\[ \vec{F}_{S/m} = m \times \underbrace{\left( -G\,\dfrac{m_S}{r^2}\,\vec{u} \right)}_{\text{ne dépend que de la source et du point } M} \]
Le facteur entre parenthèses est remarquable : il ne contient \textbf{ni} la masse test $m$, \textbf{ni} aucune propriété de l'objet placé en $M$. Il ne dépend que de la masse source $m_S$ et de la position du point $M$ (à travers $r$ et $\vec{u}$). C'est une caractéristique de l'\emph{espace} lui-même, créée par la source. On lui donne un nom : le vecteur champ de gravitation.

\begin{definition}[Champ de gravitation créé par un point matériel]
Le \cle{vecteur champ de gravitation} créé en un point $M$ par un point matériel $S$ de masse $m_S$ est :
\[ \vec{g}(M) = -G\,\dfrac{m_S}{r^2}\,\vec{u} \]
où $r = SM$ est la distance de la source au point $M$, $\vec{u}$ le vecteur unitaire dirigé de $S$ vers $M$, et $G = \num{6,67}\times 10^{-11}\ \mathrm{N\,m^2\,kg^{-2}}$ la constante de gravitation universelle.
\end{definition}

\begin{propriete}[Force subie par une masse placée dans un champ]
Toute masse $m$ placée en un point $M$ où règne le champ de gravitation $\vec{g}(M)$ subit la force :
\[ \vec{F} = m\,\vec{g}(M) \]
\emph{Démonstration (exigible, deux lignes).} D'après la loi de l'attraction universelle, $\vec{F} = -G\dfrac{m_S\, m}{r^2}\vec{u}$. En factorisant $m$ : $\vec{F} = m\left(-G\dfrac{m_S}{r^2}\vec{u}\right) = m\,\vec{g}(M)$. $\square$
\end{propriete}

\begin{attention}[Le champ existe même sans masse test]
Le champ $\vec{g}(M)$ existe en $M$ \textbf{même si aucune masse ne s'y trouve}. La masse test $m$ ne sert qu'à le \emph{révéler} : sans elle, pas de force mesurable, mais le champ est bien là. Dire « la force crée le champ » est un contresens : c'est la source $m_S$ qui crée le champ ; la force n'apparaît que lorsqu'on y place une masse.
\end{attention}

\courssub{Caractéristiques du vecteur champ de gravitation}

\begin{aretenir}[Direction, sens, norme]
Le champ de gravitation créé en $M$ par un point matériel $S$ de masse $m_S$ a pour caractéristiques :
\begin{itemize}
\item \textbf{direction} : la droite $(SM)$ — on dit que le champ est \cle{radial} ;
\item \textbf{sens} : de $M$ vers $S$ — le champ est \cle{centripète}, c'est-à-dire toujours dirigé vers la source (le signe « $-$ » devant $\vec{u}$ traduit l'attraction) ;
\item \textbf{norme} : $g(M) = G\,\dfrac{m_S}{r^2}$, qui décroît comme l'inverse du carré de la distance.
\end{itemize}
\end{aretenir}

\textbf{Unité et double interprétation.}\par
D'après $\vec{F} = m\,\vec{g}$, on a $g = F/m$ : le champ s'exprime donc en \si{\newton\per\kilogram} ($\mathrm{N\,kg^{-1}}$). Mais l'analyse dimensionnelle révèle une équivalence profonde :
\[ [g] = \dfrac{[F]}{[m]} = \dfrac{\mathrm{kg\,m\,s^{-2}}}{\mathrm{kg}} = \mathrm{m\,s^{-2}} \]
Le champ de gravitation est donc \emph{homogène à une accélération}. Ce n'est pas un hasard : si une masse $m$ en $M$ n'est soumise qu'à la gravitation, la deuxième loi de Newton donne $m\vec{a} = m\,\vec{g}(M)$, soit $\vec{a} = \vec{g}(M)$. \textbf{Le champ de gravitation en un point est l'accélération qu'acquerrait tout objet lâché sans vitesse initiale en ce point}, quelle que soit sa masse. C'est pourquoi le champ de gravitation terrestre au sol n'est autre que l'accélération de pesanteur $g_0 \approx \SI{9,8}{\newton\per\kilogram}$. Le fait que $m$ se simplifie est la célèbre \emph{équivalence entre masse grave et masse inerte}, vérifiée — dit la légende — par Galilée lâchant des objets du haut de la tour de Pise.

\courssub{Représentation du champ autour de la source}

Pour représenter le champ en plusieurs points, on applique systématiquement la même démarche.

\begin{methode}[Représenter le vecteur champ de gravitation en un point M]
\begin{enumerate}
\item Tracer la demi-droite $[SM)$ joignant la source au point $M$ : c'est la direction du champ.
\item Orienter le vecteur \textbf{vers la source $S$} (champ centripète), en prenant $M$ comme origine du vecteur.
\item Calculer la norme $g(M) = G\,m_S/r^2$ avec $r = SM$.
\item Choisir une échelle (par exemple \SI{1}{cm} pour \SI{2e-3}{\newton\per\kilogram}) et tracer le vecteur à la longueur correspondante.
\end{enumerate}
Si l'on représente le champ en plusieurs points à des distances différentes, les longueurs des flèches doivent respecter la loi en $1/r^2$ : à distance double, la flèche est \textbf{quatre fois plus courte}.
\end{methode}

\begin{popfigure}
\begin{center}
\begin{tikzpicture}[scale=1.0]
  % Source
  \fill[popink] (0,0) circle (0.14);
  \node[below=2pt, popink] at (0,-0.05) {\footnotesize $S$ ($m_S$)};
  % Points proches r=1.4 : flèche longueur 1.0
  \draw[->, >=Stealth, line width=1.1pt, poporange] (1.4,0) -- (0.4,0) node[midway, above, poporange] {\footnotesize $\vec{g}_1$};
  \fill[popdark] (1.4,0) circle (0.05) node[right=1pt, popdark] {\footnotesize $M_1$};
  \draw[->, >=Stealth, line width=1.1pt, poporange] (0,1.4) -- (0,0.4);
  \fill[popdark] (0,1.4) circle (0.05) node[above=1pt, popdark] {\footnotesize $M_2$};
  \draw[->, >=Stealth, line width=1.1pt, poporange] (-1.4,0) -- (-0.4,0);
  \fill[popdark] (-1.4,0) circle (0.05) node[left=1pt, popdark] {\footnotesize $M_3$};
  \draw[->, >=Stealth, line width=1.1pt, poporange] (0,-1.4) -- (0,-0.4);
  \fill[popdark] (0,-1.4) circle (0.05) node[below=1pt, popdark] {\footnotesize $M_4$};
  % Points éloignés r=2.8 : flèche longueur 0.25
  \draw[->, >=Stealth, line width=1.1pt, popblue] (2.8,0) -- (2.55,0) node[midway, above, popblue] {\footnotesize $\vec{g}_5$};
  \fill[popdark] (2.8,0) circle (0.05) node[right=1pt, popdark] {\footnotesize $M_5$};
  \draw[->, >=Stealth, line width=1.1pt, popblue] (0,2.8) -- (0,2.55);
  \fill[popdark] (0,2.8) circle (0.05) node[above=1pt, popdark] {\footnotesize $M_6$};
  \draw[->, >=Stealth, line width=1.1pt, popblue] (-2.8,0) -- (-2.55,0);
  \fill[popdark] (-2.8,0) circle (0.05) node[left=1pt, popdark] {\footnotesize $M_7$};
  \draw[->, >=Stealth, line width=1.1pt, popblue] (0,-2.8) -- (0,-2.55);
  \fill[popdark] (0,-2.8) circle (0.05) node[below=1pt, popdark] {\footnotesize $M_8$};
  % Cercles de distances
  \draw[dashed, popmuted] (0,0) circle (1.4);
  \draw[dashed, popmuted] (0,0) circle (2.8);
  \node[popmuted] at (1.0,1.15) {\footnotesize $r$};
  \node[popmuted] at (2.0,2.3) {\footnotesize $2r$};
\end{tikzpicture}
\end{center}
\legende{Champ de gravitation créé par un point matériel $S$. Tous les vecteurs pointent vers $S$ (champ centripète). À distance double ($M_5$ à $M_8$), la norme est divisée par quatre : les flèches bleues sont quatre fois plus courtes que les flèches orange.}
\end{popfigure}

La figure met en évidence la \textbf{symétrie sphérique} du champ : à égale distance de $S$, la norme du champ est la même, seule la direction change. L'ensemble des vecteurs forme une « gerbe » convergeant vers la source. On visualisera cela autrement dans la suite de la leçon avec les \emph{lignes de champ}, qui sont ici des demi-droites radiales dirigées vers $S$.

\begin{attention}[Erreur fréquente : le sens du vecteur]
L'erreur la plus répandue consiste à orienter $\vec{g}(M)$ vers l'extérieur, « comme si la source repoussait ». C'est faux : la gravitation est \textbf{toujours attractive}. Le vecteur champ en $M$ pointe \textbf{toujours vers la source $S$}. Le signe « $-$ » dans $-G\,m_S/r^2\,\vec{u}$ n'est pas une décoration : il est là précisément pour retourner le vecteur $\vec{u}$ (qui va de $S$ vers $M$) en un vecteur allant de $M$ vers $S$. Au Baccalauréat, un vecteur champ mal orienté fait perdre la totalité des points de la question.
\end{attention}

\courssub{Décroissance en 1/r² : étude graphique}

La norme $g(r) = G\,m_S/r^2$ est une fonction hyperbolique de la distance : elle décroît très vite quand $r$ augmente. Doubler la distance divise le champ par quatre ; le multiplier par dix le divise par cent. Cette décroissance rapide explique pourquoi l'attraction d'un objet quotidien (un baobab, un camion) est totalement négligeable devant celle de la Terre : seules les masses astronomiques créent des champs mesurables.

\begin{popfigure}
\begin{center}
\begin{tikzpicture}
\begin{axis}[
  width=0.85\linewidth, height=6.2cm,
  xlabel={$r$ (en $10^6$ m)}, ylabel={$g(r)$ (en $\mathrm{N\,kg^{-1}}$)},
  xmin=0, xmax=10, ymin=0, ymax=10.5,
  xtick={2,4,6,8,10}, ytick={0,2,4,6,8,10},
  grid=both, grid style={popline!50},
  axis lines=left, label style={popdark}, tick label style={popdark}
]
  % G(r) = 40/r^2 (échelle arbitraire, g0=10 à r=2)
  \addplot[popcurve, domain=1.98:10, samples=150] {40/x^2};
  % Points de lecture
  \addplot[only marks, mark=*, mark size=2pt, poporange] coordinates {(2,10) (4,2.5)};
  \draw[dashed, popmuted] (axis cs:2,0) -- (axis cs:2,10) -- (axis cs:0,10);
  \draw[dashed, popmuted] (axis cs:4,0) -- (axis cs:4,2.5) -- (axis cs:0,2.5);
  \node[poporange, anchor=west] at (axis cs:2.15,9.4) {\footnotesize $A\,(r_0\,;\,g_0)$};
  \node[poporange, anchor=west] at (axis cs:4.15,3.1) {\footnotesize $B\,(2r_0\,;\,g_0/4)$};
\end{axis}
\end{tikzpicture}
\end{center}
\legende{Norme du champ de gravitation en fonction de la distance à la source (loi en $1/r^2$). Lecture graphique : au point $B$, situé à distance double du point $A$, la norme du champ est divisée par quatre.}
\end{popfigure}

\begin{exemplebox}[Lecture graphique sur la courbe en $1/r^2$]
Sur le graphique ci-dessus, le point $A$ correspond à la surface d'une planète : $r_0 = \SI{2,0e6}{m}$ et $g(A) = g_0 = \SI{10}{\newton\per\kilogram}$. Le point $B$, à distance $2r_0 = \SI{4,0e6}{m}$ du centre, donne $g(B) = \SI{2,5}{\newton\per\kilogram}$. Vérification par le calcul :
\[ \frac{g(B)}{g(A)} = \frac{r_0^2}{(2r_0)^2} = \frac{1}{4} \quad\Longrightarrow\quad g(B) = \frac{10}{4} = \SI{2,5}{\newton\per\kilogram}. \]
Graphique et calcul concordent : la loi en $1/r^2$ est bien respectée.
\end{exemplebox}

\courssub{Exemple chiffré : le champ créé par la Terre à la distance de la Lune}

\begin{exoresolu}[Champ terrestre à l'orbite de la Lune]
\textbf{Énoncé.} La Terre, de masse $M_T = \num{5,97}\times 10^{24}\ \mathrm{kg}$, est assimilée à un point matériel placé en son centre. La Lune se trouve à la distance $r = \num{3,84}\times 10^{8}\ \mathrm{m}$ du centre de la Terre. On donne $G = \num{6,67}\times 10^{-11}\ \mathrm{N\,m^2\,kg^{-2}}$.
\begin{enumerate}
\item Calculer la norme du champ de gravitation terrestre $g_L$ au point où se trouve la Lune. Comparer à $g_0 = \SI{9,8}{\newton\per\kilogram}$.
\item En déduire l'accélération centripète de la Lune dans son mouvement autour de la Terre.
\end{enumerate}
\tcblower
\textbf{Solution détaillée.}
\begin{enumerate}
\item On applique l'expression du champ créé par un point matériel :
\begin{align*}
g_L &= \frac{G\,M_T}{r^2} = \frac{\num{6,67e-11}\times \num{5,97e24}}{(\num{3,84e8})^2} \\
&= \frac{\num{3,98e14}}{\num{1,475e17}} \approx \num{2,70e-3}\ \mathrm{N\,kg^{-1}}.
\end{align*}
Comparaison : $\dfrac{g_0}{g_L} = \dfrac{9,8}{\num{2,70e-3}} \approx 3630$. Le champ terrestre à la distance de la Lune est environ \textbf{3600 fois plus faible} qu'à la surface — soit $60^2$, car la Lune est à 60 rayons terrestres du centre : la loi en $1/r^2$ frappe fort, mais elle n'annule jamais le champ.
\item La Lune n'est soumise qu'à la gravitation terrestre. D'après la deuxième loi de Newton, $m_L\,\vec{a} = m_L\,\vec{g}_L$, donc :
\[ a_L = g_L \approx \SI{2,70e-3}{\meter\per\second\squared}. \]
C'est cette petite accélération, dirigée vers la Terre, qui courbe en permanence la trajectoire de la Lune et la maintient en orbite — exactement comme la pomme de Newton, mais 380 000 km plus loin.
\end{enumerate}
\end{exoresolu}

\begin{savaistu}[La pomme et la Lune : la même loi]
C'est précisément ce calcul que Newton aurait effectué vers 1666 : sachant que la Lune est à environ 60 rayons terrestres, son accélération de chute vers la Terre doit valoir $g_0/60^2 \approx \SI{2,7e-3}{\meter\per\second\squared}$ — ce que confirme l'astronomie. La même loi régit la chute d'une pomme à Mvolyé et le mouvement de la Lune : c'est toute la révolution de la gravitation \emph{universelle}. Aujourd'hui, les satellites de télécommunication qui permettent les réseaux MTN et Orange au Cameroun sont placés sur des orbites calculées avec cette même formule du champ en $1/r^2$.
\end{savaistu}

\courssub{Superposition des champs (admis)}

Que se passe-t-il lorsque plusieurs masses créent chacune un champ en un même point $M$ ? Par exemple, un satellite placé entre la Terre et la Lune « ressent » les deux astres. On admet le principe suivant, conséquence directe de l'additivité des forces.

\begin{propriete}[Principe de superposition (admis)]
Le champ de gravitation créé en un point $M$ par plusieurs masses $m_1, m_2, \ldots, m_n$ est la \textbf{somme vectorielle} des champs créés séparément par chaque masse :
\[ \vec{g}(M) = \vec{g}_1(M) + \vec{g}_2(M) + \cdots + \vec{g}_n(M). \]
\end{propriete}

Ce principe sera précieux pour la suite de la leçon : c'est lui qui permettra de calculer le champ créé par une sphère homogène, en la décomposant (mentalement) en une infinité de petites masses ponctuelles dont on additionne les contributions.

\begin{exercice}[Champ créé par deux masses]
Deux points matériels $A$ et $B$, de masses $m_A = \SI{4,0e15}{kg}$ et $m_B = \SI{1,0e15}{kg}$, sont distants de $AB = \SI{10}{km}$. Déterminer le champ de gravitation (direction, sens, norme) au point $M$, milieu du segment $[AB]$. On donne $G = \num{6,67}\times 10^{-11}\ \mathrm{N\,m^2\,kg^{-2}}$.
\end{exercice}

\begin{corrige}[Exercice : champ créé par deux masses]
$M$ est à $r = \SI{5,0}{km} = \SI{5,0e3}{m}$ de chaque masse. Les deux champs $\vec{g}_A$ et $\vec{g}_B$ sont colinéaires, portés par $(AB)$, et de sens opposés (l'un vers $A$, l'autre vers $B$). Normes :
\[ g_A = \frac{\num{6,67e-11}\times\num{4,0e15}}{(\num{5,0e3})^2} = \SI{1,07e-2}{\newton\per\kilogram}, \qquad g_B = \frac{\num{6,67e-11}\times\num{1,0e15}}{(\num{5,0e3})^2} = \SI{2,67e-3}{\newton\per\kilogram}. \]
Par superposition, le champ total est dirigé vers $A$ (le champ le plus intense l'emporte) et sa norme vaut :
\[ g(M) = g_A - g_B = \num{1,07e-2} - \num{2,67e-3} \approx \SI{8,0e-3}{\newton\per\kilogram}. \]
Remarque : il existe sur le segment $[AB]$ un point où le champ total est nul — plus proche de la masse la plus faible. Le retrouver est un excellent entraînement !
\end{corrige}

\begin{aretenir}[L'essentiel de la section]
\begin{itemize}
\item Le champ de gravitation créé en $M$ par un point matériel $S$ de masse $m_S$ est $\vec{g}(M) = -G\dfrac{m_S}{r^2}\vec{u}$ : direction radiale, sens centripète (vers $S$), norme $G\,m_S/r^2$.
\item Toute masse $m$ placée en $M$ subit $\vec{F} = m\,\vec{g}(M)$.
\item Le champ s'exprime en $\mathrm{N\,kg^{-1}}$, équivalent à $\mathrm{m\,s^{-2}}$ : c'est l'accélération de tout objet en chute libre en $M$, indépendante de sa masse.
\item Le champ existe sans masse test ; il ne dépend que de la source et de la position.
\item Les champs de plusieurs sources s'additionnent vectoriellement (superposition).
\end{itemize}
\end{aretenir}

Maintenant que le champ d'un point matériel est bien compris, une question naturelle se pose : la Terre n'est pas un point, c'est une sphère de 6400 km de rayon ! Peut-on quand même utiliser la formule en $1/r^2$ ? C'est l'objet de la section suivante, où nous montrerons le résultat remarquable obtenu par Newton : à l'extérieur, une sphère homogène se comporte exactement comme si toute sa masse était concentrée en son centre.

\coursec{Champ créé par une sphère homogène}

\courssub{Du point matériel à l'astre réel : le problème posé}

Dans la section précédente, nous avons construit le vecteur \cle{champ de gravitation} créé par un \cle{point matériel} de masse $m$ :
\[
\vec{\mathcal{G}}(M) = -\,\frac{G\,m}{r^{2}}\,\vec{u}_{OM},
\]
où $\vec{u}_{OM}$ est le vecteur unitaire dirigé de la source $O$ vers le point $M$. Ce résultat est magnifique, mais il souffre d'une objection évidente : la Terre n'est pas un point ! La Lune non plus, ni le Soleil. Quand tu te tiens à la surface de la Terre, une partie de la masse terrestre se trouve à quelques mètres sous tes pieds, une autre à \SI{12742}{km} de toi, de l'autre côté du globe. Quelle distance $r$ faut-il alors utiliser dans la formule ?

C'est précisément la question qui a longtemps bloqué Newton lui-même. Intuitivement, on peut découper la Terre en une multitude de petits morceaux de masse $m_i$, chacun créant un champ $\vec{\mathcal{G}}_i = -\dfrac{G\,m_i}{r_i^{2}}\vec{u}_i$, et additionner vectoriellement toutes ces contributions. Mais les distances $r_i$ sont toutes différentes : le calcul semble inextricable. La réponse, que Newton a fini par obtenir, est l'un des plus beaux résultats de toute la physique : \textbf{tout se passe comme si toute la masse était concentrée au centre}.

\courssub{Le théorème de la sphère homogène}

\begin{propriete}[Théorème de la sphère homogène (admis)]
Une sphère de centre $O$, de rayon $R$ et de masse $M$, dont la répartition de masse possède la \cle{symétrie sphérique} (sphère homogène, ou sphère dont la masse volumique ne dépend que de la distance au centre), crée en tout point $M$ \textbf{extérieur} à la sphère ($r = OM \geq R$) le même champ de gravitation qu'un point matériel de masse $M$ placé en $O$ :
\[
\vec{\mathcal{G}}(M) = -\,\frac{G\,M}{r^{2}}\,\vec{u}_{OM}, \qquad \mathcal{G}(r) = \frac{G\,M}{r^{2}}, \quad r \geq R.
\]
Ce théorème est \textbf{admis} : sa démonstration complète nécessite un calcul intégral hors programme. Elle n'est pas exigible au Baccalauréat.
\end{propriete}

\begin{attention}[Conditions d'application du théorème]
Ce théorème exige deux conditions qu'il ne faut jamais oublier :
\begin{itemize}
\item le point $M$ doit être \textbf{à l'extérieur} de la sphère ($r \geq R$). À l'intérieur d'une sphère creuse homogène, on peut montrer que le champ est même \emph{nul} en tout point ;
\item la répartition de masse doit être à \textbf{symétrie sphérique}. Pour une patate ou un astéroïde de forme quelconque, le théorème ne s'applique pas rigoureusement.
\end{itemize}
La Terre, dont la densité augmente vers le centre (croûte, manteau, noyau), n'est pas homogène, mais sa répartition de masse est \emph{quasi sphérique} : le théorème s'applique donc avec une excellente approximation.
\end{attention}

\courssub{Pourquoi ce résultat est-il plausible ? Justification qualitative}

La démonstration rigoureuse repose sur une somme intégrale, mais l'idée physique peut se comprendre sans calcul. Considérons un point $M$ extérieur à la sphère.

\begin{itemize}
\item \textbf{Direction du champ.} Pour chaque petit morceau de la sphère situé d'un côté de l'axe $(OM)$, il existe un morceau symétrique de même masse de l'autre côté. Les composantes de leurs champs perpendiculaires à l'axe $(OM)$ s'annulent deux à deux : par \cle{symétrie}, le champ total est nécessairement dirigé selon la droite $(OM)$, c'est-à-dire \textbf{radial}. Comme la gravitation est attractive, il pointe vers $O$.
\item \textbf{Dépendance en $r$ uniquement.} La sphère étant invariante par toute rotation autour de $O$, aucune direction de l'espace n'est privilégiée : la norme du champ ne peut dépendre que de la distance $r = OM$, pas de la latitude ni de la longitude du point $M$.
\item \textbf{Loi en $1/r^{2}$.} Le calcul intégral montre alors un miracle de compensation : les morceaux proches de $M$ attirent fort mais sont peu nombreux, les morceaux lointains attirent faiblement mais sont nombreux ; la somme totale redonne \emph{exactement} $\dfrac{GM}{r^{2}}$, comme si toute la masse était en $O$.
\end{itemize}

\begin{savaistu}[Newton et la naissance du calcul intégral]
Selon la tradition, Newton aurait élaboré sa théorie de la gravitation dès 1666, lors de la grande peste qui l'avait renvoyé de Cambridge dans sa campagne natale. Pourtant, il attendit près de vingt ans avant de publier les \emph{Principia} (1687). L'une des raisons de ce retard serait précisément qu'il ne savait pas encore démontrer qu'une sphère agit comme un point matériel : sans ce théorème, appliquer sa loi à la Terre et à la Lune n'était pas rigoureux. Pour y parvenir, il dut inventer un outil mathématique nouveau, le \textbf{calcul infinitésimal} (développé en parallèle par Leibniz), ancêtre de l'intégration que tu rencontreras à l'université. Une leçon d'humilité scientifique : une loi n'est vraiment établie que lorsqu'on sait la justifier.
\end{savaistu}

\courssub{Conséquence fondamentale : la distance se mesure entre les centres}

\begin{aretenir}[La bonne distance dans la loi de Newton]
Grâce au théorème de la sphère homogène, la loi de l'attraction universelle s'applique à deux astres sphériques (planètes, étoiles, satellites naturels) \textbf{à condition de prendre pour $r$ la distance entre leurs centres} :
\[
F = \frac{G\,M\,m}{r^{2}}, \qquad r = \text{distance centre à centre}.
\]
Pour un objet situé à l'altitude $h$ au-dessus d'une planète de rayon $R$, on a donc $r = R + h$.
\end{aretenir}

\begin{attention}[Erreur fréquente : confondre altitude et distance au centre]
L'erreur la plus classique au Baccalauréat consiste à écrire $\mathcal{G} = \dfrac{GM}{h^{2}}$ en utilisant l'altitude $h$ à la place de la distance $r = R + h$ au centre de l'astre. C'est faux ! La loi de Newton fait intervenir la distance au \emph{centre} de la sphère. Par exemple, pour un satellite à \SI{400}{km} d'altitude, $r = R_T + h \approx \SI{6770}{km}$, et non \SI{400}{km} : l'erreur sur le champ serait d'un facteur 286 ! Réflexe à acquérir : \textbf{dessiner le centre $O$ de l'astre sur ton schéma et mesurer $r$ depuis $O$}.
\end{attention}

\begin{popfigure}
\begin{center}
\begin{tikzpicture}[scale=1.0]
% Sphère homogène
\shade[ball color=popblue!60] (0,0) circle (1.2);
\draw[popdark,thick] (0,0) circle (1.2);
\fill[popdark] (0,0) circle (0.05) node[below left=1pt,popdark]{$O$};
\draw[popdark,dashed] (0,0) -- (0.85,0.85) node[midway,above left,popdark]{$R$};
\node[popdark] at (-0.5,-0.55) {$M$};
% Points extérieurs et vecteurs champ
\coordinate (A) at (3.4,0);
\coordinate (B) at (0,2.6);
\coordinate (C) at (-2.2,2.2);
\coordinate (D) at (2.0,-2.6);
\draw[popmuted,dashed] (0,0) -- (A) node[midway,below,popmuted]{$r_A$};
\draw[popmuted,dashed] (0,0) -- (B) node[midway,right,popmuted]{$r_B$};
\draw[popmuted,dashed] (0,0) -- (C);
\draw[popmuted,dashed] (0,0) -- (D);
\foreach \P/\len in {A/1.1,B/1.5,C/1.5,D/1.2}{
  \fill[poporange] (\P) circle (0.07);
  \draw[-{Stealth[length=3mm]},poporange,very thick] (\P) -- ($(\P)!\len cm!(0,0)$);
}
\node[poporange,right] at (3.45,0.15){$\vec{\mathcal{G}}_A$};
\node[poporange,right] at (0.1,2.65){$\vec{\mathcal{G}}_B$};
\node[poporange,left] at (-2.25,2.3){$\vec{\mathcal{G}}_C$};
\node[poporange,right] at (2.05,-2.75){$\vec{\mathcal{G}}_D$};
% Point matériel équivalent (encart)
\begin{scope}[xshift=8.2cm]
\fill[popdark] (0,0) circle (0.09);
\node[popdark,below] at (0,-0.1){$O$ : masse $M$};
\coordinate (A2) at (3.4,0);
\coordinate (B2) at (0,2.6);
\coordinate (C2) at (-2.2,2.2);
\coordinate (D2) at (2.0,-2.6);
\foreach \P/\len in {A2/1.1,B2/1.5,C2/1.5,D2/1.2}{
  \fill[poporange] (\P) circle (0.07);
  \draw[-{Stealth[length=3mm]},poporange,very thick] (\P) -- ($(\P)!\len cm!(0,0)$);
}
\draw[popmuted,dashed] (0,0) -- (A2);
\draw[popmuted,dashed] (0,0) -- (B2);
\draw[popmuted,dashed] (0,0) -- (C2);
\draw[popmuted,dashed] (0,0) -- (D2);
\end{scope}
\draw[-{Stealth[length=3mm]},popdark,thick] (4.6,0) -- (6.4,0) node[midway,above,popdark]{\'equivalent};
\end{tikzpicture}
\end{center}
\legende{À gauche : champ créé par une sphère homogène de centre $O$ en plusieurs points extérieurs ; les vecteurs sont radiaux, dirigés vers $O$, et leur norme décroît avec $r$. À droite : le point matériel équivalent de masse $M$ placé en $O$ crée exactement le même champ.}
\end{popfigure}

\courssub{Lignes de champ et carte de champ}

Pour visualiser d'un coup d'œil la structure d'un champ de gravitation dans tout l'espace, on utilise une représentation géométrique très parlante : les \cle{lignes de champ}.

\begin{definition}[Ligne de champ]
Une \cle{ligne de champ} de gravitation est une courbe \textbf{tangente en chacun de ses points au vecteur champ} $\vec{\mathcal{G}}$ en ce point, et \textbf{orientée dans le sens du champ} (c'est-à-dire vers la masse source, puisque la gravitation est attractive).
\end{definition}

Pour la sphère homogène, le champ est radial et centripète en tout point extérieur : les lignes de champ sont donc des \textbf{demi-droites radiales convergentes vers le centre $O$}. On dit que le champ possède la \cle{symétrie sphérique}.

\begin{popfigure}
\begin{center}
\begin{tikzpicture}[scale=0.9]
\shade[ball color=popblue!60] (0,0) circle (1.0);
\draw[popdark,thick] (0,0) circle (1.0);
\fill[popdark] (0,0) circle (0.05) node[below left=1pt,popdark]{$O$};
\node[popdark] at (-0.45,-0.5){$M$};
% Lignes de champ radiales orientées vers O
\foreach \a in {0,30,60,90,120,150,180,210,240,270,300,330}{
  \draw[-{Stealth[length=2.6mm]},popgreen,thick] (\a:4.0) -- (\a:1.35);
}
% Cercles de même norme
\draw[popmuted,dashed] (0,0) circle (2.0);
\draw[popmuted,dashed] (0,0) circle (3.2);
\node[popmuted] at (1.45,1.55){$r_1$};
\node[popmuted] at (2.3,2.45){$r_2$};
\node[popmuted,align=left] at (5.6,1.2){$\mathcal{G}(r_1) = \dfrac{GM}{r_1^{2}}$\\[2pt] $\mathcal{G}(r_2) = \dfrac{GM}{r_2^{2}} < \mathcal{G}(r_1)$};
\end{tikzpicture}
\end{center}
\legende{Carte des lignes de champ d'une sphère homogène : demi-droites radiales orientées vers le centre. Sur une sphère de rayon $r$ donné (cercles en pointillés), la norme du champ est constante ; elle diminue quand $r$ augmente.}
\end{popfigure}

\begin{propriete}[Densité des lignes et intensité du champ]
Plus les lignes de champ sont \textbf{resserrées} (proches les unes des autres) en une région, plus le champ y est \textbf{intense}. Pour la sphère homogène, les lignes radiales se resserrent lorsqu'on se rapproche de la surface : c'est la traduction géométrique de la loi en $\dfrac{1}{r^{2}}$. Inversement, loin de l'astre, les lignes s'écartent et le champ s'affaiblit.
\end{propriete}

\begin{attention}[Ne pas confondre ligne de champ et trajectoire]
Une ligne de champ n'est \textbf{pas} la trajectoire d'un corps en mouvement ! Elle indique seulement la direction de la \emph{force} (donc de l'accélération) subie par un corps placé en ce point. Un satellite lancé avec une vitesse initiale horizontale ne suit pas une ligne radiale : il décrit une orbite, alors que le champ reste radial. Nous y reviendrons dans l'étude des satellites.
\end{attention}

\courssub{Application numérique : le champ à la surface de la Terre}

\begin{exoresolu}[Vérification de la valeur de $g_0$ à la surface de la Terre]
\textbf{Énoncé.} On assimile la Terre à une sphère à répartition sphérique de masse, de masse $M_T = \SI{5,97e24}{kg}$ et de rayon moyen $R_T = \SI{6,37e6}{m}$. On donne $G = \SI{6,67e-11}{N.m^{2}.kg^{-2}}$.
\begin{enumerate}
\item Calculer la norme $\mathcal{G}_0$ du champ de gravitation à la surface de la Terre.
\item Comparer à la valeur usuelle de l'intensité de la pesanteur $g_0 = \SI{9,8}{N.kg^{-1}}$.
\end{enumerate}
\tcblower
\textbf{Solution.}
\begin{enumerate}
\item La surface de la Terre est un point \emph{extérieur} (limite $r = R_T$) : le théorème de la sphère homogène s'applique avec $r = R_T$ :
\[
\mathcal{G}_0 = \frac{G\,M_T}{R_T^{2}} = \frac{\num{6,67e-11} \times \num{5,97e24}}{(\num{6,37e6})^{2}}.
\]
Calculons : $G\,M_T = \num{6,67e-11} \times \num{5,97e24} \approx \SI{3,98e14}{m^{3}.s^{-2}}$ et $R_T^{2} \approx \SI{4,06e13}{m^{2}}$, d'où
\[
\mathcal{G}_0 \approx \frac{\num{3,98e14}}{\num{4,06e13}} \approx \SI{9,8}{N.kg^{-1}}.
\]
Vérification dimensionnelle : $\dfrac{[G]\,[M]}{[R^{2}]} = \dfrac{\si{m^{3}.kg^{-1}.s^{-2}} \times \si{kg}}{\si{m^{2}}} = \si{m.s^{-2}} = \si{N.kg^{-1}}$. L'unité est correcte.
\item On retrouve exactement la valeur mesurée de l'intensité de la pesanteur : $g_0 \approx \SI{9,8}{N.kg^{-1}}$. C'est une confirmation éclatante du théorème : toute la masse de la Terre, répartie sur \SI{12700}{km} de diamètre, agit sur toi \emph{comme si} elle était concentrée en un point situé à \SI{6370}{km} sous tes pieds. C'est ce résultat qui permettra, dans la section suivante, d'écrire le champ terrestre à l'altitude $h$ sous la forme $g(h) = g_0\,\dfrac{R_T^{2}}{(R_T+h)^{2}}$.
\end{enumerate}
\end{exoresolu}

\begin{exemplebox}[Champ créé par le Soleil à la distance de la Terre]
Le Soleil ($M_S = \SI{2,0e30}{kg}$) est une étoile quasi sphérique. À la distance $r = \SI{1,5e11}{m}$ (1 unité astronomique), le champ qu'il crée vaut
\[
\mathcal{G}_S = \frac{G\,M_S}{r^{2}} = \frac{\num{6,67e-11} \times \num{2,0e30}}{(\num{1,5e11})^{2}} \approx \SI{5,9e-3}{N.kg^{-1}}.
\]
C'est ce champ, environ 1700 fois plus faible que $g_0$, qui maintient la Terre sur son orbite : la gravitation est faible, mais elle règne sans partage à l'échelle du Système solaire.
\end{exemplebox}

\begin{methode}[Exploiter le théorème de la sphère homogène]
\begin{enumerate}
\item Vérifier que l'astre source est à symétrie sphérique (planète, étoile : toujours supposé vrai dans les exercices) et que le point étudié est \emph{extérieur} ($r \geq R$).
\item Repérer le \textbf{centre $O$} de l'astre sur un schéma et calculer $r = R + h$ si une altitude est donnée.
\item Remplacer mentalement l'astre par un point matériel de masse $M$ en $O$.
\item Écrire $\mathcal{G}(r) = \dfrac{GM}{r^{2}}$, vecteur radial dirigé vers $O$, puis effectuer l'application numérique en unités SI (distances en mètres !).
\end{enumerate}
\end{methode}

\begin{exercice}[Champ lunaire à la surface de la Lune]
La Lune a une masse $M_L = \SI{7,35e22}{kg}$ et un rayon $R_L = \SI{1,74e6}{m}$.
\begin{enumerate}
\item Calculer la norme du champ de gravitation à la surface de la Lune.
\item Un astronaute de masse $m = \SI{80}{kg}$ (avec sa combinaison) y pèse combien ? Comparer à son poids terrestre ($g_0 = \SI{9,8}{N.kg^{-1}}$).
\end{enumerate}
\end{exercice}

\begin{corrige}[Exercice : champ lunaire]
\begin{enumerate}
\item $\mathcal{G}_L = \dfrac{G\,M_L}{R_L^{2}} = \dfrac{\num{6,67e-11} \times \num{7,35e22}}{(\num{1,74e6})^{2}} \approx \dfrac{\num{4,90e12}}{\num{3,03e12}} \approx \SI{1,6}{N.kg^{-1}}$.
\item Poids lunaire : $P_L = m\,\mathcal{G}_L \approx 80 \times 1,6 \approx \SI{130}{N}$. Poids terrestre : $P_T = m\,g_0 \approx 80 \times 9,8 \approx \SI{784}{N}$. L'astronaute est environ 6 fois plus léger sur la Lune (rapport $\dfrac{g_0}{\mathcal{G}_L} \approx 6$) : c'est ce qui permet les bonds spectaculaires des missions Apollo, la masse (et donc l'inertie) restant inchangée.
\end{enumerate}
\end{corrige}

\begin{aretenir}[L'essentiel de la section]
\begin{itemize}
\item Une sphère à symétrie sphérique de masse crée, \textbf{à l'extérieur}, le même champ qu'un point matériel de même masse placé en son centre : $\mathcal{G}(r) = \dfrac{GM}{r^{2}}$ pour $r \geq R$.
\item La distance $r$ se mesure \textbf{toujours depuis le centre} de l'astre : $r = R + h$.
\item Les lignes de champ sont radiales, orientées vers le centre ; leur resserrement près de la surface traduit la croissance du champ en $\dfrac{1}{r^{2}}$.
\item À la surface de la Terre : $g_0 = \dfrac{G\,M_T}{R_T^{2}} \approx \SI{9,8}{N.kg^{-1}}$.
\end{itemize}
\end{aretenir}

\coursec{Champ de gravitation terrestre et variation avec l'altitude}

Nous venons d'établir un résultat remarquable : une sphère homogène (pleine ou creuse) crée, en tout point extérieur, le même champ de gravitation qu'un point matériel de même masse placé en son centre. La Terre, à très bonne approximation, est une telle sphère : nous pouvons donc appliquer ce théorème pour décrire le champ de gravitation terrestre, à la surface du sol comme en altitude. C'est l'objet de cette section, capitale pour la suite du cours (satellites, poids, chute libre).

\courssub{Application du théorème de la sphère homogène à la Terre}

\begin{definition}[Modèle de la Terre]
On modélise la Terre par une sphère homogène de centre $O$, de masse $M_T$ et de rayon moyen
\[ R_T = \SI{6380}{km} = \num{6,38e6}\ \si{m}. \]
D'après le théorème de la sphère homogène, le champ de gravitation terrestre en un point $M$ extérieur à la Terre, situé à la distance $r = OM$ du centre, est :
\[ \vec{\mathcal{G}}(M) = -\,G\,\frac{M_T}{r^2}\,\vec{u} \]
où $\vec{u}$ est le vecteur unitaire dirigé du centre $O$ vers le point $M$. Le champ est donc \cle{centripète} : il pointe toujours vers le centre de la Terre.
\end{definition}

\begin{popfigure}
\begin{center}
\begin{tikzpicture}[scale=1.0]
  % Terre
  \fill[popblueL] (0,0) circle (1.6);
  \draw[popblue, thick] (0,0) circle (1.6);
  \shade[ball color=popblue!40, opacity=0.6] (0,0) circle (1.6);
  \node[popdark] at (0,-0.6) {\small Terre};
  % Centre
  \fill[popdark] (0,0) circle (0.05);
  \node[below left, popdark] at (0,0) {$O$};
  % Point M
  \coordinate (M) at (3.6,0);
  \fill[poporange] (M) circle (0.07);
  \node[above, poporange] at (M) {$M$};
  % Rayon R_T
  \draw[popdark, thick, ->] (0,0) -- (1.13,1.13) node[midway, above left] {$R_T$};
  % Altitude h
  \draw[popgreen, thick, <->] (1.6,0.35) -- (3.6,0.35) node[midway, above] {$h$};
  % Distance r
  \draw[popmuted, dashed] (0,-0.35) -- (3.6,-0.35);
  \node[popmuted, below] at (1.8,-0.35) {$r = R_T + h$};
  % Vecteur champ
  \draw[popink, very thick, ->] (M) -- (2.1,0) node[midway, above] {$\vec{\mathcal{G}}(M)$};
  % Vecteur unitaire
  \draw[popmuted, ->] (0,0) -- (0.9,0) node[midway, below] {$\vec{u}$};
\end{tikzpicture}
\end{center}
\legende{Coupe de la Terre : le point $M$ à l'altitude $h$ est à la distance $r = R_T + h$ du centre $O$. Le vecteur champ de gravitation $\vec{\mathcal{G}}(M)$ est dirigé vers le centre.}
\end{popfigure}

\courssub{Champ à la surface de la Terre : la valeur $g_0$}

À la surface du sol ($h = 0$, donc $r = R_T$), le champ de gravitation vaut :
\[ g_0 = \frac{G\,M_T}{R_T^2}. \]

\begin{exemplebox}[Calcul de $g_0$ à la surface terrestre]
Données : $G = \num{6,67e-11}\ \si{N.m^2.kg^{-2}}$, $M_T = \num{5,98e24}\ \si{kg}$, $R_T = \num{6,38e6}\ \si{m}$.
\begin{align*}
g_0 &= \frac{\num{6,67e-11} \times \num{5,98e24}}{(\num{6,38e6})^2} \\
&= \frac{\num{3,99e14}}{\num{4,07e13}} \approx \SI{9,80}{N/kg}.
\end{align*}
On retrouve bien la valeur connue du champ de gravitation terrestre au niveau de la mer : $g_0 \approx \SI{9,80}{N/kg}$ (ou $\si{m/s^2}$, les deux unités étant équivalentes).
\end{exemplebox}

\begin{propriete}[Poids et force de gravitation — distinction]
Au voisinage de la surface terrestre, le \cle{poids} $\vec{P}$ d'un objet de masse $m$ est la force de gravitation $\vec{F}_g$ que la Terre exerce sur lui :
\[ \vec{P} = m\,\vec{g} \quad \text{avec} \quad \vec{g} = \vec{\mathcal{G}}_{\text{Terre}}. \]
Les termes « champ de pesanteur » et « champ de gravitation terrestre » désignent le même champ vectoriel $\vec{g}$ dès lors qu'on néglige la rotation de la Terre (effet centrifuge, non au programme). La distinction stricte entre \emph{poids} (force) et \emph{masse} (scalaire) reste essentielle : $P = m\,g$ avec $g = \|\vec{g}\|$.
\end{propriete}

\courssub{Champ à l'altitude $h$ : la relation fondamentale}

Un point $M$ situé à l'altitude $h$ au-dessus du sol est à la distance $r = R_T + h$ du centre de la Terre. Le champ de gravitation y vaut donc :
\[ g(h) = \frac{G\,M_T}{(R_T + h)^2}. \]

Cette expression fait intervenir $G$ et $M_T$, rarement données dans les exercices du Baccalauréat. On préfère l'exprimer en fonction de $g_0$ et $R_T$, qui sont les données usuelles.

\begin{propriete}[Variation du champ de gravitation avec l'altitude — démonstration exigible]
Le champ de gravitation terrestre à l'altitude $h$ s'exprime en fonction de sa valeur au sol par :
\[ g(h) = g_0\,\frac{R_T^2}{(R_T + h)^2}. \]
\textbf{Démonstration guidée.} Écrivons les deux expressions du champ :
\[ g_0 = \frac{G\,M_T}{R_T^2} \qquad \text{et} \qquad g(h) = \frac{G\,M_T}{(R_T + h)^2}. \]
Formons le rapport membre à membre afin d'éliminer le produit $G\,M_T$ :
\[ \frac{g(h)}{g_0} = \frac{G\,M_T}{(R_T+h)^2} \times \frac{R_T^2}{G\,M_T} = \frac{R_T^2}{(R_T + h)^2}, \]
d'où le résultat annoncé. \emph{Vérification dimensionnelle} : le rapport $R_T^2/(R_T+h)^2$ est sans dimension, donc $g(h)$ a bien la même unité que $g_0$.
\end{propriete}

\begin{methode}[Calculer $g$ à une altitude donnée]
\begin{enumerate}
\item \textbf{Convertir} l'altitude $h$ en mètres (attention aux km !) et exprimer $R_T$ dans la même unité.
\item \textbf{Calculer} la distance au centre : $r = R_T + h$.
\item \textbf{Appliquer} la formule $g(h) = g_0\,\dfrac{R_T^2}{(R_T + h)^2}$.
\item \textbf{Vérifier} la cohérence : on doit toujours trouver $g(h) < g_0$, et $g(h)$ d'autant plus faible que $h$ est grand.
\end{enumerate}
\end{methode}

\courssub{Étude de la fonction $g(h)$ : décroissance en $1/r^2$}

La fonction $h \mapsto g(h)$ est une fonction \cle{décroissante} de l'altitude : plus on s'éloigne de la Terre, plus le champ diminue, en suivant une loi en $1/(R_T+h)^2$. Deux valeurs remarquables à connaître :

\begin{itemize}
\item En $h = 0$ : $g(0) = g_0\,\dfrac{R_T^2}{R_T^2} = g_0$ (cohérent !).
\item En $h = R_T$ (soit à $\SI{6380}{km}$ d'altitude, à deux rayons terrestres du centre) :
\[ g(R_T) = g_0\,\frac{R_T^2}{(2R_T)^2} = \frac{g_0}{4} \approx \SI{2,45}{N/kg}. \]
Il faut donc s'élever d'un rayon terrestre entier pour diviser le champ par quatre : la gravité décroît vite, mais pas si vite !
\end{itemize}

\begin{popfigure}
\begin{center}
\begin{tikzpicture}
\begin{axis}[
    width=0.85\linewidth, height=6.5cm,
    xlabel={Altitude $h$ (km)}, ylabel={$g(h)$ (N/kg)},
    xmin=0, xmax=36000, ymin=0, ymax=10.5,
    xtick={0,400,6380,12000,20000,36000},
    x tick label style={/pgf/number format/fixed},
    grid=major, grid style={popline!50},
    axis lines=left,
    label style={popdark},
]
\addplot[popcurve, domain=0:36000, samples=200] {9.8*(6380/(6380+x))^2};
% Points remarquables
\addplot[only marks, mark=*, mark size=2.5pt, poporange] coordinates {(0,9.8)};
\addplot[only marks, mark=*, mark size=2.5pt, popgreen] coordinates {(400,8.68)};
\addplot[only marks, mark=*, mark size=2.5pt, poppurple] coordinates {(36000,0.223)};
\node[poporange, anchor=south west, font=\small] at (axis cs:100,9.5) {Surface : $g_0$};
\node[popgreen, anchor=south west, font=\small] at (axis cs:500,8.5) {ISS (400 km) : $\approx 8{,}7$};
\node[poppurple, anchor=north east, font=\small] at (axis cs:36000,0.3) {Géostationnaire (36\,000 km) : $\approx 0{,}22$};
\end{axis}
\end{tikzpicture}
\end{center}
\legende{Courbe $g(h) = g_0\,\dfrac{R_T^2}{(R_T+h)^2}$ pour $h$ de $0$ à $\SI{36000}{km}$. La décroissance est rapide au début, puis lente : même à l'altitude de l'ISS, le champ reste proche de $g_0$.}
\end{popfigure}

\courssub{Exemples concrets : du mont Cameroun à la Station spatiale}

\begin{exoresolu}[Champ de gravitation au sommet du mont Cameroun]
\textbf{Énoncé.} Le mont Cameroun, point culminant de l'Afrique centrale, s'élève à $h = \SI{4095}{m}$ au-dessus du niveau de la mer. Calculer le champ de gravitation à son sommet et le comparer à $g_0 = \SI{9,80}{N/kg}$. On donne $R_T = \SI{6380}{km}$.
\tcblower
\textbf{Solution.} Suivons la méthode.
\begin{enumerate}
\item Conversion : $h = \SI{4095}{m} = \SI{4,095}{km}$ ; $R_T = \SI{6380}{km}$.
\item Distance au centre : $r = R_T + h = 6380 + 4{,}095 = \SI{6384,095}{km}$.
\item Application numérique :
\begin{align*}
g(h) &= 9{,}80 \times \left(\frac{6380}{6384{,}095}\right)^2 \\
&= 9{,}80 \times (0{,}999358)^2 \\
&\approx 9{,}80 \times 0{,}998716 \approx \SI{9,79}{N/kg}.
\end{align*}
\item Vérification : $9{,}79 < 9{,}80$ \checkmark.
\end{enumerate}
La variation relative est $\dfrac{g_0 - g(h)}{g_0} \approx 0{,}13\ \%$ : un randonneur au sommet du mont Cameroun ne sent évidemment aucune différence de pesanteur !
\end{exoresolu}

\begin{exemplebox}[Champ de gravitation à bord de l'ISS]
La Station spatiale internationale orbite à $h = \SI{400}{km}$ d'altitude. Alors :
\[ g(400\ \text{km}) = 9{,}80 \times \left(\frac{6380}{6380 + 400}\right)^2 = 9{,}80 \times \left(\frac{6380}{6780}\right)^2 \approx 9{,}80 \times 0{,}885 \approx \SI{8,7}{N/kg}. \]
Le champ à bord de l'ISS vaut encore environ $89\ \%$ de sa valeur au sol. Les astronautes sont donc loin d'être « hors de la gravité » !
\end{exemplebox}

\begin{attention}[La gravité n'est PAS nulle dans l'ISS !]
Erreur fréquente : croire que les astronautes flottent parce qu'il n'y aurait « plus de gravité » dans l'espace. Faux ! À $\SI{400}{km}$ d'altitude, $g \approx \SI{8,7}{N/kg}$, soit $89\ \%$ de $g_0$. L'\cle{apesanteur} ressentie à bord de l'ISS provient de la \textbf{chute libre permanente} : la station et son équipage tombent ensemble, en permanence, autour de la Terre (leur vitesse horizontale les empêche de s'écraser). Tout se passe comme dans un ascenseur dont le câble aurait lâché : les objets flottent les uns par rapport aux autres, alors que la gravité agit toujours sur chacun d'eux.
\end{attention}

\begin{savaistu}[Et les satellites géostationnaires ?]
Les satellites de télécommunications (ceux qui relaient les signaux TV et les communications de CAMTEL, MTN ou Orange) sont placés sur l'orbite géostationnaire, à $h \approx \SI{36000}{km}$ d'altitude, où ils tournent à la même vitesse angulaire que la Terre et paraissent fixes dans le ciel. À cette altitude :
\[ g \approx 9{,}80 \times \left(\frac{6380}{6380 + 36000}\right)^2 \approx \SI{0,22}{N/kg}, \]
soit seulement $2{,}3\ \%$ de $g_0$. C'est pourtant ce « petit » champ qui suffit à courber la trajectoire du satellite et à le maintenir sur son orbite circulaire — nous le démontrerons dans la leçon sur les satellites.
\end{savaistu}

\courssub{Approximation au voisinage du sol}

Pour les altitudes très petites devant le rayon terrestre ($h \ll R_T$), on peut simplifier la formule exacte.

\begin{propriete}[Développement limité — complément utile, non exigible]
Si $h \ll R_T$, un développement limité au premier ordre donne :
\[ g(h) = g_0\left(1 + \frac{h}{R_T}\right)^{-2} \approx g_0\left(1 - \frac{2h}{R_T}\right). \]
La variation relative du champ est donc $\dfrac{\Delta g}{g_0} \approx -\dfrac{2h}{R_T}$.
\end{propriete}

\begin{exemplebox}[Ordre de grandeur : un avion de ligne à 10 km]
Pour $h = \SI{10}{km}$ :
\[ \frac{\Delta g}{g_0} \approx -\frac{2 \times 10}{6380} \approx -\num{3,1e-3}, \]
soit une diminution d'environ $0{,}3\ \%$ seulement. Même à bord d'un avion de ligne, le champ de gravitation est quasiment identique à celui du sol.
\end{exemplebox}

\begin{aretenir}[Conséquence pratique : champ uniforme à l'échelle locale]
À l'échelle d'une région limitée (quelques kilomètres en altitude, quelques dizaines de kilomètres en largeur — un stade, une ville comme Yaoundé ou Douala, un chantier) :
\begin{itemize}
\item la \textbf{norme} de $\vec{\mathcal{G}}$ varie de moins de $0{,}1\ \%$ : on la considère constante, $g \approx g_0$ ;
\item la \textbf{direction} des vecteurs champ (tous dirigés vers le centre de la Terre) est pratiquement la même partout : les lignes de champ sont quasi parallèles.
\end{itemize}
On dit que le champ de gravitation est \cle{uniforme} au voisinage du sol : $\vec{g}$ est un vecteur constant en norme, direction et sens. C'est cette approximation qui justifiera, dans la suite du cours, le mouvement parabolique des projectiles (ballon de football, obus, mototaxi franchissant une dune).
\end{aretenir}

\begin{exercice}[Pour s'entraîner]
On donne $g_0 = \SI{9,80}{N/kg}$ et $R_T = \SI{6380}{km}$.
\begin{enumerate}
\item Calculer le champ de gravitation à l'altitude $h = \SI{1000}{km}$.
\item À quelle altitude $h$ le champ de gravitation vaut-il $g_0/2$ ? Exprimer le résultat en fonction de $R_T$, puis le calculer.
\end{enumerate}
\end{exercice}

\begin{corrige}[Exercice : champ à 1000 km et altitude où $g = g_0/2$]
\begin{enumerate}
\item $g(1000\ \text{km}) = 9{,}80 \times \left(\dfrac{6380}{7380}\right)^2 \approx 9{,}80 \times 0{,}747 \approx \SI{7,3}{N/kg}$.
\item On résout $g_0\,\dfrac{R_T^2}{(R_T+h)^2} = \dfrac{g_0}{2}$, soit $(R_T+h)^2 = 2R_T^2$, donc $R_T + h = R_T\sqrt{2}$ et
\[ h = R_T(\sqrt{2} - 1) \approx 0{,}414\,R_T \approx \SI{2640}{km}. \]
Vérification : $2640 < 6380$, et $g_0/2$ est bien compris entre $g_0/4$ (atteint à $h = R_T$) et $g_0$ \checkmark.
\end{enumerate}
\end{corrige}

\coursec{Champ uniforme au voisinage du sol ; poids et force de gravitation}

Dans la section précédente, nous avons établi que le champ de gravitation terrestre diminue avec l'altitude selon la loi $g(h) = g_0\,\dfrac{R_T^2}{(R_T+h)^2}$. Mais ce résultat soulève une question pratique : lorsque tu étudies la chute d'une bille dans la cour du lycée ou le tir d'un ballon de football, faut-il vraiment tenir compte de la variation de $g$ entre le sol et le sommet de la trajectoire ? La réponse est non, et cette section explique pourquoi. Elle introduit la notion fondamentale de \cle{champ uniforme}, puis définit rigoureusement le \cle{poids} d'un corps et le distingue de la force de gravitation.

\courssub{Champ uniforme au voisinage du sol}

\textbf{Observation et intuition.}\par
Imagine que tu observes la Terre depuis l'espace : les lignes du champ de gravitation sont radiales, elles convergent toutes vers le centre de la planète, comme les rayons d'une roue. Maintenant, zoome sur un quartier de Yaoundé, sur une surface de quelques kilomètres carrés : à cette échelle, la courbure de la Terre devient imperceptible, les lignes de champ apparaissent parallèles entre elles, et l'intensité du champ ne varie pratiquement pas. C'est exactement ce qui se passe pour toutes les expériences réalisées à l'échelle humaine.

\begin{definition}[Champ uniforme]
Un champ de gravitation est dit \cle{uniforme} dans une région de l'espace si le vecteur champ $\vec{\mathcal{G}}$ y est \textbf{constant en direction, en sens et en norme} en tout point de cette région :
\[
\vec{\mathcal{G}}(M) = \vec{\mathcal{G}}(N) \quad \text{pour tous points } M \text{ et } N \text{ de la région.}
\]
Les lignes de champ d'un champ uniforme sont des \textbf{droites parallèles équidistantes}, orientées dans le sens du champ.
\end{definition}

\textbf{Justification quantitative.}\par
Deux arguments montrent que le champ terrestre est quasi uniforme au voisinage du sol.

\emph{Argument 1 : la direction.} Deux lignes de champ radiales issues du centre de la Terre et passant par deux points du sol distants de $d = \SI{1}{km}$ font entre elles un angle $\theta \approx \dfrac{d}{R_T} = \dfrac{10^3}{6{,}4\times 10^6} \approx \num{1,6e-4}~\text{rad}$, soit environ $0{,}009$ degré : elles sont indiscernables de droites parallèles.

\emph{Argument 2 : la norme.} Entre le sol ($h=0$) et le sommet d'un immeuble de $h = \SI{100}{m}$, la variation relative du champ est :
\[
\frac{\Delta g}{g_0} \approx \frac{2h}{R_T} = \frac{2 \times 10^2}{6{,}4\times 10^6} \approx \num{3,1e-5},
\]
soit environ $0{,}003\,\%$ : totalement négligeable devant la précision des mesures courantes.

\begin{aretenir}[Condition d'uniformité]
Le champ de gravitation terrestre peut être considéré comme \textbf{uniforme} dans toute région dont les dimensions (horizontales et verticales) sont \textbf{très petites devant le rayon terrestre} $R_T \approx \SI{6400}{km}$. C'est le cas de toutes les expériences de laboratoire, des chutes libres et des mouvements de projectiles étudiés cette année.
\end{aretenir}

\begin{popfigure}
\begin{tikzpicture}[scale=1]
% --- Gauche : échelle planétaire ---
\begin{scope}
\fill[popblue!25] (0,0) circle (1.5);
\draw[popblue, thick] (0,0) circle (1.5);
\node[popdark] at (0,0) {\textbf{Terre}};
\foreach \a in {30,60,90,120,150,210,240,270,300,330}{
  \draw[-{Stealth[length=2.5mm]}, poporange, thick] (\a:3.4) -- (\a:1.65);
}
\node[popmuted, align=center] at (0,-4.1) {\footnotesize Échelle planétaire :\\ \footnotesize lignes radiales};
\end{scope}
% --- Droite : zoom ---
\begin{scope}[xshift=8.5cm]
\fill[popgreenL] (-2.6,0) rectangle (2.6,0.35);
\draw[popgreen!70!black, thick] (-2.6,0.35) -- (2.6,0.35);
\node[popdark, anchor=west] at (-2.5,0.15) {\footnotesize sol};
\foreach \x in {-2,-1,0,1,2}{
  \draw[-{Stealth[length=2.5mm]}, poporange, thick] (\x,2.9) -- (\x,0.5);
  \fill[popink] (\x,2.9) circle (0.03);
}
\draw[-{Stealth[length=3mm]}, poppurple, very thick] (2.35,2.3) -- (2.35,1.1) node[midway, right]{$\vec{g}$};
\node[popmuted, align=center] at (0,-0.7) {\footnotesize Zoom (quelques km) : lignes verticales\\ \footnotesize équidistantes, champ uniforme};
\end{scope}
\end{tikzpicture}
\legende{À l'échelle de la planète, les lignes de champ convergent vers le centre de la Terre (gauche). À l'échelle d'une ville, elles sont quasi parallèles et équidistantes : le champ est uniforme (droite).}
\end{popfigure}

\textbf{Conséquence pratique.}\par
Dans tout ce qui suit (chute libre, mouvement des projectiles, pendule), le champ de pesanteur local sera assimilé à un champ uniforme : le vecteur $\vec{g}$ est vertical, dirigé vers le bas, de norme constante $g \approx \SI{9,8}{N.kg^{-1}}$. Cette approximation simplifie considérablement l'étude des mouvements : l'accélération d'un corps en chute libre est alors constante, ce qui conduit aux mouvements uniformément variés que tu étudieras dans les leçons de mécanique.

\courssub{Le poids d'un corps}

\textbf{Situation concrète.}\par
Sur le marché de Mokolo, une vendeuse suspend un sac d'oignons à son peson : l'aiguille indique une force. Quelle est cette force ? C'est l'attraction que la Terre exerce sur le sac : c'est son \cle{poids}.

\begin{definition}[Poids]
Le \cle{poids} $\vec{P}$ d'un corps de masse $m$ est la force d'attraction gravitationnelle exercée par la Terre sur ce corps. Il s'écrit :
\[
\boxed{\;\vec{P} = m\,\vec{g}\;}
\]
\begin{itemize}
\item \textbf{point d'application} : le centre de gravité $G$ du corps ;
\item \textbf{direction} : la verticale du lieu ;
\item \textbf{sens} : vers le bas (vers le centre de la Terre) ;
\item \textbf{norme} : $P = m\,g$, avec $P$ en newtons (N), $m$ en kilogrammes (kg) et $g$ en newtons par kilogramme ($\mathrm{N.kg^{-1}}$).
\end{itemize}
\end{definition}

\textbf{Cohérence avec la loi de Newton.}\par
Vérifions que cette définition est bien compatible avec la loi de l'attraction universelle. Pour un corps de masse $m$ situé à l'altitude $h$, la force gravitationnelle exercée par la Terre vaut :
\[
F = G\,\frac{M_T\,m}{(R_T+h)^2} = m \times \underbrace{G\,\frac{M_T}{(R_T+h)^2}}_{=\,g(h)} = m\,g(h).
\]
On retrouve bien $P = m\,g(h)$ : le poids n'est autre que la force de gravitation exprimée à l'aide du champ local. L'analyse dimensionnelle confirme la cohérence : $[m\,g] = \mathrm{kg}\times\mathrm{N.kg^{-1}} = \mathrm{N}$, une force.

\begin{popfigure}
\begin{tikzpicture}[scale=1]
% Sol
\fill[popgreenL] (-3,0) rectangle (3,0.3);
\draw[popgreen!70!black, thick] (-3,0.3) -- (3,0.3);
\foreach \x in {-2.8,-2.4,...,2.8}{\draw[popgreen!70!black] (\x,0) -- (\x+0.25,0.3);}
% Corps (caisse)
\fill[popgoldL] (-0.6,0.3) rectangle (0.6,1.5);
\draw[popdark, thick] (-0.6,0.3) rectangle (0.6,1.5);
\node[popdark] at (0,0.9) {\footnotesize masse $m$};
% Centre de gravité
\fill[popink] (0,0.9) circle (0.06);
\node[popink, right] at (0.08,0.9) {\footnotesize $G$};
% Vecteur poids
\draw[-{Stealth[length=3.5mm]}, poporange, very thick] (0,0.9) -- (0,-1.6) node[midway, left]{$\vec{P} = m\vec{g}$};
% Vecteur g à côté
\draw[-{Stealth[length=3.5mm]}, poppurple, very thick] (1.8,2.2) -- (1.8,0.9) node[midway, right]{$\vec{g}$};
% Échelle
\draw[popmuted, thick] (-2.6,-1.6) -- (-1.6,-1.6);
\draw[popmuted, thick] (-2.6,-1.7) -- (-2.6,-1.5);
\draw[popmuted, thick] (-1.6,-1.7) -- (-1.6,-1.5);
\node[popmuted, below] at (-2.1,-1.6) {\footnotesize échelle : 1 cm $\leftrightarrow$ 100 N};
\end{tikzpicture}
\legende{Corps de masse $m$ au voisinage du sol : le poids $\vec{P}$ est appliqué au centre de gravité $G$, vertical et dirigé vers le bas ; le champ $\vec{g}$ est vertical en tout point de la région.}
\end{popfigure}

\begin{attention}[Masse et poids : ne pas confondre !]
Erreur classique au Baccalauréat : confondre la \textbf{masse} et le \textbf{poids}.
\begin{itemize}
\item La \textbf{masse} $m$ est une grandeur \textbf{invariable} qui caractérise la quantité de matière ; elle se mesure en \textbf{kilogrammes} (kg) avec une balance.
\item Le \textbf{poids} $P$ est une \textbf{force} ; il se mesure en \textbf{newtons} (N) avec un dynamomètre et \textbf{dépend du lieu} puisque $P = m\,g$.
\end{itemize}
Ainsi, un élève de masse $m = \SI{70}{kg}$ pèse environ $P = 70 \times 9{,}8 = \SI{686}{N}$ à Douala. Au sommet du mont Cameroun ($h \approx \SI{4100}{m}$), sa masse est inchangée mais son poids est légèrement plus faible car $g$ y est plus petit. Dire « mon poids est de 70 kg » est physiquement incorrect !
\end{attention}

\begin{exoresolu}[Variation du poids entre la mer et le mont Cameroun]
Énoncé. Un sac de riz de masse $m = \SI{50}{kg}$ est pesé au niveau de la mer ($g_0 = \SI{9,80}{N.kg^{-1}}$) puis au sommet du mont Cameroun ($h = \SI{4100}{m}$). On donne $R_T = \SI{6,4e6}{m}$. Calculer le poids du sac dans les deux lieux et conclure.
\tcblower
Solution détaillée. Au niveau de la mer :
\[
P_0 = m\,g_0 = 50 \times 9{,}80 = \SI{490}{N}.
\]
Au sommet du mont Cameroun, le champ vaut :
\[
g(h) = g_0\,\frac{R_T^2}{(R_T+h)^2} = 9{,}80 \times \left(\frac{6{,}4\times 10^6}{6{,}4\times 10^6 + 4{,}1\times 10^3}\right)^2 \approx 9{,}80 \times (1 - 0{,}00128) \approx \SI{9,787}{N.kg^{-1}}.
\]
D'où le poids au sommet :
\[
P(h) = m\,g(h) \approx 50 \times 9{,}787 \approx \SI{489,4}{N}.
\]
La variation est $\Delta P = P_0 - P(h) \approx \SI{0,63}{N}$, soit l'équivalent du poids d'environ 64 grammes de riz : totalement imperceptible à la balance du marché ! Cela confirme que, pour les altitudes usuelles, le poids peut être considéré comme constant.
\end{exoresolu}

\courssub{Distinction entre poids et force de gravitation}

\textbf{En toute rigueur.}\par
Nous avons identifié le poids à la force de gravitation. Cette identification, suffisante au programme de Terminale, est en réalité une excellente approximation. Deux effets font que le poids diffère légèrement de la seule force de gravitation :

\begin{itemize}
\item \textbf{La rotation de la Terre.} La Terre tourne sur elle-même en 24 h. Un observateur terrestre est un référentiel en rotation : tout corps posé au sol subit, en plus de l'attraction gravitationnelle, un léger effet centrifuge qui l'éloigne de l'axe de rotation. Cet effet est maximal à l'équateur et nul aux pôles.
\item \textbf{L'aplatissement des pôles.} La Terre n'est pas parfaitement sphérique : elle est aplatie aux pôles (le rayon polaire est environ 21 km plus petit que le rayon équatorial). Un corps situé au pôle est donc plus proche du centre de la Terre qu'un corps à l'équateur, et y subit une attraction légèrement plus forte.
\end{itemize}

\begin{propriete}[Variation de g avec la latitude — admise]
La combinaison de ces deux effets fait que l'intensité de la pesanteur dépend légèrement de la latitude $\lambda$ du lieu :
\begin{center}
\begin{tabular}{|l|c|c|}
\hline
\rowcolor{popblueL} \textbf{Lieu} & \textbf{Latitude} & $g$ ($\mathrm{N.kg^{-1}}$) \\
\hline
Équateur & $0°$ & $9{,}78$ \\
\hline
Cameroun (Douala) & $\approx 4°$ N & $\approx 9{,}78$ \\
\hline
Paris & $\approx 49°$ N & $9{,}81$ \\
\hline
Pôles & $90°$ & $9{,}83$ \\
\hline
\end{tabular}
\end{center}
La variation totale entre l'équateur et les pôles n'excède pas $0{,}5\,\%$. \textbf{Au programme de Terminale, on identifie le poids à la force de gravitation} et l'on prend $g \approx \SI{9,8}{N.kg^{-1}}$ (parfois $\SI{10}{N.kg^{-1}}$ dans les exercices, si l'énoncé le précise).
\end{propriete}

\begin{savaistu}[La Terre est-elle ronde ? La réponse des expéditions géodésiques]
Au XVIII\textsuperscript{e} siècle, une querelle oppose les savants : Newton affirme que la rotation de la Terre l'aplatit aux pôles, tandis que les cartésiens français soutiennent qu'elle est allongée. Pour trancher, l'Académie des sciences envoie deux expéditions mesurer la longueur d'un degré de méridien : celle de \textbf{Maupertuis} en Laponie (1736, près du pôle) et celle de \textbf{La Condamine} au Pérou (1735, près de l'équateur). Les mesures de $g$ et des arcs de méridien confirment l'aplatissement des pôles : Newton avait raison. Voltaire, ami de Maupertuis, le surnomma ironiquement « l'aplatisseur de la Terre ». Aujourd'hui, les mesures de $g$ par gravimètres ultra-précis servent encore à la prospection pétrolière et minière, y compris au Cameroun !
\end{savaistu}

\courssub{Bilan des savoir-faire de la leçon}

\begin{methode}[Résoudre un problème de gravitation : la démarche complète]
\begin{enumerate}
\item \textbf{Identifier les corps} en interaction (masses $M$ et $m$) et la distance $d$ entre leurs centres.
\item \textbf{Choisir le bon modèle} : points matériels ou sphères homogènes (la sphère agit comme si toute sa masse était concentrée en son centre).
\item \textbf{Schématiser} : représenter les forces $\vec{F}_{A/B}$ et $\vec{F}_{B/A}$ (couple d'actions réciproques, même norme, sens opposés) ou le vecteur champ $\vec{\mathcal{G}}(M)$ dirigé vers le corps attracteur.
\item \textbf{Appliquer la loi de Newton} : $F = G\dfrac{M\,m}{d^2}$, avec $G = \SI{6,67e-11}{N.m^2.kg^{-2}}$.
\item \textbf{Pour le champ terrestre} : $g_0 = \dfrac{G\,M_T}{R_T^2}$ à la surface et $g(h) = g_0\dfrac{R_T^2}{(R_T+h)^2}$ à l'altitude $h$.
\item \textbf{Pour le poids} : $P = m\,g$, en choisissant la valeur de $g$ adaptée au lieu (surface, altitude $h$).
\item \textbf{Vérifier} : unités (force en N, champ en $\mathrm{N.kg^{-1}}$), ordre de grandeur, analyse dimensionnelle.
\end{enumerate}
\end{methode}

\begin{exercice}[Poids d'un satellite et champ en orbite]
Un satellite de masse $m = \SI{500}{kg}$ est en orbite à l'altitude $h = \SI{400}{km}$ au-dessus du Cameroun. On donne $g_0 = \SI{9,8}{N.kg^{-1}}$ et $R_T = \SI{6,4e3}{km}$.
\begin{enumerate}
\item Calculer le champ de gravitation $g(h)$ à cette altitude.
\item En déduire le poids du satellite sur son orbite et le comparer à son poids au sol.
\item Représenter sur un schéma la Terre, le satellite et le vecteur $\vec{\mathcal{G}}(h)$.
\end{enumerate}
\end{exercice}

\begin{corrige}[Exercice : poids d'un satellite]
\begin{enumerate}
\item $g(h) = g_0\dfrac{R_T^2}{(R_T+h)^2} = 9{,}8 \times \left(\dfrac{6400}{6400+400}\right)^2 = 9{,}8 \times \left(\dfrac{6400}{6800}\right)^2 \approx 9{,}8 \times 0{,}886 \approx \SI{8,7}{N.kg^{-1}}$.
\item Poids en orbite : $P(h) = m\,g(h) \approx 500 \times 8{,}7 = \SI{4350}{N}$. Poids au sol : $P_0 = 500 \times 9{,}8 = \SI{4900}{N}$. Le satellite pèse donc encore environ $89\,\%$ de son poids au sol : l'« impesanteur » ressentie en orbite n'est pas due à l'absence de gravité, mais au fait que le satellite est en chute libre permanente autour de la Terre.
\item Schéma : cercle représentant la Terre, point $S$ à la distance $R_T + h$ du centre, vecteur $\vec{\mathcal{G}}(h)$ partant de $S$ et dirigé radialement vers le centre de la Terre.
\end{enumerate}
\end{corrige}

\begin{aretenir}[L'essentiel de la section]
\begin{itemize}
\item Un champ est \textbf{uniforme} si $\vec{\mathcal{G}}$ y est constant en direction, sens et norme ; ses lignes de champ sont parallèles équidistantes.
\item Au voisinage du sol (dimensions $\ll R_T$), le champ terrestre est quasi uniforme : $g \approx \SI{9,8}{N.kg^{-1}}$, vertical, vers le bas.
\item Le \textbf{poids} d'un corps est $\vec{P} = m\vec{g}$, appliqué au centre de gravité, vertical vers le bas ; $P$ en newtons.
\item La \textbf{masse} (kg) est invariable ; le \textbf{poids} (N) dépend du lieu via $g$.
\item En toute rigueur, le poids diffère légèrement de la force de gravitation (rotation terrestre, aplatissement des pôles) : $g$ varie de $9{,}78$ à $\SI{9,83}{N.kg^{-1}}$ de l'équateur aux pôles. Au programme, on identifie poids et force de gravitation.
\end{itemize}
\end{aretenir}

\newpage

\coursec{Activité d'intégration}

\courssub{Objectifs de l'activité}

À l'issue de cette activité d'intégration, tu seras capable de :
\begin{itemize}
\item mobiliser la \cle{loi de l'attraction universelle} et le \cle{théorème de la sphère homogène} pour établir l'expression du champ de gravitation terrestre à la surface du globe ;
\item relier le \cle{poids} d'un objet à la force de gravitation et en déduire une méthode permettant de « peser la Terre », c'est-à-dire de déterminer sa masse $M_T$ ;
\item exploiter la relation $g(h) = g_0\,\dfrac{R_T^2}{(R_T+h)^2}$ pour calculer le champ de gravitation à différentes altitudes et interpréter sa décroissance ;
\item comparer des valeurs expérimentales ou calculées à des valeurs de référence, discuter les écarts et formuler une conclusion argumentée, comme le ferait un ingénieur ou un chercheur.
\end{itemize}

\courssub{Situation}

Lors d'une sortie pédagogique organisée par le club scientifique de ton lycée, un enseignant lance un défi aux élèves de Terminale C : « Newton a compris que la même force qui fait tomber une mangue du manguier maintient la Lune sur son orbite. Savez-vous que, sans quitter la salle de classe, vous pouvez \emph{peser la Terre} ? »

Intrigués, les élèves se souviennent qu'en TP de mécanique, ils ont mesuré l'intensité du champ de pesanteur à la surface du sol : un objet de masse $m = \SI{1,00}{kg}$ suspendu à un dynamomètre indique une force de $\SI{9,81}{N}$. Par ailleurs, un opérateur de téléphonie mobile (MTN, Orange ou CAMTEL) envisage d'installer une antenne relais au sommet du \cle{mont Cameroun} ($\SI{4095}{m}$ d'altitude) afin d'améliorer la couverture réseau entre Douala et Buéa. Le chef de projet se demande si le poids des équipements, transportés depuis Douala (altitude quasi nulle) jusqu'au sommet, varie suffisamment pour être pris en compte dans le dimensionnement des fixations. Les élèves décident de répondre à ces deux questions : déterminer la masse de la Terre, puis étudier comment le champ de gravitation varie avec l'altitude.

\textbf{Données :}
\begin{itemize}
\item intensité du champ de gravitation à la surface de la Terre : $g_0 = \SI{9,81}{N.kg^{-1}}$ ;
\item rayon moyen de la Terre (assimilée à une sphère homogène) : $R_T = \SI{6380}{km} = \SI{6,38e6}{m}$ ;
\item constante de gravitation universelle : $G = \SI{6,67e-11}{m^3.kg^{-1}.s^{-2}}$ ;
\item altitude du mont Cameroun : $h_1 = \SI{4095}{m}$ ;
\item altitude de la Station spatiale internationale (ISS) : $h_2 = \SI{400}{km}$ ;
\item masse de l'équipement de l'antenne relais : $m = \SI{200}{kg}$ ;
\item valeur officielle de la masse de la Terre : $M_T = \SI{5,97e24}{kg}$.
\end{itemize}

\courssub{Consignes}

Le travail se fait par groupes de quatre élèves. Chaque groupe rédige un compte rendu répondant aux questions suivantes.

\begin{enumerate}
\item \textbf{Peser la Terre.} En assimilant la Terre à une sphère homogène de masse $M_T$ et de rayon $R_T$, écrire l'expression de la force de gravitation exercée par la Terre sur un objet de masse $m$ posé à sa surface. En identifiant cette force au poids $P = m\,g_0$ de l'objet, établir l'expression littérale $M_T = \dfrac{g_0\,R_T^2}{G}$, puis calculer $M_T$. Comparer à la valeur officielle et commenter.
\item \textbf{Le champ de gravitation en altitude.} À l'aide de la relation $g(h) = g_0\,\dfrac{R_T^2}{(R_T+h)^2}$, calculer $g$ au sommet du mont Cameroun puis à l'altitude de l'ISS. Tracer (à la main sur papier millimétré ou sur tableur) la courbe $g(h)$ pour $h$ variant de $0$ à $\SI{40000}{km}$ (prendre par exemple $h = 0$ ; $400$ ; $2000$ ; $5000$ ; $10000$ ; $20000$ ; $40000$ km). Décrire l'allure de la courbe.
\item \textbf{L'antenne relais du mont Cameroun.} Calculer le poids de l'équipement de $\SI{200}{kg}$ à Douala ($h \approx 0$) puis au sommet du mont Cameroun. En déduire la diminution de poids $\Delta P$, en newtons puis en pourcentage. Conclure : l'ingénieur doit-il en tenir compte dans le dimensionnement des fixations ?
\item \textbf{Restitution.} Chaque groupe présente sa démarche au tableau en cinq minutes. La classe confronte les résultats à la valeur officielle $M_T = \SI{5,97e24}{kg}$ et discute les sources d'écart possibles.
\end{enumerate}

\courssub{Ressources à mobiliser}

\begin{aretenir}[Boîte à outils de l'activité]
\begin{itemize}
\item \textbf{Loi de Newton} : deux points matériels de masses $M$ et $m$, distants de $d$, s'attirent avec une force d'intensité
\[ F = G\,\dfrac{M\,m}{d^2}. \]
\item \textbf{Théorème de la sphère homogène} : à l'extérieur d'une sphère homogène (pleine ou creuse), le champ de gravitation est le même que si toute la masse était concentrée au centre : $\vec{g}(M) = -G\,\dfrac{M_T}{r^2}\,\vec{u}_r$ avec $r \geq R_T$.
\item \textbf{Poids et gravitation} : au voisinage de la Terre, le poids d'un objet s'identifie à la force de gravitation : $\vec{P} = m\,\vec{g}$, d'où
\[ g(h) = \dfrac{G\,M_T}{(R_T+h)^2} = g_0\,\dfrac{R_T^2}{(R_T+h)^2} \quad \text{avec} \quad g_0 = \dfrac{G\,M_T}{R_T^2}. \]
\item \textbf{Analyse dimensionnelle} : $[G] = \mathrm{m^3.kg^{-1}.s^{-2}}$, donc $\left[\dfrac{g_0 R_T^2}{G}\right] = \dfrac{\mathrm{m.s^{-2}}\times\mathrm{m^2}}{\mathrm{m^3.kg^{-1}.s^{-2}}} = \mathrm{kg}$ : la formule est homogène à une masse.
\end{itemize}
\end{aretenir}

\courssub{Démarche et corrigé commenté}

\begin{exoresolu}[Peser la Terre depuis la salle de classe]
\textbf{Énoncé.} 1) Établir $M_T = \dfrac{g_0 R_T^2}{G}$ et calculer $M_T$. 2) Calculer $g$ au sommet du mont Cameroun et à l'altitude de l'ISS ; décrire la courbe $g(h)$. 3) Calculer la diminution de poids d'un équipement de $\SI{200}{kg}$ transporté de Douala au sommet du mont Cameroun et conclure. 4) Discuter l'écart avec la valeur officielle $M_T = \SI{5,97e24}{kg}$.
\tcblower
\textbf{Solution détaillée.}

\textbf{1) Masse de la Terre.} D'après le théorème de la sphère homogène, la Terre attire un objet de masse $m$ situé à sa surface comme si toute sa masse $M_T$ était concentrée en son centre, à la distance $d = R_T$ :
\[ F = G\,\dfrac{M_T\,m}{R_T^2}. \]
Or le poids de l'objet à la surface est $P = m\,g_0$. En identifiant $F = P$ :
\[ G\,\dfrac{M_T\,m}{R_T^2} = m\,g_0 \quad \Longrightarrow \quad M_T = \dfrac{g_0\,R_T^2}{G}. \]
Application numérique (toutes les grandeurs en unités SI) :
\begin{align*}
M_T &= \dfrac{9{,}81 \times (6{,}38\times 10^{6})^2}{6{,}67\times 10^{-11}} \\
&= \dfrac{9{,}81 \times 4{,}07044\times 10^{13}}{6{,}67\times 10^{-11}} \\
&= \dfrac{3{,}9931\times 10^{14}}{6{,}67\times 10^{-11}} \approx \SI{5,99e24}{kg}.
\end{align*}
L'écart relatif avec la valeur officielle est $\dfrac{5{,}99-5{,}97}{5{,}97} \approx 0{,}3\,\%$ : excellent accord ! Nous venons littéralement de \emph{peser la Terre} avec un dynamomètre et une règle.

\textbf{2) Champ de gravitation en altitude.}
\begin{itemize}
\item Mont Cameroun : $h_1 = \SI{4095}{m} = \SI{4,095e3}{m}$, donc $R_T + h_1 = \SI{6,384095e6}{m}$ :
\[ g(h_1) = 9{,}81 \times \left(\dfrac{6{,}380\times 10^{6}}{6{,}384095\times 10^{6}}\right)^2 = 9{,}81 \times (0{,}999358)^2 \approx \SI{9,80}{N.kg^{-1}}. \]
La diminution n'est que d'environ $\SI{0,013}{N.kg^{-1}}$, soit $0{,}13\,\%$.
\item ISS : $h_2 = \SI{400}{km}$, donc $R_T + h_2 = \SI{6780}{km}$ :
\[ g(h_2) = 9{,}81 \times \left(\dfrac{6380}{6780}\right)^2 = 9{,}81 \times 0{,}8855 \approx \SI{8,69}{N.kg^{-1}}. \]
\end{itemize}
Quelques points de la courbe $g(h)$ (valeurs arrondies à $10^{-2}$ près) :

\begin{center}
\begin{tabularx}{\linewidth}{|l|X|X|X|X|X|X|X|}\hline
\rowcolor{popblueL}
$h$ (km) & $0$ & $400$ & $2000$ & $5000$ & $10000$ & $20000$ & $40000$ \\ \hline
$g$ (N/kg) & $9,81$ & $8,69$ & $5,70$ & $3,08$ & $1,49$ & $0,57$ & $0,19$ \\ \hline
\end{tabularx}
\end{center}

\begin{popfigure}
\begin{tikzpicture}
\begin{axis}[
    width=\linewidth,
    height=0.5\linewidth,
    xlabel={Altitude $h$ (km)},
    ylabel={Champ de gravitation $g$ (N/kg)},
    domain=0:40000,
    samples=200,
    xmin=0, xmax=40000,
    ymin=0, ymax=10,
    grid=major,
    grid style={popline},
    tick label style={/pgf/number format/use comma, /pgf/number format/1000 sep={\,}},
    x tick label style={/pgf/number format/fixed},
    y tick label style={/pgf/number format/fixed, /pgf/number format/precision=1},
    legend pos=north east,
    legend style={fill=popcream, draw=popdark},
]
\addplot[popcurve, thick] {9.81*(6380/(6380+x))^2};
\addlegendentry{$g(h) = g_0 \left(\frac{R_T}{R_T+h}\right)^2$}
% Points remarquables
\addplot[only marks, mark=*, mark size=2.5, poporange] coordinates {
    (0,9.81)
    (400,8.69)
    (2000,5.70)
    (5000,3.08)
    (10000,1.49)
    (20000,0.57)
    (40000,0.19)
};
\end{axis}
\end{tikzpicture}
\legende{Évolution du champ de gravitation terrestre $g(h)$ en fonction de l'altitude $h$. La décroissance suit une loi en $1/(R_T+h)^2$.}
\end{popfigure}

La courbe est une \cle{décroissance en $1/r^2$} (avec $r = R_T+h$) : $g$ diminue rapidement au début, puis de plus en plus lentement, sans jamais s'annuler. À $\SI{40000}{km}$ d'altitude (au-delà de l'orbite géostationnaire à $\approx \SI{36000}{km}$), le champ vaut encore environ $1,9\,\%$ de sa valeur au sol.

\textbf{3) Poids de l'équipement.}
\begin{itemize}
\item À Douala ($h \approx 0$) : $P_0 = m\,g_0 = 200 \times 9{,}81 = \SI{1962}{N}$.
\item Au sommet du mont Cameroun : $P_1 = m\,g(h_1) = 200 \times 9{,}797 \approx \SI{1959}{N}$.
\end{itemize}
Diminution : $\Delta P = P_0 - P_1 \approx \SI{2,6}{N}$, soit $\dfrac{\Delta P}{P_0} \approx 0{,}13\,\%$.

\textbf{Conclusion} : la diminution de poids est négligeable (équivalent au poids d'un objet de $\SI{260}{g}$ au sol). L'ingénieur n'a pas besoin d'en tenir compte pour les fixations de l'antenne ; en revanche, pour un satellite en orbite (ISS, géostationnaire), la variation de $g$ devient considérable et doit impérativement être prise en compte dans la mécanique orbitale.

\textbf{4) Sources d'écart sur $M_T$.} L'écart de $0{,}3\,\%$ provient principalement :
\begin{itemize}
\item de la valeur arrondie de $g_0$ (qui varie légèrement avec la latitude — effet de rotation — et la géologie locale) ;
\item de l'approximation de la Terre par une sphère parfaite (elle est aplatie aux pôles : rayon équatorial $\approx \SI{6378}{km}$, polaire $\approx \SI{6357}{km}$) ;
\item de l'hypothèse d'homogénéité (le noyau terrestre est bien plus dense que la croûte) ;
\item de la négligence de la rotation terrestre, qui distingue le poids (force mesurée) de la force de gravitation (force d'attraction pure).
\end{itemize}
\end{exoresolu}

\courssub{Bilan de l'activité}

Cette activité a permis de mettre en évidence plusieurs idées fondamentales de la leçon.

\begin{itemize}
\item \textbf{La puissance de la loi de Newton} : une loi universelle, valable pour la mangue qui tombe comme pour la Lune, permet, combinée au théorème de la sphère homogène, de déterminer la masse d'une planète entière à partir d'une simple mesure de $g_0$ au sol. C'est exactement la méthode historique rendue possible après la mesure de $G$ par Cavendish en 1798.
\item \textbf{La décroissance en $1/r^2$} : le champ de gravitation diminue avec l'altitude selon $g(h) = g_0\,\dfrac{R_T^2}{(R_T+h)^2}$. Cette décroissance est très lente à l'échelle humaine : au sommet du mont Cameroun, $g$ n'a baissé que de $0{,}13\,\%$, ce qui justifie l'approximation du \cle{champ uniforme} au voisinage du sol.
\item \textbf{Le sens physique de l'« impesanteur »} : à l'altitude de l'ISS, $g$ vaut encore $\SI{8,69}{N.kg^{-1}}$, soit près de $89\,\%$ de $g_0$ ! Les astronautes ne flottent donc pas parce que la gravité aurait disparu, mais parce qu'ils sont en chute libre permanente autour de la Terre.
\item \textbf{L'esprit critique scientifique} : comparer un résultat à une valeur de référence, quantifier l'écart relatif et identifier ses sources (modèle de la sphère homogène, aplatissement de la Terre, rotation, précision des mesures) est une compétence essentielle du physicien comme de l'ingénieur.
\end{itemize}

\begin{propriete}[Champ uniforme et distinction Poids / Gravitation]
Au voisinage immédiat de la surface terrestre ($h \ll R_T$), le champ de gravitation est pratiquement constant en intensité ($g \approx g_0$), en direction (verticale locale) et en sens (vers le centre de la Terre) : on parle de \cle{champ de pesanteur uniforme}.
\begin{itemize}
\item \textbf{Force de gravitation} $\vec{F}_g = m\,\vec{g}$ : force d'attraction exercée par la Terre (dépend de $M_T$, $R_T$, $G$).
\item \textbf{Poids} $\vec{P} = m\,\vec{g}_{\text{eff}}$ : force que l'objet exerce sur son support (mesurée par un dynamomètre).
\end{itemize}
Du fait de la rotation terrestre, le référentiel terrestre n'est pas galiléen. Le poids $\vec{P}$ diffère de la force de gravitation $\vec{F}_g$ de la force d'inertie centrifuge : $\vec{P} = \vec{F}_g + \vec{F}_{\text{cf}}$. À l'équateur, $P = F_g - m\,\omega^2 R_T$ ; aux pôles, $\vec{P} = \vec{F}_g$. Cette distinction est \emph{exigible au Bac} pour expliquer la variation de $g_0$ avec la latitude.
\end{propriete}

\begin{savaistu}[Henry Cavendish, l'homme qui a « pesé la Terre »]
En 1798, plus d'un siècle après Newton, le physicien anglais Henry Cavendish mesura la constante $G$ à l'aide d'une balance de torsion extrêmement sensible, capable de détecter l'attraction gravitationnelle entre des sphères de plomb de quelques kilogrammes. Connaissant $G$, $g_0$ et $R_T$, il en déduisit la masse de la Terre avec une précision remarquable pour l'époque. Son expérience est d'ailleurs intitulée dans son mémoire : « Expériences pour déterminer la densité de la Terre ». Aujourd'hui, les mêmes principes permettent de déterminer la masse des planètes lointaines et même des exoplanètes, en observant le mouvement de leurs satellites (lois de Kepler + loi de Newton).
\end{savaistu}

\newpage

\coursec{Méthodes et exercices résolus}

\coursec{Interaction gravitationnelle et champ de gravitation}

\courssub{La loi de l'attraction universelle de Newton}

\begin{propriete}[Loi de l'attraction universelle (Newton)]
Deux masses ponctuelles $m_A$ et $m_B$ (ou deux sphères homogènes de centres $A$ et $B$) s'attirent mutuellement avec des forces de même intensité :
\[ F_{A/B} = F_{B/A} = G\,\dfrac{m_A\,m_B}{d^2}, \]
où $d = AB$ est la distance entre les centres (en \si{\metre}), $G = 6,67\times 10^{-11}\ \mathrm{N\cdot m^2\cdot kg^{-2}}$ est la constante de gravitation universelle.
\textbf{Forme vectorielle :} $\vec{F}_{A/B} = -G\,\dfrac{m_A m_B}{d^3}\,\vec{AB} = G\,\dfrac{m_A m_B}{d^2}\,\vec{u}_{BA}$ (attractive, dirigée de $A$ vers $B$).
\textbf{Analyse dimensionnelle :} $[G] = \dfrac{[F][d^2]}{[m]^2} = \mathrm{N\cdot m^2\cdot kg^{-2}}$. \emph{(Démonstration non exigible au Bac, mais la loi et son utilisation le sont.)}
\end{propriete}

\begin{methode}[Calculer une force d'attraction gravitationnelle]
\begin{enumerate}
\item \textbf{Identifier} les deux corps : leurs masses $m_A$ et $m_B$ (en kilogrammes) et la distance $d$ entre leurs centres (en mètres). Convertir toutes les données en unités SI avant tout calcul.
\item \textbf{Schématiser} la situation : placer les deux corps, tracer le segment joignant leurs centres, représenter les deux forces $\vec{F}_{A/B}$ (exercée par $B$ sur $A$) et $\vec{F}_{B/A}$ (exercée par $A$ sur $B$), dirigées \emph{l'une vers l'autre} le long de ce segment (forces attractives).
\item \textbf{Écrire} la loi de Newton :
\[ F = G\,\dfrac{m_A\, m_B}{d^2} \quad \text{avec} \quad G = 6{,}67\times 10^{-11}\ \mathrm{N\cdot m^2\cdot kg^{-2}}. \]
\item \textbf{Calculer} en séparant les puissances de 10, puis conclure par une phrase avec l'unité et un nombre raisonnable de chiffres significatifs.
\item \textbf{Vérifier} la cohérence : les deux forces ont la même norme (principe des actions réciproques) ; la force diminue si la distance augmente.
\end{enumerate}
\end{methode}

\begin{exoresolu}[Attraction entre la Terre et la Lune]
Énoncé. La Terre (masse $M_T = 5{,}98\times 10^{24}\ \mathrm{kg}$) et la Lune (masse $M_L = 7{,}34\times 10^{22}\ \mathrm{kg}$) ont leurs centres distants de $d = 3{,}84\times 10^{8}\ \mathrm{m}$. On donne $G = 6{,}67\times 10^{-11}\ \mathrm{N\cdot m^2\cdot kg^{-2}}$.
\begin{enumerate}
\item Représenter sur un schéma les forces d'interaction gravitationnelle entre la Terre et la Lune.
\item Calculer l'intensité de la force exercée par la Terre sur la Lune.
\item Comparer à la force exercée par la Lune sur la Terre. Justifier.
\end{enumerate}
\tcblower
\textbf{1. Schéma.} Les deux forces sont portées par la droite joignant les centres, attractives, de même norme, de sens opposés :
\begin{popfigure}
\begin{tikzpicture}[scale=1]
\shade[ball color=popblue!70] (0,0) circle (0.5);
\node[below] at (0,-0.55) {Terre};
\shade[ball color=popmuted!60] (5,0) circle (0.25);
\node[below] at (5,-0.35) {Lune};
\draw[popline,dashed] (0.5,0) -- (4.75,0);
\draw[-{Stealth[length=3mm]},poporange,very thick] (0.55,0.25) -- (2.2,0.25) node[midway,above,popdark]{$\vec{F}_{L/T}$};
\draw[-{Stealth[length=3mm]},poporange,very thick] (4.45,-0.25) -- (2.8,-0.25) node[midway,below,popdark]{$\vec{F}_{T/L}$};
\end{tikzpicture}
\legende{Forces d'attraction mutuelle Terre--Lune (3\textsuperscript{e} loi de Newton).}
\end{popfigure}
\textbf{2. Intensité de la force.} La loi de l'attraction universelle s'écrit :
\[ F_{T/L} = G\,\dfrac{M_T\, M_L}{d^2}. \]
Application numérique, en regroupant les puissances de 10 :
\begin{align*}
F_{T/L} &= 6{,}67\times 10^{-11}\times \dfrac{5{,}98\times 10^{24}\times 7{,}34\times 10^{22}}{(3{,}84\times 10^{8})^2} \\
&= 6{,}67\times 10^{-11}\times \dfrac{43{,}9\times 10^{46}}{14{,}75\times 10^{16}} \\
&= 6{,}67\times 10^{-11}\times 2{,}976\times 10^{30} \\
&\approx 1{,}98\times 10^{20}\ \mathrm{N}.
\end{align*}
\textbf{Conclusion :} la Terre exerce sur la Lune une force attractive d'intensité $F_{T/L} \approx 1{,}98\times 10^{20}\ \mathrm{N}$. C'est cette force qui maintient la Lune sur son orbite.

\textbf{3. Comparaison.} D'après le principe des actions réciproques (troisième loi de Newton), $\vec{F}_{T/L} = -\vec{F}_{L/T}$ : les deux forces ont \textbf{même norme} $1{,}98\times 10^{20}\ \mathrm{N}$, même direction et des sens opposés. La Lune attire la Terre aussi fort que la Terre attire la Lune, même si la Terre est bien plus massive : c'est la force qui est commune aux deux corps, pas l'accélération qu'elle leur communique.
\end{exoresolu}

\courssub{Champ de gravitation créé par un point matériel}

\begin{definition}[Champ de gravitation]
Le \cle{champ de gravitation} $\vec{\mathcal{G}}$ en un point $M$ de l'espace est la force de gravitation qu'éprouverait une masse test unitaire ($1\ \mathrm{kg}$) placée en $M$, toutes les autres masses restant fixes. C'est un champ vectoriel : $\vec{\mathcal{G}}(M) = \dfrac{\vec{F}(M)}{m_{\text{test}}}$. Son unité SI est le \si{\newton\per\kilogram} ($\mathrm{N\cdot kg^{-1}}$), équivalent au \si{\metre\per\second\squared}.
\end{definition}

\begin{propriete}[Champ créé par un point matériel (ou sphère homogène, à l'extérieur)]
Un point matériel $S$ de masse $M$ crée en un point $M$ (distance $r = SM$) un champ de gravitation :
\[ \vec{\mathcal{G}}(M) = -\,G\,\dfrac{M}{r^2}\,\vec{u}, \qquad \mathcal{G}(M) = G\,\dfrac{M}{r^2}, \]
où $\vec{u}$ est le vecteur unitaire dirigé de $S$ vers $M$. Le signe $-$ indique que le champ est \textbf{attractif} (dirigé de $M$ vers $S$, centripète).
Les \textbf{lignes de champ} sont des demi-droites radiales convergentes vers $S$ ; les \textbf{surfaces équipotentielles} sont des sphères centrées sur $S$.
\end{propriete}

\begin{methode}[Représenter $\vec{\mathcal{G}}$ créé par un point matériel ou une sphère homogène]
\begin{enumerate}
\item \textbf{Placer} le corps source $S$ de masse $M$ (point matériel ou sphère homogène : dans les deux cas, tout se passe \emph{à l'extérieur} comme si toute la masse était concentrée au centre) et le point $M$ où l'on veut le champ.
\item \textbf{Tracer} la demi-droite $[SM)$ : le champ $\vec{\mathcal{G}}(M)$ est porté par cette droite, \textbf{dirigé de $M$ vers $S$} (champ attractif, centripète).
\item \textbf{Écrire} l'expression vectorielle en introduisant le vecteur unitaire $\vec{u}$ dirigé de $S$ vers $M$ :
\[ \vec{\mathcal{G}}(M) = -\,G\,\dfrac{M}{r^2}\,\vec{u}, \qquad \mathcal{G} = G\dfrac{M}{r^2} \quad (\mathrm{N\cdot kg^{-1}}). \]
\item \textbf{Choisir une échelle} (par exemple $1\ \mathrm{cm} \leftrightarrow 2\ \mathrm{N\cdot kg^{-1}}$) et tracer le vecteur en respectant la longueur calculée.
\item \textbf{Exploiter la symétrie} : les lignes de champ sont des demi-droites radiales convergentes vers $S$ ; le champ a même norme en tout point d'une sphère de centre $S$.
\end{enumerate}
\end{methode}

\begin{exoresolu}[Champ de gravitation créé par un astre]
Énoncé. Un astre sphérique homogène $S$, de masse $M = 6{,}0\times 10^{23}\ \mathrm{kg}$, est assimilé à un point matériel placé en son centre $O$. Soit un point $P$ situé à la distance $r = 4{,}0\times 10^{6}\ \mathrm{m}$ de $O$. On donne $G = 6{,}67\times 10^{-11}\ \mathrm{N\cdot m^2\cdot kg^{-2}}$.
\begin{enumerate}
\item Donner l'expression vectorielle du champ de gravitation $\vec{\mathcal{G}}(P)$ créé par l'astre en $P$. Le représenter à l'échelle $1\ \mathrm{cm} \leftrightarrow 1{,}0\ \mathrm{N\cdot kg^{-1}}$.
\item Calculer la norme de $\vec{\mathcal{G}}(P)$.
\item Quelle est la force subie par un satellite de masse $m = 500\ \mathrm{kg}$ placé en $P$ ?
\end{enumerate}
\tcblower
\textbf{1. Expression vectorielle et représentation.} Soit $\vec{u}$ le vecteur unitaire dirigé de $O$ vers $P$. Le champ créé par un corps à symétrie sphérique en un point extérieur est :
\[ \vec{\mathcal{G}}(P) = -\,G\,\dfrac{M}{r^2}\,\vec{u}. \]
Le signe $-$ traduit le caractère attractif : $\vec{\mathcal{G}}(P)$ est dirigé de $P$ vers $O$.
\begin{popfigure}
\begin{tikzpicture}[scale=1]
\shade[ball color=poppurple!70] (0,0) circle (0.35);
\node[below left,popdark] at (-0.1,-0.3) {$O$ (astre)};
\filldraw[popdark] (4,0) circle (1.5pt) node[below right]{$P$};
\draw[-{Stealth[length=2.5mm]},popline,dashed] (0.4,0) -- (4,0);
\draw[-{Stealth[length=3mm]},poporange,very thick] (4,0.3) -- (1.5,0.3) node[midway,above,popdark]{$\vec{\mathcal{G}}(P)$};
\draw[-{Stealth[length=2.5mm]},popblue,thick] (0.45,-0.35) -- (1.45,-0.35) node[midway,below,popdark]{$\vec{u}$};
\end{tikzpicture}
\legende{Champ de gravitation créé par un astre sphérique en un point extérieur $P$.}
\end{popfigure}
\textbf{2. Norme du champ.}
\begin{align*}
\mathcal{G}(P) &= G\,\dfrac{M}{r^2} = 6{,}67\times 10^{-11}\times \dfrac{6{,}0\times 10^{23}}{(4{,}0\times 10^{6})^2} \\
&= 6{,}67\times 10^{-11}\times \dfrac{6{,}0\times 10^{23}}{1{,}6\times 10^{13}} \\
&= 6{,}67\times 10^{-11}\times 3{,}75\times 10^{10} \approx 2{,}5\ \mathrm{N\cdot kg^{-1}}.
\end{align*}
À l'échelle $1\ \mathrm{cm} \leftrightarrow 1{,}0\ \mathrm{N\cdot kg^{-1}}$, le vecteur mesure $2{,}5\ \mathrm{cm}$, fléché de $P$ vers $O$.

\textbf{3. Force sur le satellite.} Par définition du champ de gravitation, $\vec{F} = m\,\vec{\mathcal{G}}(P)$, donc :
\[ F = m\,\mathcal{G}(P) = 500\times 2{,}5 = 1{,}25\times 10^{3}\ \mathrm{N}. \]
\textbf{Conclusion :} le satellite subit une force attractive d'intensité $F \approx 1{,}3\times 10^{3}\ \mathrm{N}$, dirigée vers le centre de l'astre. On retrouve le même résultat par la loi de Newton $F = GMm/r^2$, ce qui valide la cohérence des deux approches.
\end{exoresolu}

\courssub{Champ créé par une sphère homogène}

\begin{propriete}[Théorème de Newton pour une sphère homogène (résultat admis)]
\begin{itemize}
\item \textbf{À l'extérieur} ($r \geq R$) : une sphère homogène (pleine ou creuse) de masse $M$ et de rayon $R$ crée le même champ qu'un point matériel de masse $M$ placé en son centre $O$ :
\[ \mathcal{G}(M) = G\,\dfrac{M}{r^2} \quad (r = OM \geq R). \]
\item \textbf{À la surface} ($r = R$) : $\mathcal{G}_0 = \dfrac{GM}{R^2}$.
\item \textbf{À l'intérieur d'une sphère creuse homogène} : le champ est \textbf{nul} en tout point.
\item \textbf{À l'intérieur d'une sphère pleine homogène} (complément, non exigible au Bac mais utile) : $\mathcal{G}(r) = \dfrac{GM}{R^3}\,r$ (proportionnel à $r$, nul au centre).
\end{itemize}
\end{propriete}

\begin{methode}[Sphère homogène : se ramener à un point matériel]
\begin{enumerate}
\item \textbf{Vérifier la condition d'application} : le point $M$ doit être \emph{à l'extérieur} de la sphère ($r \geq R$). Le résultat vaut pour une sphère pleine homogène \emph{et} pour une sphère creuse homogène (couche sphérique).
\item \textbf{Remplacer} mentalement la sphère par un point matériel de même masse $M$ placé en son centre $O$.
\item \textbf{Appliquer} la formule du point matériel avec $r = OM$ (distance au \emph{centre}, pas à la surface !) :
\[ \mathcal{G}(M) = G\,\dfrac{M}{r^2} \qquad (r \geq R). \]
\item \textbf{Cas particulier} : à la surface ($r = R$), $\mathcal{G}_0 = \dfrac{GM}{R^2}$.
\item \textbf{À l'intérieur} d'une sphère creuse homogène : le champ est \emph{nul} en tout point (résultat admis, mais exigible comme constat).
\end{enumerate}
\textbf{Réflexe :} toujours mesurer $r$ depuis le centre $O$ ; confondre $r$ avec l'altitude $h = r - R$ est l'erreur la plus fréquente.
\end{methode}

\begin{exoresolu}[Champ à la surface d'une planète]
Énoncé. La planète Mars est assimilée à une sphère homogène de masse $M = 6{,}42\times 10^{23}\ \mathrm{kg}$ et de rayon $R = 3{,}39\times 10^{6}\ \mathrm{m}$. On donne $G = 6{,}67\times 10^{-11}\ \mathrm{N\cdot m^2\cdot kg^{-2}}$.
\begin{enumerate}
\item Justifier que le champ de gravitation à la surface de Mars s'écrit $g_{0M} = \dfrac{GM}{R^2}$.
\item Calculer $g_{0M}$ et comparer à $g_0 = 9{,}81\ \mathrm{N\cdot kg^{-1}}$ à la surface de la Terre.
\item Un robot de masse $m = 180\ \mathrm{kg}$ est posé au sol martien. Calculer la force de gravitation qu'il subit.
\end{enumerate}
\tcblower
\textbf{1. Justification.} Mars étant assimilée à une sphère homogène, son champ de gravitation en tout point \emph{extérieur} est le même que celui d'un point matériel de masse $M$ placé en son centre $O$. Un point de la surface est à la distance $r = R$ du centre, donc :
\[ g_{0M} = \mathcal{G}(R) = G\,\dfrac{M}{R^2}. \]
\textbf{2. Calcul.}
\begin{align*}
g_{0M} &= 6{,}67\times 10^{-11}\times \dfrac{6{,}42\times 10^{23}}{(3{,}39\times 10^{6})^2} \\
&= 6{,}67\times 10^{-11}\times \dfrac{6{,}42\times 10^{23}}{1{,}149\times 10^{13}} \\
&= 6{,}67\times 10^{-11}\times 5{,}59\times 10^{10} \approx 3{,}73\ \mathrm{N\cdot kg^{-1}}.
\end{align*}
Comparaison : $\dfrac{g_{0M}}{g_0} = \dfrac{3{,}73}{9{,}81} \approx 0{,}38$. La gravité à la surface de Mars vaut environ $38\ \%$ de celle de la Terre : un objet y pèse environ $2{,}6$ fois moins.

\textbf{3. Force sur le robot.}
\[ F = m\,g_{0M} = 180\times 3{,}73 \approx 6{,}7\times 10^{2}\ \mathrm{N}. \]
\textbf{Conclusion :} le robot subit sur Mars une force de gravitation d'environ $671\ \mathrm{N}$, contre $180\times 9{,}81 \approx 1{,}77\times 10^{3}\ \mathrm{N}$ sur Terre. Sa \emph{masse} ($180\ \mathrm{kg}$) est en revanche identique sur les deux planètes : la masse est invariante, pas la force de gravitation.
\end{exoresolu}

\courssub{Champ de gravitation terrestre et variation avec l'altitude}

\begin{propriete}[Champ de gravitation terrestre]
La Terre est modélisée par une sphère homogène de masse $M_T$ et de rayon $R_T \approx 6,38\times 10^{6}\ \mathrm{m}$.
\begin{itemize}
\item À la surface ($h=0$) : $g_0 = \dfrac{GM_T}{R_T^2} \approx 9{,}81\ \mathrm{N\cdot kg^{-1}}$.
\item À l'altitude $h$ (distance au centre $r = R_T + h$) :
\[ g(h) = \dfrac{GM_T}{(R_T+h)^2} = g_0\,\dfrac{R_T^2}{(R_T+h)^2}. \]
\end{itemize}
Cette formule montre que $g$ diminue avec l'altitude, mais reste non nul quelle que soit la distance.
\end{propriete}

\begin{methode}[Établir et utiliser $g(h) = g_0\,\dfrac{R_T^2}{(R_T+h)^2}$]
\begin{enumerate}
\item \textbf{Poser le cadre} : la Terre est une sphère homogène de masse $M_T$ et de rayon $R_T$ ; un point à l'altitude $h$ est à la distance $r = R_T + h$ du centre.
\item \textbf{Écrire les deux expressions} du champ :
\[ g(h) = \dfrac{G M_T}{(R_T+h)^2} \quad \text{et} \quad g_0 = \dfrac{G M_T}{R_T^2} \quad (\text{à la surface}). \]
\item \textbf{Éliminer $GM_T$} en faisant le rapport membre à membre :
\[ \dfrac{g(h)}{g_0} = \dfrac{R_T^2}{(R_T+h)^2} \quad\Longrightarrow\quad g(h) = g_0\,\dfrac{R_T^2}{(R_T+h)^2}. \]
\item \textbf{Application numérique} : convertir $h$ en mètres, calculer le rapport $\dfrac{R_T}{R_T+h}$ \emph{avant} de l'élever au carré (moins d'erreurs d'arrondi).
\item \textbf{Inversion} : pour trouver $h$ connaissant $g(h)$, isoler : $R_T + h = R_T\sqrt{\dfrac{g_0}{g(h)}}$, puis $h = R_T\left(\sqrt{\dfrac{g_0}{g(h)}} - 1\right)$.
\end{enumerate}
\end{methode}

\begin{exoresolu}[Gravité à bord de la station spatiale]
Énoncé. La Station spatiale internationale (ISS) orbite à l'altitude $h = 400\ \mathrm{km}$. On donne $g_0 = 9{,}81\ \mathrm{N\cdot kg^{-1}}$ et $R_T = 6{,}38\times 10^{6}\ \mathrm{m}$.
\begin{enumerate}
\item Établir l'expression de $g(h)$ en fonction de $g_0$, $R_T$ et $h$.
\item Calculer $g$ à l'altitude de l'ISS.
\item Un élève affirme : « Les astronautes flottent parce qu'à cette altitude, la gravité est nulle. » Commenter.
\end{enumerate}
\tcblower
\textbf{1. Établissement.} La Terre étant assimilée à une sphère homogène, le champ à la distance $r = R_T + h$ de son centre vaut $g(h) = \dfrac{GM_T}{(R_T+h)^2}$. À la surface : $g_0 = \dfrac{GM_T}{R_T^2}$, d'où $GM_T = g_0 R_T^2$. En reportant :
\[ g(h) = g_0\,\dfrac{R_T^2}{(R_T+h)^2}. \]
\textbf{2. Application numérique.} Conversion : $h = 400\ \mathrm{km} = 4{,}00\times 10^{5}\ \mathrm{m}$, donc $R_T + h = 6{,}38\times 10^{6} + 0{,}40\times 10^{6} = 6{,}78\times 10^{6}\ \mathrm{m}$.
\begin{align*}
g(h) &= 9{,}81\times \left(\dfrac{6{,}38\times 10^{6}}{6{,}78\times 10^{6}}\right)^2 \\
&= 9{,}81\times (0{,}9410)^2 \\
&= 9{,}81\times 0{,}8855 \approx 8{,}69\ \mathrm{N\cdot kg^{-1}}.
\end{align*}
\textbf{Conclusion :} à $400\ \mathrm{km}$ d'altitude, $g \approx 8{,}69\ \mathrm{N\cdot kg^{-1}}$.

\textbf{3. Commentaire.} L'affirmation est \textbf{fausse} : à l'altitude de l'ISS, la gravité vaut encore environ $89\ \%$ de sa valeur au sol. Si les astronautes « flottent », ce n'est pas parce que la gravité a disparu, mais parce que la station et son équipage sont en \emph{chute libre permanente} autour de la Terre : tout tombe ensemble, d'où la sensation d'impesanteur. La gravité est d'ailleurs indispensable : c'est elle qui joue le rôle de force centripète maintenant la station sur son orbite.
\end{exoresolu}

\begin{savaistu}[Le Cameroun et l'espace : télécommunications et observation]
Le Cameroun, situé près de l'équateur, est une zone stratégique pour les satellites géostationnaires (MTN, Orange, Camtel utilisent des capacités sur des satellites comme \emph{Intelsat}, \emph{Eutelsat}, \emph{RASCOM}). L'altitude géostationnaire est $h \approx 35\,786\ \mathrm{km}$ ($r \approx 6,6\,R_T$), où $g \approx 0,22\ \mathrm{N\cdot kg^{-1}}$.
Les missions spatiales \emph{GRACE} et \emph{GOCE} ont cartographié le champ de gravité terrestre avec une précision centimétrique, révélant les variations de masse d'eau (nappes phréatiques, lacs comme Lom Pangar, barrage de Nachtigal) et aidant à la gestion des ressources hydrauliques et à la prévention des inondations. La gravimétrie spatiale est ainsi un outil de développement pour l'agriculture et l'énergie hydraulique au Cameroun.
\end{savaistu}

\courssub{Champ uniforme au voisinage du sol ; poids et force de gravitation}

\begin{propriete}[Champ uniforme au voisinage du sol]
Si l'altitude $h$ est faible devant le rayon terrestre ($h \ll R_T$, typiquement $h < 10\ \mathrm{km}$ devant $R_T \approx 6\,380\ \mathrm{km}$), alors $R_T + h \approx R_T$.
On en déduit $g(h) \approx g_0$ : le champ garde une norme quasi constante ; de plus, toutes les verticales locales étant quasi parallèles à l'échelle d'une région, la direction du champ est constante.
\textbf{Conclusion :} au voisinage du sol, le champ de gravitation est \emph{uniforme} (même direction, même sens, même norme) ; les lignes de champ sont des droites verticales parallèles, dirigées vers le bas.
\end{propriete}

\begin{definition}[Poids et force de gravitation terrestre]
La \cle{force de gravitation} $\vec{F} = -\dfrac{GM_T m}{r^2}\vec{u}$ est l'attraction exercée par la Terre, dirigée vers son centre.
Le \cle{poids} $\vec{P} = m\vec{g}$ est la force mesurée dans le référentiel terrestre (lié à la Terre tournante) : il résulte de la force de gravitation \emph{et} de la force d'inertie d'entraînement (force centrifuge) due à la rotation de la Terre sur elle-même.
\begin{itemize}
\item Le poids n'est donc pas rigoureusement dirigé vers le centre de la Terre (sauf aux pôles et à l'équateur).
\item Sa norme dépend légèrement de la latitude (maximale aux pôles, minimale à l'équateur : correction $\approx 0,3\ \%$).
\item Dans les exercices de mécanique au voisinage du sol, on confond les deux ($\vec{P} \approx \vec{F}$), ce qui est une excellente approximation.
\end{itemize}
\end{definition}

\begin{methode}[Justifier l'approximation du champ uniforme]
\begin{enumerate}
\item \textbf{Comparer} l'altitude $h$ au rayon terrestre : si $h \ll R_T$ (typiquement $h < 10\ \mathrm{km}$ devant $R_T \approx 6\,380\ \mathrm{km}$), alors $R_T + h \approx R_T$.
\item \textbf{En déduire} $g(h) \approx g_0$ : le champ garde une norme quasi constante ; de plus, toutes les verticales locales étant quasi parallèles à l'échelle d'une région, la direction du champ est constante.
\item \textbf{Conclure} : au voisinage du sol, le champ de gravitation est \emph{uniforme} (même direction, même sens, même norme) ; les lignes de champ sont des droites verticales parallèles, dirigées vers le bas.
\item \textbf{Distinction poids / gravitation} : le \cle{poids} $\vec{P} = m\vec{g}$ est la force de gravitation terrestre \emph{corrigée} des effets de la rotation de la Terre (force d'inertie d'entraînement). À l'équateur, cette correction est maximale mais reste faible (environ $0,3\ \%$). Au lycée, on identifie souvent $\vec{P}$ à la force de gravitation.
\end{enumerate}
\end{methode}

\begin{exoresolu}[Le champ est-il uniforme entre Yaoundé et le sommet du Mont Cameroun ?]
Énoncé. Le Mont Cameroun culmine à $h = 4{,}10\ \mathrm{km}$ d'altitude. On donne $g_0 = 9{,}81\ \mathrm{N\cdot kg^{-1}}$ et $R_T = 6{,}38\times 10^{6}\ \mathrm{m}$.
\begin{enumerate}
\item Calculer la variation relative $\dfrac{g_0 - g(h)}{g_0}$ entre le niveau de la mer et le sommet du Mont Cameroun. Conclure sur l'uniformité du champ.
\item Un randonneur de masse $m = 75\ \mathrm{kg}$ gravit le sommet. Calculer la variation de la force de gravitation qu'il subit.
\item Rappeler la distinction entre le poids et la force de gravitation terrestre.
\end{enumerate}
\tcblower
\textbf{1. Variation relative.} Avec $h = 4{,}10\times 10^{3}\ \mathrm{m}$ :
\begin{align*}
\dfrac{g(h)}{g_0} &= \left(\dfrac{R_T}{R_T+h}\right)^2 = \left(\dfrac{6{,}38\times 10^{6}}{6{,}38\times 10^{6}+4{,}10\times 10^{3}}\right)^2 \\
&= \left(\dfrac{6{,}3800\times 10^{6}}{6{,}3841\times 10^{6}}\right)^2 = (0{,}99936)^2 \approx 0{,}99872.
\end{align*}
Donc $\dfrac{g_0 - g(h)}{g_0} = 1 - 0{,}99872 \approx 1{,}3\times 10^{-3}$, soit environ $0{,}13\ \%$.
\textbf{Conclusion :} la variation est inférieure à $1\ \%$ : à l'échelle du relief camerounais, le champ de gravitation peut être considéré comme \textbf{uniforme} ($g \approx g_0$, direction verticale locale partout quasi parallèle).

\textbf{2. Variation de la force.} Au niveau de la mer : $F_0 = m g_0 = 75\times 9{,}81 = 735{,}75\ \mathrm{N}$. Au sommet : $F = m g(h) = 75\times 9{,}81\times 0{,}99872 \approx 734{,}8\ \mathrm{N}$. La variation est :
\[ \Delta F = F_0 - F \approx 0{,}9\ \mathrm{N}, \]
soit l'équivalent du poids d'une masse d'environ $95\ \mathrm{g}$ : totalement imperceptible pour le randonneur.

\textbf{3. Distinction.} La \cle{force de gravitation} $\vec{F} = -\dfrac{GM_T m}{r^2}\vec{u}$ est l'attraction exercée par la Terre, dirigée vers son centre. Le \cle{poids} $\vec{P} = m\vec{g}$ est la force mesurée dans le référentiel terrestre : il résulte de la force de gravitation \emph{et} de la force d'inertie d'entraînement due à la rotation de la Terre sur elle-même. Le poids n'est donc pas rigoureusement dirigé vers le centre de la Terre (sauf aux pôles et à l'équateur) et sa norme dépend légèrement de la latitude. Dans les exercices de mécanique au voisinage du sol, on confond les deux, ce qui est une excellente approximation.
\end{exoresolu}

\courssub{Synthèse : Points d'équilibre et comparaisons}

\begin{methode}[Problèmes de synthèse : distances, rapports, points d'équilibre]
\begin{enumerate}
\item \textbf{Traduire} géométriquement l'énoncé : poser les distances inconnues ($x$, $d-x$, etc.) sur un schéma.
\item \textbf{Écrire l'égalité des normes} quand deux champs ou deux forces se compensent : $\dfrac{GM_1}{r_1^2} = \dfrac{GM_2}{r_2^2}$ ; $G$ se simplifie toujours.
\item \textbf{Résoudre} l'équation obtenue (souvent $\dfrac{r_1}{r_2} = \sqrt{\dfrac{M_1}{M_2}}$) en prenant la racine carrée du rapport des masses.
\item \textbf{Vérifier} que la solution est physiquement acceptable (position entre les deux corps, distance positive...).
\item \textbf{Contrôler} l'homogénéité de chaque formule avant l'application numérique.
\end{enumerate}
\end{methode}

\begin{exoresolu}[Point d'équilibre gravitationnel entre la Terre et la Lune]
Énoncé. On cherche le point $P$ du segment joignant le centre de la Terre au centre de la Lune où les champs de gravitation des deux astres se compensent exactement. On donne $M_T = 81\,M_L$ et la distance Terre--Lune $d = 3{,}84\times 10^{8}\ \mathrm{m}$.
\begin{enumerate}
\item Faire un schéma et justifier que $P$ est situé entre les deux astres.
\item Déterminer la position de $P$ par rapport au centre de la Terre.
\end{enumerate}
\tcblower
\textbf{1. Schéma et justification.} Entre les deux astres, le champ terrestre $\vec{\mathcal{G}}_T$ pointe vers la Terre et le champ lunaire $\vec{\mathcal{G}}_L$ vers la Lune : ils sont de \emph{sens opposés} et peuvent s'annuler. À l'extérieur du segment, les deux champs pointent du même côté (vers l'astre le plus proche... en fait tous deux vers le segment) : ils s'ajoutent et ne peuvent se compenser. Le point $P$ est donc nécessairement entre les deux corps, et plus proche de la Lune (moins massive).
\begin{popfigure}
\begin{tikzpicture}[scale=1]
\shade[ball color=popblue!70] (0,0) circle (0.4);
\node[below] at (0,-0.45) {Terre};
\shade[ball color=popmuted!60] (7,0) circle (0.2);
\node[below] at (7,-0.3) {Lune};
\filldraw[poporange] (6.3,0) circle (2pt) node[above,popdark]{$P$};
\draw[popline] (0.45,0) -- (6.8,0);
\draw[-{Stealth[length=2.5mm]},popblue,thick] (6.3,0.35) -- (5.1,0.35) node[midway,above,popdark]{$\vec{\mathcal{G}}_T$};
\draw[-{Stealth[length=2.5mm]},poppurple,thick] (6.3,-0.35) -- (7.0,-0.35) node[midway,below,popdark]{$\vec{\mathcal{G}}_L$};
\draw[{Stealth[length=2mm]}-{Stealth[length=2mm]},popdark] (0,-0.9) -- (6.3,-0.9) node[midway,below]{$x$};
\draw[{Stealth[length=2mm]}-{Stealth[length=2mm]},popdark] (6.3,-0.9) -- (7,-0.9) node[midway,below]{$d-x$};
\end{tikzpicture}
\legende{Point d'équilibre gravitationnel Terre--Lune.}
\end{popfigure}
\textbf{2. Position de $P$.} Notons $x = TP$ ; alors $PL = d - x$. La compensation s'écrit $\mathcal{G}_T = \mathcal{G}_L$ :
\[ \dfrac{GM_T}{x^2} = \dfrac{GM_L}{(d-x)^2} \quad\Longrightarrow\quad \dfrac{(d-x)^2}{x^2} = \dfrac{M_L}{M_T} = \dfrac{1}{81}. \]
En prenant la racine carrée (les distances sont positives) :
\[ \dfrac{d-x}{x} = \dfrac{1}{9} \quad\Longrightarrow\quad 9(d-x) = x \quad\Longrightarrow\quad 9d = 10x \quad\Longrightarrow\quad x = \dfrac{9}{10}d. \]
Application numérique :
\[ x = 0{,}9\times 3{,}84\times 10^{8} = 3{,}46\times 10^{8}\ \mathrm{m}. \]
\textbf{Conclusion :} le point d'équilibre gravitationnel est situé à $x \approx 3{,}46\times 10^{8}\ \mathrm{m}$ du centre de la Terre, soit à environ $3{,}8\times 10^{7}\ \mathrm{m}$ de la Lune : il est bien entre les deux astres, neuf fois plus loin de la Terre que de la Lune, ce qui est cohérent avec le rapport des masses ($\sqrt{81} = 9$).
\end{exoresolu}

\begin{attention}[Les 5 erreurs qui coûtent des points au Bac]
\begin{enumerate}
\item \textbf{Oublier le carré sur la distance.} La loi de Newton est en $1/d^2$, pas en $1/d$ : écrire $F = GMm/d$ est une erreur de cours éliminatoire. Vérifie chaque fois : « la force est-elle divisée par $d^2$ ? ».
\item \textbf{Confondre distance au centre et altitude.} Dans $g(h) = \dfrac{GM_T}{(R_T+h)^2}$, le dénominateur contient $(R_T + h)^2$ et non $h^2$. À la surface, $h = 0$ et l'on retrouve $g_0 = GM_T/R_T^2$ : fais ce test de cohérence systématiquement.
\item \textbf{Mal convertir les unités.} Les distances doivent être en \emph{mètres} : $400\ \mathrm{km} = 4{,}00\times 10^{5}\ \mathrm{m}$. Une distance laissée en kilomètres fausse le résultat d'un facteur $10^{6}$. Écris la conversion explicitement sur ta copie.
\item \textbf{Tracer le champ ou la force dans le mauvais sens.} Le champ de gravitation $\vec{\mathcal{G}}$ et la force $\vec{F}$ sont \emph{attractifs} : dirigés vers le corps qui les crée. Un vecteur fléché vers l'extérieur est faux. Précise aussi le vecteur unitaire $\vec{u}$ (de la source vers le point) pour justifier le signe $-$ dans $\vec{\mathcal{G}} = -\dfrac{GM}{r^2}\vec{u}$.
\item \textbf{Confondre masse et poids, ou oublier l'égalité des forces mutuelles.} La masse (kg) est invariante ; le poids (N) dépend de l'astre. Et quelle que soit la différence de masses, $F_{A/B} = F_{B/A}$ : écrire que « la Terre attire la Lune plus fort que l'inverse » est une faute de physique grave, sanctionnée par le correcteur.
\end{enumerate}
\end{attention}

\newpage

\coursec{Exercices}

\courssub{Je vérifie mes connaissances}

\begin{exercice}[QCM : la loi de Newton]
Pour chaque affirmation, choisir la (ou les) bonne(s) réponse(s) en justifiant brièvement.

\textbf{1.} La force gravitationnelle exercée par un corps A de masse $m_A$ sur un corps B de masse $m_B$, distants de $d$, a pour intensité :
\begin{itemize}
\item[a)] $F = G\,\dfrac{m_A + m_B}{d^2}$
\item[b)] $F = G\,\dfrac{m_A\, m_B}{d}$
\item[c)] $F = G\,\dfrac{m_A\, m_B}{d^2}$
\item[d)] $F = \dfrac{m_A\, m_B}{G\, d^2}$
\end{itemize}

\textbf{2.} Si l'on double la distance entre deux corps ponctuels, l'intensité de la force gravitationnelle entre eux est :
\begin{itemize}
\item[a)] divisée par 2
\item[b)] divisée par 4
\item[c)] multipliée par 2
\item[d)] inchangée
\end{itemize}

\textbf{3.} La constante de gravitation universelle $G$ vaut :
\begin{itemize}
\item[a)] $\SI{6,67e-11}{N.m^2.kg^{-2}}$
\item[b)] $\SI{9,81}{m.s^{-2}}$
\item[c)] $\SI{6,67e-11}{N.kg^2.m^{-2}}$
\item[d)] $\SI{6,02e23}{mol^{-1}}$
\end{itemize}

\textbf{4.} Le champ de gravitation créé par une sphère homogène de masse $M$ et de rayon $R$ en un point extérieur situé à la distance $r$ de son centre ($r \geq R$) est :
\begin{itemize}
\item[a)] $\mathcal{G} = G\,\dfrac{M}{r}$
\item[b)] $\mathcal{G} = G\,\dfrac{M}{r^2}$, dirigé vers le centre de la sphère
\item[c)] $\mathcal{G} = G\,\dfrac{M}{R^2}$, quelle que soit la valeur de $r$
\item[d)] $\mathcal{G} = G\,\dfrac{M}{r^2}$, dirigé vers l'extérieur de la sphère
\end{itemize}
\end{exercice}

\begin{exercice}[Vrai ou faux ?]
Répondre par \textbf{vrai} ou \textbf{faux} en justifiant chaque réponse par une phrase ou un calcul.

\textbf{1.} La Terre attire la Lune plus fortement que la Lune n'attire la Terre.

\textbf{2.} Le vecteur champ de gravitation créé par un point matériel en un point M est toujours dirigé de M vers le point matériel.

\textbf{3.} À la surface de la Terre, le champ de gravitation a pour valeur $g_0 = \dfrac{G\,M_T}{R_T^2}$.

\textbf{4.} Le champ de gravitation terrestre augmente lorsque l'altitude augmente.

\textbf{5.} Au sommet du mont Cameroun (altitude $\SI{4095}{m}$), le champ de gravitation est rigoureusement égal à sa valeur au niveau de la mer.
\end{exercice}

\begin{exercice}[Définitions et restitution de cours]
\textbf{1.} Énoncer la loi de l'attraction universelle de Newton (expression vectorielle et conditions d'application).

\textbf{2.} Définir le vecteur champ de gravitation $\vec{\mathcal{G}}(M)$ créé par un point matériel de masse $m$ en un point M. Donner son expression vectorielle en introduisant le vecteur unitaire $\vec{u}$ dirigé du point matériel vers M.

\textbf{3.} Énoncer le théorème de Gauss appliqué à une sphère homogène : que vaut le champ créé par une sphère homogène (pleine ou creuse) en un point extérieur ?

\textbf{4.} Établir, à partir de l'expression du champ de gravitation terrestre, la relation $g(h) = g_0\,\dfrac{R_T^2}{(R_T + h)^2}$.
\end{exercice}

\begin{exercice}[Calculs directs]
On donne : $G = \SI{6,67e-11}{N.m^2.kg^{-2}}$ ; $M_T = \SI{5,97e24}{kg}$ ; $R_T = \SI{6370}{km}$.

\textbf{1.} Calculer l'intensité de la force gravitationnelle exercée par la Terre sur un sac de riz de masse $m = \SI{50}{kg}$ posé au sol à Yaoundé (on assimilera la Terre à une sphère homogène).

\textbf{2.} Calculer la valeur $g_0$ du champ de gravitation à la surface de la Terre. Comparer à la valeur usuelle $\SI{9,81}{N.kg^{-1}}$.

\textbf{3.} Calculer l'intensité de la force gravitationnelle entre deux élèves de masse $m_1 = m_2 = \SI{60}{kg}$ assis à $d = \SI{1,0}{m}$ l'un de l'autre (on les assimilera à des points matériels). Commenter le résultat : pourquoi ne la ressent-on pas ?
\end{exercice}

\courssub{Je m'entraîne}

\begin{exercice}[Représentation du champ terrestre à différentes distances]
On donne : $g_0 = \SI{9,81}{N.kg^{-1}}$ ; $R_T = \SI{6370}{km}$.

\textbf{1.} Calculer la valeur du champ de gravitation terrestre aux points A, B et C situés respectivement aux distances $r_A = R_T$, $r_B = 2\,R_T$ et $r_C = 3\,R_T$ du centre de la Terre.

\textbf{2.} Sur un schéma, représenter la Terre (cercle de rayon $R_T$ à l'échelle $\SI{1}{cm}$ pour $R_T$) et les points A, B et C alignés sur une même demi-droite issue du centre.

\textbf{3.} Reprenter les vecteurs champ de gravitation $\vec{\mathcal{G}}_A$, $\vec{\mathcal{G}}_B$ et $\vec{\mathcal{G}}_C$ à l'échelle $\SI{1}{cm}$ pour $\SI{2,5}{N.kg^{-1}}$.

\textbf{4.} Que deviennent ces vecteurs si l'on choisit trois autres points situés aux mêmes distances mais sur une autre demi-droite ? Quelle propriété du champ de gravitation terrestre cela illustre-t-il ?
\end{exercice}

\begin{exercice}[Altitude où le champ est divisé par 2 puis par 4]
On donne : $R_T = \SI{6370}{km}$.

\textbf{1.} Montrer que le champ de gravitation terrestre à l'altitude $h$ vaut $g(h) = \dfrac{g_0}{\left(1 + \dfrac{h}{R_T}\right)^2}$.

\textbf{2.} Déterminer l'altitude $h_1$ à laquelle $g(h_1) = \dfrac{g_0}{2}$. Exprimer d'abord $h_1$ en fonction de $R_T$, puis donner sa valeur numérique en kilomètres.

\textbf{3.} Déterminer de même l'altitude $h_2$ à laquelle $g(h_2) = \dfrac{g_0}{4}$. Exprimer $h_2$ en fonction de $R_T$, puis calculer sa valeur.

\textbf{4.} La Station spatiale internationale évolue à l'altitude $h = \SI{400}{km}$. Calculer $g$ à cette altitude et commenter l'affirmation fréquente : « les astronautes flottent parce qu'il n'y a plus de gravité là-haut ».
\end{exercice}

\begin{exercice}[Gravité à la surface de Mars]
On donne pour Mars : $M_M = \SI{6,4e23}{kg}$ ; $R_M = \SI{3390}{km}$ ; $G = \SI{6,67e-11}{N.m^2.kg^{-2}}$. Sur Terre : $g_0 = \SI{9,81}{N.kg^{-1}}$.

\textbf{1.} Donner l'expression littérale du champ de gravitation $g_M$ à la surface de Mars, en justifiant l'assimilation de Mars à une sphère homogène.

\textbf{2.} Calculer la valeur de $g_M$.

\textbf{3.} Calculer le poids d'un astronaute de masse $m = \SI{80}{kg}$ à la surface de Mars, puis à la surface de la Terre.

\textbf{4.} Un même sac de matériel pèse-t-il plus lourd sur Mars ou sur Terre ? Sa masse change-t-elle ? Expliquer pourquoi les sauts des astronautes seraient plus faciles sur Mars.
\end{exercice}

\begin{exercice}[Approximation du champ au voisinage du sol]
\textbf{1.} En posant $g(h) = g_0\left(1 + \dfrac{h}{R_T}\right)^{-2}$ et en utilisant l'approximation $(1 + \varepsilon)^{-2} \approx 1 - 2\varepsilon$ valable pour $\varepsilon \ll 1$, montrer que pour $h \ll R_T$ :
\[ g(h) \approx g_0\left(1 - \dfrac{2h}{R_T}\right). \]

\textbf{2.} Calculer la variation relative $\dfrac{g_0 - g(h)}{g_0}$ au sommet du mont Cameroun ($h = \SI{4095}{m}$). On donne $R_T = \SI{6370}{km}$. L'approximation $g \approx g_0$ est-elle justifiée à cette altitude ?

\textbf{3.} On considère que l'approximation $g \approx g_0$ cesse d'être acceptable lorsque la variation relative dépasse $1\,\%$. Déterminer l'altitude limite $h_{\text{lim}}$ correspondante.

\textbf{4.} Comparer $h_{\text{lim}}$ à l'altitude de croisière d'un avion de ligne (environ $\SI{10}{km}$) et à celle d'un satellite en orbite basse (environ $\SI{400}{km}$). Conclure sur le domaine de validité du « champ uniforme ».
\end{exercice}

\courssub{Je me prépare au Bac}

\begin{exercice}[Soleil, Lune et marées — type Bac]
On donne :
\begin{itemize}
\item Constante de gravitation : $G = \SI{6,67e-11}{N.m^2.kg^{-2}}$ ;
\item Masse de la Terre : $M_T = \SI{5,97e24}{kg}$ ; rayon terrestre : $R_T = \SI{6370}{km}$ ;
\item Masse du Soleil : $M_S = \SI{1,99e30}{kg}$ ; distance Terre--Soleil : $d_{TS} = \SI{1,50e8}{km}$ ;
\item Masse de la Lune : $M_L = \SI{7,35e22}{kg}$ ; distance Terre--Lune : $d_{TL} = \SI{3,84e5}{km}$.
\end{itemize}

On assimilera le Soleil, la Terre et la Lune à des corps à répartition sphérique de masse.

\textbf{1.} Rappeler l'expression vectorielle de la force de gravitation exercée par un corps A de masse $m_A$ sur un corps B de masse $m_B$, le vecteur unitaire $\vec{u}$ étant dirigé de A vers B. Faire un schéma soigné représentant les forces d'interaction Terre--Lune.

\textbf{2.} Calculer l'intensité $F_{S/T}$ de la force gravitationnelle exercée par le Soleil sur la Terre.

\textbf{3.} Calculer l'intensité $F_{L/T}$ de la force gravitationnelle exercée par la Lune sur la Terre. Comparer $F_{S/T}$ et $F_{L/T}$ en calculant le rapport $\dfrac{F_{S/T}}{F_{L/T}}$.

\textbf{4.} Les marées ne dépendent pas de la force totale exercée sur la Terre, mais de la \emph{différence} du champ de gravitation entre le point de la surface terrestre le plus proche de l'astre attracteur et le centre de la Terre. On admet que l'effet de marée dû à un astre de masse $M$ situé à la distance $d$ est proportionnel à $\dfrac{M}{d^3}$.
\begin{itemize}
\item[a)] Calculer les rapports $\dfrac{M_S}{d_{TS}^3}$ et $\dfrac{M_L}{d_{TL}^3}$ (on donnera les valeurs numériques avec leurs unités).
\item[b)] En déduire pourquoi, bien que la force exercée par le Soleil sur la Terre soit bien plus grande, les marées observées sur les côtes camerounaises (Kribi, Limbé) sont principalement sensibles à la Lune.
\end{itemize}

\textbf{5.} Calculer la valeur du champ de gravitation créé par le Soleil au voisinage de la Terre, puis celle du champ créé par la Lune au voisinage de la Terre. Comparer à $g_0 = \SI{9,81}{N.kg^{-1}}$ et conclure.
\end{exercice}

\begin{exercice}[Satellite de télécommunications au-dessus du Cameroun — type Bac]
Un satellite de masse $m = \SI{500}{kg}$, utilisé pour les télécommunications, évolue sur une orbite circulaire à l'altitude $h = \SI{800}{km}$ au-dessus de la surface terrestre.

On donne : $G = \SI{6,67e-11}{N.m^2.kg^{-2}}$ ; $M_T = \SI{5,97e24}{kg}$ ; $R_T = \SI{6370}{km}$ ; $g_0 = \SI{9,81}{N.kg^{-1}}$.

\textbf{1.} Faire un schéma représentant la Terre, le satellite S et la force gravitationnelle $\vec{F}$ exercée par la Terre sur le satellite. Préciser la direction et le sens de cette force.

\textbf{2.} Donner l'expression littérale du champ de gravitation terrestre $g(h)$ à l'altitude $h$ en fonction de $G$, $M_T$, $R_T$ et $h$, puis en fonction de $g_0$, $R_T$ et $h$.

\textbf{3.} Calculer la valeur de $g(h)$ à l'altitude $h = \SI{800}{km}$. Exprimer le rapport $\dfrac{g(h)}{g_0}$ en pourcentage.

\textbf{4.} Calculer l'intensité $F$ de la force gravitationnelle subie par le satellite sur son orbite.

\textbf{5.} Calculer le poids $P_0$ du satellite au sol (avant son lancement). Comparer $F$ et $P_0$.

\textbf{6.} Un journaliste affirme : « à $\SI{800}{km}$ d'altitude, le satellite est en impesanteur car la gravité terrestre n'y agit presque plus ». Cette affirmation est-elle correcte ? Justifier à l'aide des résultats précédents et proposer une explication correcte de l'état d'« apesanteur » ressenti à bord.
\end{exercice}

\begin{exercice}[Étude du champ de gravitation terrestre — type Bac]
On donne : $G = \SI{6,67e-11}{N.m^2.kg^{-2}}$ ; masse de la Terre $M_T = \SI{5,97e24}{kg}$ ; rayon terrestre $R_T = \SI{6370}{km}$. On assimile la Terre à une sphère homogène.

\textbf{Partie A : champ à la surface et en altitude}

\textbf{1.} Énoncer le théorème qui permet d'affirmer qu'à l'extérieur d'une sphère homogène, le champ de gravitation est le même que celui créé par un point matériel de même masse placé au centre de la sphère.

\textbf{2.} En déduire l'expression du champ de gravitation terrestre $\mathcal{G}(r)$ en un point situé à la distance $r \geq R_T$ du centre de la Terre. Préciser sa direction et son sens.

\textbf{3.} Montrer qu'à l'altitude $h$, on peut écrire $g(h) = g_0\,\dfrac{R_T^2}{(R_T + h)^2}$ où $g_0$ est la valeur du champ à la surface. Calculer $g_0$.

\textbf{4.} Compléter le tableau suivant (valeurs de $g$ en \si{N.kg^{-1}}, arrondies au centième) :

\begin{center}
\begin{tabularx}{\linewidth}{|l|*{5}{X|}}
\hline
\rowcolor{popblueL}
$h$ (km) & 0 & 100 & 1000 & 6370 & 35786 \\
\hline
$g(h)$ (\si{N.kg^{-1}}) & & & & & \\
\hline
\end{tabularx}
\end{center}

\textbf{5.} L'altitude $h = \SI{35786}{km}$ correspond à l'orbite géostationnaire utilisée par les satellites de télécommunications. Que vaut le rapport $\dfrac{g}{g_0}$ à cette altitude ?

\textbf{Partie B : pesée de précision}

\textbf{6.} Un laboratoire de métrologie à Douala (altitude $h_1 \approx 0$) mesure le poids d'une masse étalon $m = \SI{1,000}{kg}$. Le même étalon est ensuite transporté à l'observatoire situé à l'altitude $h_2 = \SI{2000}{m}$. En utilisant l'approximation $g(h) \approx g_0\left(1 - \dfrac{2h}{R_T}\right)$, calculer la variation de poids $\Delta P = P_1 - P_2$ entre les deux lieux. Commenter.

\textbf{7.} Expliquer qualitativement pourquoi la valeur réelle de $g$ à la surface de la Terre dépend aussi légèrement de la latitude (on pourra évoquer la forme de la Terre et sa rotation).
\end{exercice}

\newpage

\coursec{Corrigés des exercices}

\begin{corrige}[Exercice 1 — QCM : la loi de Newton]
\textbf{1. Réponse c).} La loi de l'attraction universelle s'écrit $F = G\,\dfrac{m_A\,m_B}{d^2}$. Vérification par analyse dimensionnelle :
\[
[G]\times\frac{[m_A][m_B]}{[d]^2} = \si{N.m^2.kg^{-2}} \times \frac{\si{kg^2}}{\si{m^2}} = \si{N}. \quad \checkmark
\]
Les propositions a) (somme des masses au lieu du produit), b) (distance au dénominateur à la puissance 1 au lieu de 2) et d) ($G$ au dénominateur) sont dimensionnellement ou physiquement fausses.

\textbf{2. Réponse b).} La force varie en $1/d^2$. Si $d' = 2d$ :
\[
F' = G\,\frac{m_A m_B}{(2d)^2} = \frac{1}{4}\,G\,\frac{m_A m_B}{d^2} = \frac{F}{4}.
\]
La force est divisée par 4.

\textbf{3. Réponse a).} $G = \SI{6,67e-11}{N.m^2.kg^{-2}}$ (constante universelle de gravitation). La proposition b) est la valeur de $g_0$ (champ de pesanteur terrestre, qui n'est pas une constante universelle), c) a des unités inversées, d) est le nombre d'Avogadro $N_A$ (grandeur de chimie).

\textbf{4. Réponse b).} D'après le théorème de Newton, une sphère homogène crée à l'extérieur le même champ qu'un point matériel de masse $M$ placé en son centre : $\mathcal{G} = G\,\dfrac{M}{r^2}$. Le champ est \emph{attractif}, donc dirigé \textbf{vers le centre} de la sphère (ce qui élimine d). Les propositions a) et c) sont dimensionnellement ou physiquement incorrectes.
\end{corrige}

\begin{corrige}[Exercice 2 — Vrai ou faux ?]
\textbf{1. FAUX.} D'après la troisième loi de Newton (principe des actions réciproques) :
\[
\vec{F}_{\text{Terre}\to\text{Lune}} = -\vec{F}_{\text{Lune}\to\text{Terre}}.
\]
Les deux forces ont la \emph{même intensité} $F = G\,\dfrac{M_T M_L}{d_{TL}^2}$, quelles que soient les masses. C'est l'\emph{accélération} subie qui diffère ($a = F/m$) : la Lune, beaucoup moins massive, est bien plus accélérée que la Terre.

\textbf{2. VRAI.} La gravitation est toujours attractive : $\vec{\mathcal{G}}(M) = -G\,\dfrac{m}{AM^2}\,\vec{u}_{AM}$. Le signe $-$ devant le vecteur unitaire dirigé de $A$ vers $M$ signifie que $\vec{\mathcal{G}}$ pointe de $M$ vers $A$, c'est-à-dire \emph{vers la masse source}.

\textbf{3. VRAI.} La Terre étant assimilée à une sphère homogène, le théorème de Newton donne, à sa surface ($r = R_T$) :
\[
g_0 = \frac{G\,M_T}{R_T^2} = \frac{\num{6,67e-11}\times\num{5,97e24}}{(\num{6,37e6})^2} \approx \SI{9,81}{N.kg^{-1}}.
\]

\textbf{4. FAUX.} $g(h) = g_0\,\dfrac{R_T^2}{(R_T+h)^2}$ : quand $h$ augmente, le dénominateur $(R_T+h)^2$ augmente, donc $g$ \emph{diminue} (loi en $1/r^2$).

\textbf{5. FAUX.} Au sommet du mont Cameroun, $h = \SI{4095}{m} = \SI{4,095}{km}$ :
\[
\frac{g_0 - g(h)}{g_0} \approx \frac{2h}{R_T} = \frac{2\times 4,095}{6370} \approx \num{1,29e-3} \approx \SI{0,13}{\%}.
\]
Le champ y vaut $g \approx 9,81\times(1 - 0,00129) \approx \SI{9,80}{N.kg^{-1}}$ : il est \emph{légèrement plus faible} qu'au niveau de la mer. L'écart est faible mais non nul ; « rigoureusement égal » est donc faux.
\end{corrige}

\begin{corrige}[Exercice 3 — Définitions et restitution de cours]
\textbf{1. Loi de l'attraction universelle.} Deux points matériels $A$ et $B$, de masses $m_A$ et $m_B$, distants de $d = AB$, exercent l'un sur l'autre des forces attractives de même intensité. La force exercée par $A$ sur $B$ s'écrit :
\[
\vec{F}_{A\to B} = -\,G\,\frac{m_A\,m_B}{d^2}\,\vec{u}_{AB},
\]
où $\vec{u}_{AB}$ est le vecteur unitaire dirigé de $A$ vers $B$ et $G = \SI{6,67e-11}{N.m^2.kg^{-2}}$. \emph{Conditions d'application :} corps ponctuels ou à répartition de masse à symétrie sphérique (la distance $d$ est alors prise entre les centres).

\textbf{2. Champ de gravitation d'un point matériel.} Le vecteur champ de gravitation créé en $M$ par un point matériel $A$ de masse $m$ est la force qu'exercerait $A$ sur l'unité de masse placée en $M$ :
\[
\vec{\mathcal{G}}(M) = \frac{\vec{F}_{A\to M}}{m_0} = -\,G\,\frac{m}{AM^2}\,\vec{u},
\]
où $\vec{u}$ est le vecteur unitaire dirigé de $A$ vers $M$. Direction : radiale ; sens : vers $A$ ; intensité : $\mathcal{G} = G\,m/r^2$ ; unité : \si{N.kg^{-1}}.

\textbf{3. Théorème de Newton (théorème des coques).} En tout point $M$ \emph{extérieur} à une sphère homogène (pleine ou creuse) de masse $M$ et de centre $O$ ($r = OM \geq R$), le champ de gravitation est le même que celui créé par un point matériel de masse $M$ placé en $O$ :
\[
\vec{\mathcal{G}}(M) = -\,G\,\frac{M}{r^2}\,\vec{u}_{OM}.
\]
(Complément : à l'intérieur d'une sphère creuse homogène, le champ est nul.)

\textbf{4. Démonstration de $g(h)$.} À l'altitude $h$, la distance au centre est $r = R_T + h$. En appliquant le théorème de Newton :
\[
g(h) = \frac{G\,M_T}{(R_T+h)^2}.
\]
On fait apparaître $g_0 = \dfrac{G M_T}{R_T^2}$ en multipliant numérateur et dénominateur par $R_T^2$ :
\[
g(h) = \frac{G\,M_T}{R_T^2}\times\frac{R_T^2}{(R_T+h)^2} = g_0\,\frac{R_T^2}{(R_T+h)^2}. \quad \blacksquare
\]
\end{corrige}

\begin{corrige}[Exercice 4 — Calculs directs]
\textbf{1. Force sur le sac de riz.} La Terre étant assimilée à une sphère homogène, la distance à utiliser est $R_T = \SI{6,37e6}{m}$ :
\[
F = \frac{G\,M_T\,m}{R_T^2} = \frac{\num{6,67e-11}\times\num{5,97e24}\times 50}{(\num{6,37e6})^2}.
\]
Calcul intermédiaire : $G M_T = \num{6,67e-11}\times\num{5,97e24} = \num{3,98e14}{N.m^2.kg^{-1}}$ et $R_T^2 = \num{4,06e13}{m^2}$.
\[
F = \frac{\num{3,98e14}\times 50}{\num{4,06e13}} = \frac{\num{1,99e16}}{\num{4,06e13}} \approx \SI{4,91e2}{N} \approx \SI{491}{N}.
\]

\textbf{2. Valeur de $g_0$.}
\[
g_0 = \frac{G\,M_T}{R_T^2} = \frac{\num{3,98e14}}{\num{4,06e13}} \approx \SI{9,81}{N.kg^{-1}}.
\]
On retrouve exactement la valeur usuelle (3 chiffres significatifs) : le modèle de la sphère homogène est excellent. On vérifie au passage la cohérence : $F = m\,g_0 = 50\times 9,81 \approx \SI{491}{N}$, conforme à la question 1.

\textbf{3. Force entre les deux élèves.}
\[
F = G\,\frac{m_1 m_2}{d^2} = \num{6,67e-11}\times\frac{60\times 60}{(1,0)^2} = \num{6,67e-11}\times 3600 \approx \SI{2,4e-7}{N}.
\]
\emph{Commentaire :} cette force représente environ $\dfrac{\num{2,4e-7}}{60\times 9,81} \approx \num{4e-10}$ du poids d'un élève. Elle est totalement négligeable devant le poids et les forces de contact (frottements) : on ne la ressent pas. L'interaction gravitationnelle ne devient macroscopiquement perceptible que lorsqu'au moins l'une des masses est astronomique (planète, étoile).
\end{corrige}

\begin{corrige}[Exercice 5 — Représentation du champ terrestre à différentes distances]
\textbf{1. Valeurs du champ.} D'après le théorème de Newton, $\mathcal{G}(r) = \dfrac{G M_T}{r^2} = g_0\,\dfrac{R_T^2}{r^2}$.
\begin{align*}
\mathcal{G}_A &= g_0\,\frac{R_T^2}{R_T^2} = g_0 = \SI{9,81}{N.kg^{-1}},\\
\mathcal{G}_B &= g_0\,\frac{R_T^2}{(2R_T)^2} = \frac{g_0}{4} \approx \SI{2,45}{N.kg^{-1}},\\
\mathcal{G}_C &= g_0\,\frac{R_T^2}{(3R_T)^2} = \frac{g_0}{9} \approx \SI{1,09}{N.kg^{-1}}.
\end{align*}

\textbf{2. Schéma.} On trace un cercle de centre $O$ et de rayon $\SI{1}{cm}$ (la Terre), puis une demi-droite horizontale issue de $O$ ; on place $A$ à $\SI{1}{cm}$, $B$ à $\SI{2}{cm}$ et $C$ à $\SI{3}{cm}$ de $O$.

\textbf{3. Vecteurs.} Échelle : $\SI{1}{cm} \leftrightarrow \SI{2,5}{N.kg^{-1}}$.
\begin{itemize}
\item $\vec{\mathcal{G}}_A$ : longueur $\dfrac{9,81}{2,5} \approx \SI{3,9}{cm}$, porté par $(OA)$, dirigé de $A$ vers $O$ ;
\item $\vec{\mathcal{G}}_B$ : longueur $\dfrac{2,45}{2,5} \approx \SI{1,0}{cm}$, dirigé de $B$ vers $O$ ;
\item $\vec{\mathcal{G}}_C$ : longueur $\dfrac{1,09}{2,5} \approx \SI{0,44}{cm}$, dirigé de $C$ vers $O$.
\end{itemize}

\begin{popfigure}
\begin{tikzpicture}[scale=1.1]
\draw[popblue, thick, fill=popblue!10] (0,0) circle (1cm);
\node[popdark] at (0,0) {\small Terre};
\draw[popline] (0,0) -- (3.4,0);
\foreach \r/\l/\g in {1/A/3.9, 2/B/0.98, 3/C/0.44} {
  \fill[popdark] (\r,0) circle (1.5pt) node[above=3pt] {$\l$};
  \draw[->, poporange, very thick] (\r,0) -- (\r,0) ++(-\g*0.4,0);
}
\node[poporange, below] at (0.6,0) {$\vec{\mathcal{G}}_A$};
\node[poporange, below] at (1.8,0) {$\vec{\mathcal{G}}_B$};
\node[poporange, below] at (2.9,0) {$\vec{\mathcal{G}}_C$};
\end{tikzpicture}
\legende{Vecteurs champ de gravitation en A, B, C : radiaux, centripètes, d'intensité décroissante en $1/r^2$.}
\end{popfigure}

\textbf{4. Autre demi-droite.} Les trois nouveaux vecteurs auraient les \emph{mêmes normes} ($9,81$ ; $2,45$ ; $\SI{1,09}{N.kg^{-1}}$) et seraient toujours dirigés vers $O$ le long de la nouvelle demi-droite. Cela illustre la \cle{symétrie sphérique} du champ : son intensité ne dépend que de la distance $r$ au centre, pas de la direction. Les lignes de champ sont des demi-droites radiales convergentes vers $O$.
\end{corrige}

\begin{corrige}[Exercice 6 — Altitude où le champ est divisé par 2 puis par 4]
\textbf{1. Expression de $g(h)$.} À l'altitude $h$, $r = R_T + h$, d'où :
\[
g(h) = \frac{G M_T}{(R_T+h)^2} = \frac{G M_T}{R_T^2}\times\frac{1}{\left(1+\dfrac{h}{R_T}\right)^2} = \frac{g_0}{\left(1+\dfrac{h}{R_T}\right)^2}. \quad \blacksquare
\]

\textbf{2. Altitude $h_1$ telle que $g = g_0/2$.}
\[
\frac{g_0}{\left(1+\dfrac{h_1}{R_T}\right)^2} = \frac{g_0}{2}
\;\Longrightarrow\; \left(1+\frac{h_1}{R_T}\right)^2 = 2
\;\Longrightarrow\; 1+\frac{h_1}{R_T} = \sqrt{2}.
\]
\[
\boxed{h_1 = R_T\left(\sqrt{2}-1\right)} \approx 0,414 \times 6370 \approx \SI{2638}{km}.
\]

\textbf{3. Altitude $h_2$ telle que $g = g_0/4$.}
\[
\left(1+\frac{h_2}{R_T}\right)^2 = 4 \;\Longrightarrow\; 1+\frac{h_2}{R_T} = 2 \;\Longrightarrow\; \boxed{h_2 = R_T = \SI{6370}{km}}.
\]
Vérification : à $r = 2R_T$, $g = g_0/4$ (loi en $1/r^2$). Cohérent.

\textbf{4. ISS à $h = \SI{400}{km}$.}
\[
g = 9,81\times\left(\frac{6370}{6370+400}\right)^2 = 9,81\times\left(\frac{6370}{6770}\right)^2 = 9,81\times(0,9409)^2 \approx \SI{8,69}{N.kg^{-1}}.
\]
Soit encore $\dfrac{g}{g_0} \approx \SI{88,5}{\%}$ de la valeur au sol ! \textbf{L'affirmation est fausse} : la gravité est bien présente à bord de l'ISS. Les astronautes « flottent » parce que la station et son contenu sont en \cle{chute libre} permanente autour de la Terre : tout tombe avec la même accélération $g$, donc aucune force de contact (réaction du plancher) ne s'exerce sur les corps. C'est l'\emph{apesanteur} (absence de force d'appui), et non l'absence de gravité.
\end{corrige}

\begin{corrige}[Exercice 7 — Gravité à la surface de Mars]
\textbf{1. Expression littérale.} Mars est une planète quasi sphérique dont la répartition de masse est à symétrie sphérique en bonne approximation (couches concentriques : croûte, manteau, noyau). Le théorème de Newton s'applique donc à sa surface ($r = R_M$) :
\[
g_M = \frac{G\,M_M}{R_M^2}.
\]

\textbf{2. Application numérique.} $R_M = \SI{3,39e6}{m}$ :
\[
g_M = \frac{\num{6,67e-11}\times\num{6,4e23}}{(\num{3,39e6})^2} = \frac{\num{4,27e13}}{\num{1,149e13}} \approx \SI{3,7}{N.kg^{-1}}.
\]

\textbf{3. Poids de l'astronaute.}
\[
P_{\text{Mars}} = m\,g_M = 80\times 3,7 \approx \SI{2,97e2}{N} \approx \SI{297}{N} ;
\qquad
P_{\text{Terre}} = m\,g_0 = 80\times 9,81 \approx \SI{785}{N}.
\]

\textbf{4. Comparaison.} Le sac pèse \emph{plus lourd sur Terre} ($P_T > P_M$ car $g_0 > g_M$). Sa \textbf{masse ne change pas} : la masse est une grandeur intrinsèque (quantité de matière, inertie), indépendante du lieu ; seul le \emph{poids}, force d'attraction locale, dépend de $g$.
Pour le saut : à vitesse initiale $v_0$ identique, la hauteur atteinte est $h = \dfrac{v_0^2}{2g}$. Comme $g_M \approx g_0/2,6$, on saute environ \textbf{2,6 fois plus haut} sur Mars, et la durée du saut ($t = 2v_0/g$) est aussi 2,6 fois plus longue : d'où la sensation de « flotter ».
\end{corrige}

\begin{corrige}[Exercice 8 — Approximation du champ au voisinage du sol]
\textbf{1. Démonstration.} On pose $\varepsilon = \dfrac{h}{R_T}$ avec $\varepsilon \ll 1$ et on utilise le développement limité $(1+\varepsilon)^{-2} \approx 1 - 2\varepsilon$ :
\[
g(h) = g_0\left(1+\frac{h}{R_T}\right)^{-2} = g_0(1+\varepsilon)^{-2} \approx g_0(1-2\varepsilon) = g_0\left(1-\frac{2h}{R_T}\right). \quad \blacksquare
\]

\textbf{2. Mont Cameroun.} $h = \SI{4095}{m} = \SI{4,095}{km}$ :
\[
\frac{g_0-g(h)}{g_0} \approx \frac{2h}{R_T} = \frac{2\times 4,095}{6370} \approx \num{1,29e-3} \approx \SI{0,13}{\%}.
\]
La variation est de l'ordre du millième : l'approximation $g \approx g_0$ est \textbf{tout à fait justifiée} au sommet du mont Cameroun (et a fortiori partout au sol).

\textbf{3. Altitude limite.} On impose $\dfrac{2h_{\text{lim}}}{R_T} = 1\,\% = 0,01$ :
\[
h_{\text{lim}} = \frac{0,01\times R_T}{2} = \frac{0,01\times 6370}{2} \approx \SI{32}{km}.
\]

\textbf{4. Comparaison et conclusion.}
\begin{itemize}
\item Avion de ligne ($h \approx \SI{10}{km} < h_{\text{lim}}$) : variation relative $\approx \dfrac{2\times 10}{6370} \approx \SI{0,31}{\%}$ : le modèle du champ uniforme reste \textbf{valide}.
\item Satellite en orbite basse ($h \approx \SI{400}{km} \gg h_{\text{lim}}$) : variation $\approx \dfrac{2\times 400}{6370} \approx \SI{12,6}{\%}$ : le modèle du champ uniforme est \textbf{caduc} ; il faut utiliser la formule exacte $g(h) = g_0\,\dfrac{R_T^2}{(R_T+h)^2}$.
\end{itemize}
\textbf{Conclusion :} le « champ uniforme » ($\vec{g}$ constant en norme, direction et sens) est une excellente approximation pour tous les mouvements au voisinage du sol (chutes, projectiles, sport, BTP, aviation), mais ne convient plus dès que l'altitude dépasse quelques dizaines de kilomètres (ballons stratosphériques, fusées, satellites).
\end{corrige}

\begin{corrige}[Exercice 9 — Soleil, Lune et marées — type Bac]
\emph{Barème indicatif (sur 10) : 1. (1,5 pt) ; 2. (1,5 pt) ; 3. (2 pts) ; 4. (3 pts) ; 5. (2 pts).}

\textbf{1. Expression vectorielle et schéma.} La force exercée par A sur B s'écrit :
\[
\vec{F}_{A\to B} = -\,G\,\frac{m_A\,m_B}{d^2}\,\vec{u}_{AB},
\]
avec $\vec{u}_{AB}$ unitaire dirigé de A vers B ; le signe $-$ traduit le caractère attractif (force dirigée de B vers A).

\begin{popfigure}
\begin{tikzpicture}[scale=0.9]
\draw[popblue, thick, fill=popblue!15] (-3.5,0) circle (0.9cm);
\node[popdark] at (-3.5,0) {\small Terre};
\draw[popmuted, thick, fill=popmuted!25] (3.5,0) circle (0.35cm);
\node[popdark] at (3.5,0.7) {\small Lune};
\draw[<->, popdark] (-2.6,-1.3) -- (3.15,-1.3) node[midway, below] {$d_{TL}$};
\draw[->, poporange, very thick] (3.5,0) -- (1.9,0) node[midway, above] {$\vec{F}_{T\to L}$};
\draw[->, popgreen, very thick] (-3.5,0) -- (-1.9,0) node[midway, above] {$\vec{F}_{L\to T}$};
\end{tikzpicture}
\legende{Interaction Terre--Lune : forces opposées, de même intensité (3\up{e} loi de Newton).}
\end{popfigure}

\textbf{2. Force Soleil--Terre.} $d_{TS} = \SI{1,50e8}{km} = \SI{1,50e11}{m}$ :
\[
F_{S/T} = G\,\frac{M_S M_T}{d_{TS}^2} = \frac{\num{6,67e-11}\times\num{1,99e30}\times\num{5,97e24}}{(\num{1,50e11})^2}
= \frac{\num{7,92e44}}{\num{2,25e22}} \approx \SI{3,52e22}{N}.
\]

\textbf{3. Force Lune--Terre.} $d_{TL} = \SI{3,84e5}{km} = \SI{3,84e8}{m}$ :
\[
F_{L/T} = G\,\frac{M_L M_T}{d_{TL}^2} = \frac{\num{6,67e-11}\times\num{7,35e22}\times\num{5,97e24}}{(\num{3,84e8})^2}
= \frac{\num{2,93e37}}{\num{1,47e17}} \approx \SI{1,99e20}{N}.
\]
\[
\frac{F_{S/T}}{F_{L/T}} = \frac{\num{3,52e22}}{\num{1,99e20}} \approx \num{177}.
\]
Le Soleil attire la Terre environ \textbf{177 fois plus fort} que la Lune.

\textbf{4. Effet de marée.} a) Les rapports s'expriment en \si{kg.m^{-3}} :
\[
\frac{M_S}{d_{TS}^3} = \frac{\num{1,99e30}}{(\num{1,50e11})^3} = \frac{\num{1,99e30}}{\num{3,375e33}} \approx \SI{5,9e-4}{kg.m^{-3}},
\]
\[
\frac{M_L}{d_{TL}^3} = \frac{\num{7,35e22}}{(\num{3,84e8})^3} = \frac{\num{7,35e22}}{\num{5,66e25}} \approx \SI{1,3e-3}{kg.m^{-3}}.
\]
b) Rapport des effets de marée :
\[
\frac{M_L/d_{TL}^3}{M_S/d_{TS}^3} \approx \frac{\num{1,3e-3}}{\num{5,9e-4}} \approx 2,2.
\]
L'effet de marée lunaire est donc environ \textbf{2,2 fois plus important} que l'effet solaire. Les marées dépendent de la \emph{différence} d'attraction entre les deux faces de la Terre (terme en $M/d^3$) et non de la force totale (en $M/d^2$) : la proximité de la Lune compense largement sa faible masse. C'est pourquoi les marées observées à Kribi ou Limbé suivent principalement le rythme lunaire (le Soleil ne faisant que moduler l'amplitude : vives-eaux quand Soleil, Terre et Lune sont alignés, mortes-eaux en quadrature).

\textbf{5. Champs créés au voisinage de la Terre.}
\[
\mathcal{G}_S = \frac{G M_S}{d_{TS}^2} = \frac{F_{S/T}}{M_T} = \frac{\num{3,52e22}}{\num{5,97e24}} \approx \SI{5,9e-3}{N.kg^{-1}},
\]
\[
\mathcal{G}_L = \frac{G M_L}{d_{TL}^2} = \frac{F_{L/T}}{M_T} = \frac{\num{1,99e20}}{\num{5,97e24}} \approx \SI{3,3e-5}{N.kg^{-1}}.
\]
Comparaison à $g_0 = \SI{9,81}{N.kg^{-1}}$ : $\mathcal{G}_S \approx 0,06\,\%\,g_0$ et $\mathcal{G}_L \approx 0,0003\,\%\,g_0$. Ces champs sont \textbf{totalement négligeables} devant le champ terrestre pour les mouvements au sol : ils ne modifient pas le poids des objets ; leurs seuls effets mesurables sont les effets différentiels (marées) et le maintien des orbites.

\begin{attention}[Points de vigilance — type Bac]
\begin{itemize}
\item Convertir \emph{systématiquement} les distances en mètres avant tout calcul (les distances sont données en km !).
\item Ne pas confondre force totale ($\propto M/d^2$) et effet de marée ($\propto M/d^3$) : c'est le cœur du paradoxe de l'exercice.
\item Citer le théorème de Newton pour justifier l'assimilation des astres à des points matériels.
\item Donner les puissances de dix avec soin et conclure chaque calcul par une valeur arrondie avec unité.
\end{itemize}
\end{attention}
\end{corrige}

\begin{corrige}[Exercice 10 — Satellite de télécommunications au-dessus du Cameroun — type Bac]
\emph{Barème indicatif (sur 10) : 1. (1,5 pt) ; 2. (1,5 pt) ; 3. (2 pts) ; 4. (1 pt) ; 5. (1,5 pt) ; 6. (2,5 pts).}

\textbf{1. Schéma.} La force $\vec{F}$ exercée par la Terre sur le satellite S est \emph{radiale} (portée par la droite OS) et \emph{centripète} (dirigée de S vers le centre O de la Terre).

\begin{popfigure}
\begin{tikzpicture}[scale=0.85]
\draw[popblue, thick, fill=popblue!10] (0,0) circle (1.6cm);
\node[popdark] at (0,0) {\small Terre};
\draw[popline, dashed] (0,0) circle (2.4cm);
\fill[popdark] (2.4,0) circle (2.5pt) node[right=3pt] {S};
\draw[->, poporange, very thick] (2.4,0) -- (1.2,0) node[midway, above] {$\vec{F}$};
\draw[<->, popdark] (0,-0.35) -- (1.6,-0.35) node[midway, below] {\small $R_T$};
\draw[<->, popdark] (1.6,0.35) -- (2.4,0.35) node[midway, above] {\small $h$};
\end{tikzpicture}
\legende{Force gravitationnelle subie par le satellite : radiale et centripète.}
\end{popfigure}

\textbf{2. Expressions de $g(h)$.} À l'altitude $h$, $r = R_T + h$ :
\[
g(h) = \frac{G\,M_T}{(R_T+h)^2} = g_0\,\frac{R_T^2}{(R_T+h)^2} \quad \text{avec} \quad g_0 = \frac{G M_T}{R_T^2}.
\]

\textbf{3. Valeur à $h = \SI{800}{km}$.} $r = 6370 + 800 = \SI{7170}{km}$ :
\[
g(h) = 9,81\times\left(\frac{6370}{7170}\right)^2 = 9,81\times(0,8884)^2 = 9,81\times 0,7893 \approx \SI{7,74}{N.kg^{-1}}.
\]
\[
\frac{g(h)}{g_0} = 0,789 \approx \SI{79}{\%}.
\]

\textbf{4. Force subie par le satellite.}
\[
F = m\,g(h) = 500\times 7,74 \approx \SI{3,87e3}{N} \approx \SI{3870}{N}.
\]

\textbf{5. Poids au sol et comparaison.}
\[
P_0 = m\,g_0 = 500\times 9,81 \approx \SI{4905}{N}.
\]
\[
\frac{F}{P_0} = \frac{3870}{4905} \approx 0,79 : \quad \text{la force en orbite vaut encore } \approx 79\,\% \text{ du poids au sol.}
\]

\textbf{6. Analyse de l'affirmation du journaliste.} L'affirmation est \textbf{fausse} : à $\SI{800}{km}$ d'altitude, le champ de gravitation vaut encore $\SI{7,74}{N.kg^{-1}}$, soit 79\,\% de sa valeur au sol — loin d'être « presque nul ». C'est d'ailleurs cette force gravitationnelle qui joue le rôle de \emph{force centripète} et maintient le satellite sur son orbite circulaire : sans elle, le satellite partirait en ligne droite (première loi de Newton).
\textbf{Explication correcte :} le satellite et tout son contenu sont en \cle{chute libre} permanente autour de la Terre : ils possèdent tous la même accélération $\vec{g}(h)$. Aucun objet ne s'appuie sur le plancher de la cabine, donc aucune force de contact ne se manifeste : c'est l'\cle{apesanteur} (absence de pesée, pas absence de gravité).

\begin{attention}[Points de vigilance — type Bac]
\begin{itemize}
\item Sur le schéma, la flèche $\vec{F}$ doit partir du satellite et pointer \emph{vers le centre de la Terre} (pas vers la surface, pas tangentielle).
\item Justifier l'usage de $r = R_T + h$ (distance au \emph{centre}, pas à la surface).
\item La distinction impesanteur / apesanteur est un classique du Bac : l'apesanteur résulte de la chute libre, jamais d'une « disparition » de la gravité.
\end{itemize}
\end{attention}
\end{corrige}

\begin{corrige}[Exercice 11 — Étude du champ de gravitation terrestre — type Bac]
\emph{Barème indicatif (sur 12) : Partie A : 1. (1 pt) ; 2. (1,5 pt) ; 3. (2 pts) ; 4. (2,5 pts) ; 5. (1 pt). Partie B : 6. (2,5 pts) ; 7. (1,5 pt).}

\textbf{Partie A}

\textbf{1.} C'est le \cle{théorème de Newton} (ou théorème des coques sphériques) : à l'extérieur d'une sphère à répartition de masse à symétrie sphérique, le champ de gravitation est identique à celui d'un point matériel de même masse placé au centre de la sphère.

\textbf{2.} Pour $r \geq R_T$ :
\[
\vec{\mathcal{G}}(r) = -\,G\,\frac{M_T}{r^2}\,\vec{u}_r, \qquad \mathcal{G}(r) = \frac{G\,M_T}{r^2},
\]
où $\vec{u}_r$ est le vecteur unitaire dirigé du centre O vers le point considéré. \textbf{Direction :} radiale (droite passant par O). \textbf{Sens :} vers le centre O de la Terre (champ attractif, centripète).

\textbf{3.} À l'altitude $h$, $r = R_T + h$ :
\[
g(h) = \frac{G M_T}{(R_T+h)^2} = \underbrace{\frac{G M_T}{R_T^2}}_{g_0}\times\frac{R_T^2}{(R_T+h)^2} = g_0\,\frac{R_T^2}{(R_T+h)^2}. \quad \blacksquare
\]
\[
g_0 = \frac{\num{6,67e-11}\times\num{5,97e24}}{(\num{6,37e6})^2} = \frac{\num{3,98e14}}{\num{4,06e13}} \approx \SI{9,81}{N.kg^{-1}}.
\]

\textbf{4. Tableau complété.} Avec $g(h) = 9,81\times\left(\dfrac{6370}{6370+h}\right)^2$ :
\begin{itemize}
\item $h = \SI{100}{km}$ : $g = 9,81\times(6370/6470)^2 = 9,81\times 0,9692 \approx \SI{9,51}{N.kg^{-1}}$ ;
\item $h = \SI{1000}{km}$ : $g = 9,81\times(6370/7370)^2 = 9,81\times 0,7471 \approx \SI{7,33}{N.kg^{-1}}$ ;
\item $h = \SI{6370}{km}$ ($r = 2R_T$) : $g = 9,81/4 \approx \SI{2,45}{N.kg^{-1}}$ ;
\item $h = \SI{35786}{km}$ ($r = \SI{42156}{km}$) : $g = 9,81\times(6370/42156)^2 = 9,81\times 0,02283 \approx \SI{0,22}{N.kg^{-1}}$.
\end{itemize}

\begin{center}
\begin{tabularx}{\linewidth}{|l|*{5}{X|}}
\hline
\rowcolor{popblueL}
$h$ (km) & 0 & 100 & 1000 & 6370 & 35786 \\
\hline
$g(h)$ (\si{N.kg^{-1}}) & 9,81 & 9,51 & 7,33 & 2,45 & 0,22 \\
\hline
\end{tabularx}
\end{center}

\textbf{5. Orbite géostationnaire.}
\[
\frac{g}{g_0} = \left(\frac{R_T}{R_T+h}\right)^2 = \left(\frac{6370}{42156}\right)^2 \approx (0,1511)^2 \approx 0,0228,
\]
soit seulement $\approx \SI{2,3}{\%}$ de la valeur au sol. Le champ y est faible, mais \emph{non nul} : c'est lui qui maintient le satellite géostationnaire sur son orbite.

\textbf{Partie B}

\textbf{6. Variation de poids.} À Douala ($h_1 \approx 0$) : $P_1 = m g_0 = 1,000\times 9,81 = \SI{9,810}{N}$.
À $h_2 = \SI{2000}{m}$, avec l'approximation :
\[
g_2 \approx g_0\left(1-\frac{2h_2}{R_T}\right) = 9,81\times\left(1-\frac{2\times 2}{6370}\right) = 9,81\times(1-\num{6,28e-4}) \approx \SI{9,8038}{N.kg^{-1}}.
\]
\[
P_2 = m g_2 \approx \SI{9,8038}{N} ; \qquad \Delta P = P_1 - P_2 \approx \SI{6,2e-3}{N} = \SI{6,2}{mN}.
\]
\emph{Commentaire :} la variation est minime ($\approx \SI{0,06}{\%}$), imperceptible dans la vie courante, mais parfaitement détectable par une balance de précision : les laboratoires de métrologie doivent corriger leurs mesures en fonction de l'altitude (et de la latitude).

\textbf{7. Dépendance en latitude.} Deux causes principales :
\begin{itemize}
\item \textbf{Forme de la Terre :} la Terre est aplatie aux pôles ($R_{\text{équateur}} \approx \SI{6378}{km} > R_{\text{pôle}} \approx \SI{6357}{km}$). Comme $g \propto 1/R^2$, $g$ est plus fort aux pôles qu'à l'équateur.
\item \textbf{Rotation terrestre :} dans le référentiel terrestre, tout corps subit en outre l'accélération d'entraînement centrifuge, maximale à l'équateur et nulle aux pôles, qui diminue légèrement le champ de pesanteur effectif $\vec{g}$ par rapport au champ de gravitation $\vec{\mathcal{G}}$.
\end{itemize}
Au total, $g$ varie d'environ $\SI{9,78}{N.kg^{-1}}$ (équateur) à $\SI{9,83}{N.kg^{-1}}$ (pôles), soit environ 0,5\,\%.

\begin{attention}[Points de vigilance — type Bac]
\begin{itemize}
\item Le théorème de Newton doit être \emph{cité explicitement} avant d'assimiler la Terre à un point matériel : c'est une justification attendue.
\item Dans le tableau, chaque valeur doit être accompagnée de son unité et arrondie de façon cohérente (au centième ici).
\item Pour la question 6, bien préciser que la \emph{masse} étalon ne change pas : seul le \emph{poids} varie.
\item Pour la question 7, une réponse qualitative claire (aplatissement + rotation) suffit ; inutile de quantifier l'accélération centrifuge.
\end{itemize}
\end{attention}
\end{corrige}

\newpage

\coursec{Fiche bilan}

\begin{popfigure}
\begin{tikzpicture}[mindmap, grow cyclic, every node/.style={concept, text=white, font=\small\bfseries},
  root concept/.append style={concept color=popink, font=\large\bfseries},
  level 1/.append style={level distance=3.4cm, sibling angle=72},
  level 1 concept/.append style={font=\small\bfseries}]
\node[root concept]{Gravitation}
  child[concept color=popblue]{ node[align=center] {Loi de Newton\\ $F=G\dfrac{m\,m'}{d^{2}}$} }
  child[concept color=poporange]{ node[align=center] {Champ d'un point\\ $\vec{\mathcal{G}}=-G\dfrac{M}{r^{2}}\vec{u}$} }
  child[concept color=popgreen]{ node[align=center] {Sphère homogène\\ $\mathcal{G}=G\dfrac{M}{r^{2}}$ ($r\geq R$)} }
  child[concept color=poppurple]{ node[align=center] {Champ terrestre\\ $g(h)=g_{0}\dfrac{R_{T}^{2}}{(R_{T}+h)^{2}}$} }
  child[concept color=poppink]{ node[align=center] {Champ uniforme\\ $h\ll R_{T}$ : $g\approx g_{0}$} };
\end{tikzpicture}
\legende{Carte mentale de la leçon : les cinq idées forces de l'interaction gravitationnelle.}
\end{popfigure}

\begin{bilan}
\textbf{Définitions clés}
\begin{itemize}
  \item \cle{Force de gravitation} : force attractive, à distance, exercée entre deux corps massifs ; dirigée selon la droite qui joint leurs centres.
  \item \cle{Champ de gravitation} $\vec{\mathcal{G}}(M)$ en un point $M$ : grandeur vectorielle telle qu'une masse $m$ placée en $M$ subit $\vec{F}=m\,\vec{\mathcal{G}}$. Unité : N/kg (ou m/s$^{2}$).
  \item \cle{Lignes de champ} : courbes tangentes en chaque point à $\vec{\mathcal{G}}$, orientées vers la masse source (champ \emph{centripète}).
  \item \cle{Poids} $\vec{P}=m\,\vec{g}$ : au voisinage du sol, il s'identifie pratiquement à la force de gravitation terrestre.
\end{itemize}

\textbf{Formules à connaître par c\oe ur}
\begin{center}
\begin{tabularx}{\linewidth}{|l|X|l|}\hline
\rowcolor{popblueL} \textbf{Grandeur} & \textbf{Expression} & \textbf{Conditions} \\ \hline
Loi de Newton & $F=G\dfrac{m\,m'}{d^{2}}$ ; $G=\SI{6,67e-11}{N.m^{2}.kg^{-2}}$ & corps à symétrie sphérique ou ponctuels \\ \hline
Champ d'un point matériel & $\vec{\mathcal{G}}=-G\dfrac{M}{r^{2}}\,\vec{u}$ ($\vec{u}$ unitaire sortant) & masse $M$ au point $O$ \\ \hline
Sphère homogène (extérieur) & $\mathcal{G}=G\dfrac{M}{r^{2}}$, tout se passe comme si $M$ était au centre & $r\geq R$ \\ \hline
Champ terrestre & $g(h)=g_{0}\dfrac{R_{T}^{2}}{(R_{T}+h)^{2}}$ ; $g_{0}=\dfrac{G\,M_{T}}{R_{T}^{2}}\approx\SI{9,8}{N/kg}$ & altitude $h$ \\ \hline
Ordre $h\ll R_{T}$ & $g(h)\approx g_{0}\left(1-\dfrac{2h}{R_{T}}\right)\approx g_{0}$ & champ uniforme localement \\ \hline
\end{tabularx}
\end{center}

\textbf{Méthodes express}
\begin{itemize}
  \item \emph{Représenter $\vec{F}$ ou $\vec{\mathcal{G}}$} : point d'application au corps attiré, direction = droite des centres, sens = vers la masse attractive, longueur proportionnelle à la norme (échelle imposée).
  \item \emph{Calculer une force} : convertir les distances en mètres, les masses en kilogrammes, puis appliquer $F=Gmm'/d^{2}$ ; vérifier l'homogénéité.
  \item \emph{Champ en altitude} : écrire $g(h)/g_{0}=R_{T}^{2}/(R_{T}+h)^{2}$ ; pour $h$ petit devant $R_{T}$, conclure que $g$ est quasi constant (champ uniforme).
  \item \emph{Comparaison de champs} : utiliser les rapports pour éliminer $G$ et $M$ (méthode des quotients).
\end{itemize}
\end{bilan}

\begin{aretenir}[Checklist avant le Bac]
\begin{itemize}
  \item[$\square$] Je sais énoncer la loi de l'attraction universelle avec la valeur de $G$ et son unité.
  \item[$\square$] Je sais représenter correctement la force gravitationnelle (direction, sens, point d'application).
  \item[$\square$] Je sais que la force est toujours \emph{attractive} et que les deux forces forment un couple action--réaction.
  \item[$\square$] Je sais écrire $\vec{\mathcal{G}}=-G\dfrac{M}{r^{2}}\vec{u}$ et justifier le signe moins.
  \item[$\square$] Je sais qu'une sphère homogène agit, à l'extérieur, comme une masse ponctuelle placée en son centre.
  \item[$\square$] Je sais démontrer $g(h)=g_{0}\dfrac{R_{T}^{2}}{(R_{T}+h)^{2}}$ à partir de $g_{0}=GM_{T}/R_{T}^{2}$.
  \item[$\square$] Je convertis toujours les distances en mètres avant tout calcul numérique.
  \item[$\square$] Je sais distinguer poids ($\vec{P}=m\vec{g}$, local) et force de gravitation (loi universelle).
  \item[$\square$] Je sais justifier que le champ est uniforme au voisinage du sol ($h\ll R_{T}$).
  \item[$\square$] Je vérifie l'homogénéité de mes résultats et je donne un nombre raisonnable de chiffres significatifs.
  \item[$\square$] Je sais tracer des lignes de champ radiales, orientées vers la masse source.
  \item[$\square$] Je sais exploiter un rapport de champs ou de forces pour éliminer les constantes inconnues.
\end{itemize}
\end{aretenir}

\findecours

\end{document}
